% Project: NYCU GIA 碩士論文模板
% Acknowledgement: 此模板參考交大前輩 https://github.com/indigo40123/NYCU-Thesis-Template 提供的 overleaf 版本 https://www.overleaf.com/latex/templates/national-yang-ming-chiao-tung-university-nycu-thesis-template/zgjfrxhcdcmj
% Author: Yun-Chen Lee yclee@arch.nycu.edu.tw
% Date: 2024/07/03
% Description: 修改為建築所適用的格式，使用 APA 參考文獻格式
%-------------------------------------------------------------------------------

% yclee 文件設定：有浮水印/最終可印刷版本(保留左側裝訂留白距離)/中文/不限制章節從奇數頁開始 (無須修改)
% \documentclass[watermark,final,zh,openany,man]{thesis}  % final mode (version for library)
\documentclass[watermark,final,zh,openany,american]{thesis}  % final mode (version for library)


%-------------------------------------------------------------------------------
% Package Loading (無須修改)
%-------------------------------------------------------------------------------

\usepackage{multirow}     % for multi-row table
\usepackage{booktabs}     % table style used in books
\usepackage{tabularx}     % 允許表格欄位依版面寬度自動換行（避免長中文描述撐破表格）
\usepackage{parskip}      % yclee 取消段落縮排
% \usepackage[]{apacite}
% \usepackage[natbibapa]{apacite}
% \usepackage[american]{babel}
% \usepackage{csquotes}
% \usepackage[style=apa,sortcites=true,sorting=nyt,backend=biber]{biblatex}
% \DeclareLanguageMapping{american}{american-apa}
% \addbibresource{bib/bibliography.bib}
% \usepackage{hyperref} % Required for inserting clickable links.
% \usepackage{natbib} % Required for APA-style citations.
% \usepackage[style=apa, backend=biber]{biblatex}


% \usepackage[T1]{fontenc}
% \usepackage[utf8]{inputenc}
% \usepackage{babel}
% \usepackage{csquotes}
% \usepackage[style=apa, sortcites=true, sorting=nyt, backend=biber]{biblatex}
% \addbibresource{bib/bibliography.bib}

\usepackage{natbib}
\usepackage{amsmath,amssymb}
\usetikzlibrary{arrows.meta,positioning,shapes.geometric,fit,backgrounds,calc}

% 避免內文中夾雜之長程式碼／檔案路徑（\texttt{}，不可斷行）造成 overfull hbox
% （行末溢出版面），允許段落於必要時使用額外的字間彈性空間吸收該長度
\setlength{\emergencystretch}{3em}

% yclee 以 URL (網址連結) 的設定：皆顯示為黑色 (無須修改)
\PassOptionsToPackage{hyphens}{url}\usepackage{hyperref}
\hypersetup{colorlinks, citecolor=black, linkcolor=black, urlcolor=black}
\renewcommand\UrlFont{\color{black}\rmfamily}

% 修改圖片說明文字大小
\usepackage{caption}
\captionsetup[figure]{font=small}
\captionsetup[table]{font=small}
% scriptsize
% footnotesize
% small (建議大小)
% normalsize
% large
% Large


%-------------------------------------------------------------------------------
% Configuration
%-------------------------------------------------------------------------------

% 填寫題目, 作者, 指導教授, 學校, 系所, 日期等資訊 (可到此文件修改這些資訊)
% Title
\title{Optimization of Open Space Pedestrian Routes and Building Morphology Shading Design Considering Urban Pedestrian Thermal Comfort}
\titlezh{考量都市行人熱舒適之開放空間行走路徑與建築型態遮陽設計最佳化}

% Author
\author{Wang, Cheng-An}
\titleauthor{Cheng-An Wang}
\authorzh{王誠安}

% Advisor
\advisor{Hou, June-Hao}
\titleadvisor{June-Hao Hou}
\advisorzh{侯君昊}

% First co-advisor (Leave {} empty if you don't have a co-advisor)
\coadvisorA{}
\coadvisorAzh{}

% Second co-advisor (Leave {} empty if you don't have a second co-advisor)
\coadvisorB{}
\coadvisorBzh{}

% Field
\field{Architecture}

% Institute
\institute{Institute of Architecture}
\institutezh{建築研究所}

% College
\college{College of Humanities and Social Sciences}

% University
\university{National Yang Ming Chiao Tung University}
\universityzh{國立陽明交通大學}

% Location
\location{Taiwan, Republic of China}

% Date of final submission
\degreemonth{July}
\degreeyear{2026}
\degreeyearzh{中華民國~ 一一五年七月}

% Date of final submission
\nation{中華民國}
\yearmonth{一一五年七月}

% Watermark %At present(v.1 2021.Sep.29, v.2 2022.Feb.25), not water wark for NYCU yet.
% \watermark{figures/logo.png}
\watermark{figures/nycu_logo.png}





% 修改 thesis.cls 的預設字型 (無須修改)
\input{config/fonts}

% 修改目錄中「章」的編號格式，使其與內文標題（第X章）一致，避免僅顯示裸數字
\renewcommand{\cftchappresnum}{第}
\renewcommand{\cftchapaftersnum}{章}
\setlength{\cftchapnumwidth}{4em}

% yclee 修改圖目錄格式：不區分章節 (無須修改) ---------------------------------------
\makeatletter
\@removefromreset{figure}{chapter}
\renewcommand{\thefigure}{\arabic{figure}}
\makeatother

\renewcommand{\cftfigpresnum}{圖}
\renewcommand{\cftfigaftersnum}{\hspace{1em}} % 在數字和標題之間增加間距
\setlength{\cftfignumwidth}{3.6em} % 設置圖目錄中數字部分的寬度（放寬以容納兩位數圖號，避免 overfull hbox）

% 修改表目錄格式
\makeatletter
\@removefromreset{table}{chapter}
\renewcommand{\thetable}{\arabic{table}}
\makeatother

\renewcommand{\cfttabpresnum}{表}
\renewcommand{\cfttabaftersnum}{\hspace{1em}} % 在數字和標題之間增加間距
\setlength{\cfttabnumwidth}{3em} % 設置圖目錄中數字部分的寬度

% 經典文獻圖之預留插入位置框：待補入正式授權之文獻圖後，以 \includegraphics 取代 \litfigbox 即可。
\newcommand{\litfigbox}[1]{%
  \fbox{\parbox[c][4.2cm][c]{0.82\textwidth}{\centering\small\itshape%
  【經典文獻圖—預留插入位置，待補入正式授權圖檔】\\[4pt] #1}}}


%-------------------------------------------------------------------------------

\begin{document}

% Generate cover, title, authorization, approval and copyright pages
\maketitle

%-------------------------------------------------------------------------------
% Abstract
%-------------------------------------------------------------------------------

% 中文摘要
\begin{abstractzh}

% TODO: 請在此處撰寫中文摘要（約 300-500 字）

本研究針對台灣高溫夏濕都市環境中行人熱壓力問題，建構一套整合物理資訊時空圖神經網路（PI-ST-GNN）、微氣候加權 $A^*$ 路徑演算法與多目標遺傳演算法（NSGA-II）的生成式都市設計最佳化框架。透過異質圖神經網路取代傳統高運算成本的計算流體力學（CFD）模擬，實現通用熱氣候指數（UTCI）之低延遲即時預測（驗證 $R^2 \approx 0.9966$），並以此驅動考量熱暴露成本的行人路徑規劃與 NSGA-II 多目標形態最佳化，同步輸出最佳化建築形態配置與涼適步行路徑，為氣候調適型都市設計提供即時決策支援。

關鍵詞：圖神經網路、都市微氣候、通用熱氣候指數、路徑最佳化、多目標遺傳演算法、生成式設計

\end{abstractzh}


% 英文摘要
\begin{abstract}

% TODO: Please write the English abstract here (approximately 300-500 words)

This study addresses the problem of extreme pedestrian thermal stress in Taiwan's hot-humid urban environments by developing an integrated generative urban design optimization framework combining Physics-Informed Spatio-Temporal Graph Neural Networks (PI-ST-GNN), microclimate-weighted $A^*$ pathfinding, and the Non-dominated Sorting Genetic Algorithm II (NSGA-II). A heterogeneous graph neural network replaces computationally expensive Computational Fluid Dynamics (CFD) simulations to achieve millisecond-latency real-time prediction of the Universal Thermal Climate Index (UTCI), validated at $R^2 \approx 0.9966$. The predicted thermal field drives pedestrian route planning with thermal exposure cost weighting and NSGA-II multi-objective morphological optimization, simultaneously outputting optimal building configurations and thermally comfortable walking paths to support real-time decision-making in climate-adaptive urban design.

Keywords: Graph Neural Network, Urban Microclimate, Universal Thermal Climate Index, Path Optimization, Multi-Objective Genetic Algorithm, Generative Design

\end{abstract}


% yclee 誌謝：加浮水印
% \begin{acknowledgement}%

終於來到這一步，真是太感動了，謝謝家人，老師和同學。


\thispagestyle{empty}
\end{acknowledgement}

%-------------------------------------------------------------------------------
% 目錄
%-------------------------------------------------------------------------------

% Generate 'Table of Contents', 'List of Figures', and 'List of Tables', and
% Set page numbering to 'arabic'
\maketocs

%-------------------------------------------------------------------------------
% Contents
%-------------------------------------------------------------------------------

% Set page numbering to 'arabic' (1, 2, 3, ...)
\mainmatter

% 內文, 請依照章節順序擺放
\chapter{緒論}
\label{chapter:intro}

\section{研究背景與動機}
\label{sec:background}

台灣都市環境地處副熱帶高溫夏濕氣候帶，其戶外熱舒適威脅隨全球氣候變遷加劇而日趨嚴峻。二十一世紀以來，全球氣候變遷與極端天氣事件的發生頻率與強度顯著加劇；根據聯合國跨政府氣候變化評估委員會（IPCC）的評估報告，全球地表溫度較工業革命前已大幅上升，特別是在熱帶與副熱帶氣候帶，極端高溫與持久性熱浪已逐漸轉變為常態性的氣候背景狀態。與此同時，全球都市化進程以空前速度推進，高密度的不透水鋪面、高蓄熱人工材料、密集多層建築群體，以及大量交通與空調系統釋放的人為熱，共同促成了嚴峻的都市熱島效應（Urban Heat Island, UHI），使都市邊界層之風場與輻射場發生質的改變 \citep{oke1987boundary}。台灣地理位於副熱帶與熱帶交界，屬於典型的高溫夏濕氣候，夏季高相對濕度大幅抑制人體表面的蒸散冷卻效率，使得戶外熱壓力常態性地達到不舒適乃至危險的水準 \citep{erell2011urban}。

然而，都市熱島與微氣候研究長期以來多聚焦於氣象學與環境工程視角的量化描述，對於建築師與都市設計者在實際設計現場所面對的具體困境著墨有限。當設計師嘗試以量體配置、遮蔭策略或綠化佈局回應戶外熱舒適課題時，最直接的痛點在於：任何一次幾何修改的微氣候後果，都必須等待數十分鐘至數小時的模擬運算才能得知，使得「設計—評估—修正」的迭代迴圈在方案發展的早期階段幾乎無法即時進行。這並非單純的運算效能問題，而是根本性地限制了微氣候考量能否真正融入設計思考的節奏——當回饋延遲遠超過設計師調整方案的直覺反應時間，微氣候評估便被迫退居為方案定案後的「事後驗證」角色，而非驅動幾何決策的「前置引導」角色。都市意象研究先驅 Lynch 早在其經典著作《城市意象》中即指出，市民對都市空間品質的感知，建立於可辨識、可體驗的實體環境元素之上 \citep{lynch1960image}；本研究延伸此一設計思維脈絡，主張都市戶外熱舒適同樣應被視為設計師在方案發展早期即可感知、可即時操作的空間性能維度，而非事後才能檢核的工程指標。

如何透過都市實體空間形態的量化研究與優化、開放空間的幾何優化以及永續綠化設施的空間佈局，以實質緩解夏季公共空間的熱壓力並縮短「設計—評估」迴圈的等待時間，已成為數位建築學與都市資訊學（Urban Informatics）等跨域研究社群的重要研究課題 \citep{aman2023ai, westermann2019surrogate}。

傳統的環境評估標準大多仰賴靜態、定點式的空間採樣，然而人類在都市公共空間中的空間行為在本質上是具有高度時空複雜性的移動軌跡，此一落差構成本研究第二個切入點。行人在通勤、購物、休憩漫步的過程中，持續經歷著在二維幾何空間與一維時間軸上連續交疊的動態熱暴露；學術界已指出，適當的路徑遮蔭規劃與微氣候設計能大幅改善都市開放空間的利用率與社會活動品質 \citep{coccolo2018thermal}。若戶外空間長期暴露於極端熱壓力下，將導致都市活動大幅向室內收縮，不僅引發空調能源的惡性循環，更侵奪都市居民享受健康公共生活的權利。因此，如何利用運算設計（Computational Design）工具，將行人動線上的遮蔭幾何優化與都市開放空間的形態配置進行深度聯結，是實踐氣候調適型都市設計的核心出發點。


\section{研究問題}
\label{sec:problem}

傳統高擬真度數值模擬引擎（High-fidelity Simulation Engines）難以支援即時設計迭代，這是現行數位建築設計工作流評估戶外微氣候熱舒適度時的核心運算瓶頸。計算流體力學（CFD）軟體（如 OpenFOAM、ANSYS Fluent）、全尺度微氣候熱力學模擬軟體（如 ENVI-met），以及基於 Rhino/Grasshopper 環境的 Ladybug Tools \citep{sadeghipour2013ladybug}，是學術界與實務界最常使用的評估工具，然而這些工具在應用於生成式探索或即時設計迭代時，存在著難以跨越的運算瓶頸：

\begin{enumerate}
    \item \textbf{非即時性與高時間成本：}為求解穩態納維-斯托克斯方程（Steady-state Navier-Stokes Equations）與複雜三維輻射傳熱方程，傳統 CFD 與環境模擬軟體往往需耗費數十分鐘乃至數日才能完成單一方案的模擬計算。
    \item \textbf{極度依賴高性能硬體：}網格細化之解析度直接影響模擬準確性，高精度網格計算對運算資源有極高需求。
    \item \textbf{難以整合於早期設計階段：}高延遲、難以貫通的模擬反饋機制，導致微氣候分析在實務上淪為方案完成後的事後驗證手段，難以在幾何概念設計階段提供即時的決策引導。
\end{enumerate}

上述三項限制——非即時性與高時間成本、對高性能硬體的高度依賴、以及難以整合於早期設計階段——彼此並非各自獨立的技術瑕疵，而是同一個結構性成因（模擬引擎以求解偏微分方程為核心運算邏輯，其運算複雜度與幾何精細度呈非線性增長）所導致的三種外顯症狀。這三者共同構成本研究\textbf{第一個研究缺口}：因為傳統高擬真度模擬引擎的運算延遲從根本上排除了其作為早期設計迭代工具的可能性，所以微氣候評估始終無法真正進入方案發展的前段，僅能扮演事後驗證的角色。

即使上述運算延遲問題透過某種代理模型獲得緩解，現有大量針對都市形態微氣候最佳化的學術研究，多仍聚焦於「全空間平均最佳化」（Global Spatial Average Optimization）之框架，卻缺乏結合實際行人移動軌跡的動態評估，構成本研究\textbf{第二個研究缺口}。此類方法通常將場地離散化為均質網格，並以整個研究範圍內平均 UTCI 或平均輻射溫度之極小化作為最佳化目標，存在兩項研究盲點：其一，忽視了人類步行距離之行為閾值——行人在夏季戶外極端熱壓力下的活動距離通常限制在 400 至 800 公尺（約 5 至 10 分鐘路程）內，其感受取決於路徑上各路段熱壓力之暴露歷程，而非整個場域之平均網格數值；其二，全空間平均最佳化易導致「空間品質的稀釋」——演算法為追求全局指標最小，可能將建築物與植栽稀疏佈置於整個場地，而非集中強化行人實際通行之帶狀廊道，導致優化後的方案在統計數據上表現優異，實際面對高密度人流的動線仍暴露於極端高溫下。因為現有最佳化框架未將「三維實體幾何配置」與「人行動線熱舒適品質」相互串聯比較，所以目前的研究缺乏一種能同時回應兩者的閉環運算架構。


\section{研究目的}
\label{sec:objectives}

承上節所述之兩項研究缺口，本研究提出對應之解決策略：因為傳統高擬真度模擬引擎之運算延遲難以支援即時設計迭代（第一個研究缺口），所以本研究建構低延遲之深度學習代理模型，取代傳統模擬引擎作為早期設計階段之評估工具；因為現有全空間平均最佳化方法缺乏行人動線之動態熱暴露評估（第二個研究缺口），所以本研究另行定義以人行動線為對象之熱舒適評估指標，並將兩者整合進多目標基因演算法之決策框架，使「形態幾何」與「人行動線熱舒適品質」得以在同一套運算流程中被權衡比較。具體而言，本研究之三項研究目的分別回應如下：

\noindent\textbf{研究目的一：建立低延遲全域 UTCI 預測代理模型（回應第一個研究缺口）}

因為傳統高擬真度模擬引擎之運算瓶頸源自其以求解偏微分方程為核心之運算邏輯，所以本研究不採用傳統像素型或體素型卷積神經網路（CNN）架構，改為建構異質圖神經網路（Heterogeneous Graph Neural Network）作為微氣候預測代理模型 \citep{kipf2017semi}：將都市空間中的空氣節點、建築物等實體定義為異質結構中的相異節點類型，並以邊將空間物理風場的流體聯通性與輻射場的天空視角因子嵌入其中，即時輸出網格點的 UTCI 分布。為進一步區分不同物理因果關係之訊號貢獻，本研究採用關係圖卷積網路（Relational GCN）取代同質圖卷積 \citep{schlichtkrull2018modeling}，並以真實都市建物輪廓與樓高資料取代單純程序化生成之幾何，使模型驗證得以錨定於實際都市紋理，詳見第~\ref{sec:building_extraction}節。

\noindent\textbf{研究目的二：建立熱舒適人行動線設計評估指標（回應第二個研究缺口）}

因為全空間平均最佳化無法反映行人沿路徑之動態熱暴露歷程，所以本研究以 GNN 代理模型輸出之高擬真即時空間 UTCI 數據場為基礎，建立\textbf{熱舒適人行動線設計}（Thermal Comfort-Based Walkway Design）評估指標，將行人動線網絡之熱暴露評估直接整合為多目標最佳化之適應度函數，取代僅評估單一起訖點之傳統尋路演算法（如 $A^*$）：都市開放空間設計決策之對象通常是整體人行動線系統在一整日各時段之熱舒適表現，而非單一使用者的單次行程，此一評估尺度之調整，詳見第~\ref{subsec:from_astar_to_walkway}節。

\noindent\textbf{研究目的三：整合多目標基因演算法之生成設計工作流（回應目的一、二之技術整合需求）}

因為研究目的一、二分別解決了「即時預測」與「動線評估指標」兩項獨立問題，但兩者仍須被統一至同一套設計決策流程才能實際支援方案生成，所以本研究將前兩項深度學習代理模型與熱舒適人行動線評估指標，深度整合至第二代非支配排序遺傳演算法（NSGA-II）\citep{deb2002fast} 之運算框架中，建構完整的多目標生成式決策支援系統：演算法優化多層建築物之樓層平面幾何形態組合、建築量體與高度，以及行道樹之各項參數分佈。透過多目標最佳化之迭代，系統自主找出並輸出非支配的 Pareto 最佳設計解集，供設計師在「全域熱舒適最大化」「人行動線熱舒適最大化」與「永續綠地率」之間進行理性的空間決策。

圖~\ref{fig:overview_blueprint} 呈現上述三項研究目的與後續各章節之對應關係，作為本論文之整體研究藍圖。

\begin{figure}[htbp]
    \centering
    \includegraphics[width=0.92\textwidth]{../urban-thermal-gnn/04_training/figures/fig_arch_clean.png}
    \caption{研究藍圖}
    \label{fig:overview_blueprint}
\end{figure}

上圖由基本 GIS 資訊建構（第三章）與都市形態模擬取得訓練資料，經異質圖轉譯後由 PI-ST-GNN 進行低延遲推理（第三、四章），最終驅動 NSGA-II 多目標最佳化與熱舒適人行動線設計評估，並以 Grasshopper 介面回饋予設計者（第三、五章）。


\section{研究方法概述}
\label{sec:method_overview}

本研究採用之研究方法可歸納為四個階段，完整技術細節於第三章「系統分析與設計」中依序展開：（一）以政府開放 GIS 圖資與 IoT 感測器資料建立研究場域之地理資訊基礎；（二）以 Ladybug Tools／EnergyPlus 執行都市微氣候逐時模擬，取得監督訓練所需之熱壓力標籤；（三）建構異質圖神經網路並以物理引導損失函數約束訓練，使模型輸出兼具數值精度與熱力學一致性；（四）將訓練完成之代理模型嵌入 NSGA-II 多目標最佳化與熱舒適人行動線設計評估，並透過 FastAPI／WebSocket \citep{rfc6455} 與 Rhino Grasshopper 介面整合，形成可即時互動之閉環設計決策系統。

上述四個階段與前節三項研究目的之對應關係如下：階段（一）至（三）——地理資訊建構、微氣候模擬取樣、異質圖神經網路訓練——共同支撐\textbf{研究目的一}之低延遲全域 UTCI 預測代理模型；階段（四）中之熱舒適人行動線評估機制對應\textbf{研究目的二}，NSGA-II 多目標最佳化與 Grasshopper 介面整合則對應\textbf{研究目的三}。此一方法論安排之邏輯，在於使每一階段之技術選擇均直接回應第~\ref{sec:problem}節所述之研究缺口，而非孤立的技術堆疊。


\section{研究範圍與限制}
\label{sec:scope}

本研究在空間尺度的假定上，嚴格聚焦於都市「行區尺度」（Neighborhood Scale / District Scale），具體運算幾何規模範圍設定在邊長約 \textbf{500 公尺至 800 公尺}的正方形都市行區範疇。此一尺度的選定具有充足的科學與實務依據：

\begin{enumerate}
    \item \textbf{環境學依據：}根據都市微氣候學奠基者 Oke 的理論，行區尺度是都市微氣候中「微氣候混合層」與「街道峽谷（Urban Canyons）」次系統的核心交點。在此尺度上，都市表面幾何形態（如建築高度豐富質性、天空視角因子 SVF）與局部風場的空氣動力學粗糙長度（Roughness Length）產生最強的交互作用，是直接決定 pedestrian-scale（人體尺度）微氣候特徵的關鍵尺度 \citep{oke1987boundary}。
    \item \textbf{運算限制：}儘管本研究採用了高效的 GNN 代理模型，但下游 NSGA-II 多目標最佳化之族群個體迭代，仍受到節點數量之指數增長運算負載限制。若下至單棟建築實體尺度，則無法覆蓋都市行區風廓與群體遮蔭的巨觀行為；若上擴至都市總體規劃尺度，則將超出當前計算幾何在 Rhino/Grasshopper 前端資料傳輸 Payload 的可能極限。因此，500m--800m 的行區範疇是兼顧環境行為學實用性與計算幾何可行性的最佳尺度。
\end{enumerate}

在時間軸與邊界條件的設定條件上，本研究聚焦以夏季夏至熱浪日（6 月 21 日前後）之日間時段（08:00 至 18:00）。此一時間限制主要依據台灣微氣候學研究針對熱帶與副熱帶高密度都市空間之熱舒適研究結論：台灣高溫夏濕環境中的最高熱壓力與都市熱島強度峰值，高度集中於夏至日前後正午至下午之間，此時段的太陽高度角最大，直接輻射與地表二次長波輻射釋放已達飽和，戶外 UTCI 常態性突破 $38\,^\circ\text{C}$，達到強烈乃至極端的熱壓力風險，且為都市行人暴露於紫外線與高溫危害、引發急性公共健康風險的關鍵時段 \citep{su2023outdoor}。以該時段的微氣候邊界條件作為 GNN 代理模型的靜態氣象特徵，對於評估極端環境適應性、建立高防禦性的生成式設計規範，具有統計代表性與極端工程設計（Worst-case Scenario Engineering）的科學效度。

\chapter{文獻回顧}
\label{chapter:relatedwork}

\section{都市空間品質與微氣候之理論脈絡}
\label{sec:urban_theory_context}

都市設計理論對「空間品質應如何被感知與評估」的討論，早於當代微氣候量化研究數十年。Lynch 在其經典著作《城市意象》（The Image of the City）中提出，市民對都市環境的認知建立於路徑（Paths）、邊界（Edges）、區域（Districts）、節點（Nodes）與地標（Landmarks）五類可辨識元素之上，主張都市空間品質是使用者實際體驗、逐步建構而成的認知過程，而非僅由抽象的統計指標所定義 \citep{lynch1960image}。此一觀點對本研究具有雙重啟發：其一，它說明了為何本研究選擇以人行動線（對應 Lynch 之「路徑」概念）而非全場域平均值作為熱舒適評估的空間單元——行人對熱環境品質之感知本即依附於其實際移動路徑，而非場地整體的統計平均；其二，它提示都市微氣候研究若僅停留於工程指標的極小化，將難以回應設計者與使用者對空間品質之真實感知，此一缺口正是本研究第~\ref{subsec:from_astar_to_walkway}節提出熱舒適人行動線設計評估指標之理論動機。

本研究之案例場域為台灣新竹市，其地理位置及氣候特性所帶來之熱舒適挑戰，構成本研究之地域脈絡：新竹市地處台灣西北部沿海平原，屬副熱帶季風氣候，夏季受西南氣流影響，市區高密度商業與居住混合街廓之戶外熱壓力與台北、台中等台灣主要都市具有相近的物理成因，惟其風場受季風與地形交互影響，局部微氣候變異更為顯著。近年台灣本土之戶外熱舒適實測研究已逐步累積不同都市密度街廓之 UTCI 現地量測資料，例如針對台北市不同街道峽谷高寬比之戶外熱舒適比較研究，證實高密度街廓雖天空視角因子較低，但因風速增強與遮蔭效果，其正午熱舒適表現可能優於低密度開闊街廓 \citep{su2023outdoor}，此一發現與本研究第~\ref{sec:physics_loss}節風場阻滯尾流約束之設計動機互為呼應——都市幾何對熱舒適之影響並非單調的「密度越低越涼爽」，而是輻射遮蔽與風場阻滯兩種機制相互競爭的結果，此一非單調性正是純幾何規則難以捕捉、而需要資料驅動代理模型學習之空間關係。


\section{都市微氣候與行人熱舒適度 (UTCI)}
\label{sec:utci_review}

\subsection{都市形態對微氣候之影響因素與熱力機制}
\label{subsec:urban_form_microclimate}

都市微氣候（Urban Microclimate）是指都市地表物理特徵、幾何形態和人類活動干擾後，在近地面大氣邊界層內所形成的局部氣候型態 \citep{oke1987boundary}。在數位建築學與都市形態學（Urban Morphology）的交叉學域中，實體空間的幾何三維配置已被認識為主導局部微氣候變異的核心驅動力 \citep{aman2023ai}。都市建築群的增長與形態，直接改變了地表與大氣之間的輻射、熱量與水分交換機制，進而驅動了微氣候環境的空間演變 \citep{xin2025urbangraph}。

在現有幾何形態分析數之中，天空視角因子（Sky View Factor, SVF）是評估都市開放空間熱輻射行為最具代表性的指標。SVF 定義為某一觀測點上方暴露之立體角佔整個球面角的比例；當都市開發密度越高、建築體越高或越深際時，空間的 SVF 數值便愈低 \citep{coccolo2018thermal}。低 SVF 空間在日間雖可部分遮擋太陽直達短波輻射，但在夜間卻引發嚴重的「輻射長留效應」（Radiation Trapping Effect），即建築多棟與地面相互釋放的長波輻射（Longwave Radiation）在特窄的行道峽谷（Urban Canyons）內發生多重反射，難以向高空散逸，這是導致都市熱島效應（UHI）夜間持繼高溫的主要物理機制 \citep{erell2011urban}。此外，街道峽谷的高寬比（Aspect Ratio, $H/W$）亦與局部風場的空氣動力學特徵密切相關。當 $H/W$ 比例失衡時，近地面風速（$v_{\text{a}}$）極易在迎風面產生強烈的分流與下沖流，或在背風面形成流體分滯的渦流帶（Wake Recirculation Zone），阻礙都市大氣的熱量對流擴散 \citep{xin2025urbangraph}。上述都市型態、夜間熱島之熱遲滯與街谷空氣動力學三者之交互機制，可整合為圖~\ref{fig:lit_urbanform} 之經典理論示意。

\begin{figure}[htbp]
    \centering
    \includegraphics[width=0.92\textwidth]{../urban-thermal-gnn/04_training/figures/fig_lit_urbanform.png}
    \caption[都市街谷幾何、熱島熱遲滯與空氣動力學機制示意]{都市街谷幾何、熱島熱遲滯與空氣動力學機制之經典理論示意：(a) 天空視角因子（SVF）之短波遮擋與夜間長波輻射多重反射所致之熱遲滯效應；(b) 街谷高寬比 $H/W$ 所主導之迎風下沖流與背風渦流帶風場結構。原創繪製，理論基礎參考 \citet{oke1987boundary} 與 \citet{erell2011urban}。}
    \label{fig:lit_urbanform}
\end{figure}

除幾何特徵外，都市表層材料的熱物理屬性亦扮演著關鍵角色。表面反射率（Albedo）決定了反射太陽短波輻射的反射係數，而發射率（Emissivity）則決定了其在吸熱後向外輻射釋放熱能的能力 \citep{liu2022impacts}。為了抵消人工鋪面帶來的熱效應，引入永續綠底設施（Green Infrastructure）已被廣泛認定為調適都市微氣候最有效的手段。樹木等植被覆蓋對都市微氣候的調節作用主要透過兩大機制實現：第一為直接的「幾何遮蔭（Solar Shading）」，透過樹葉層衰減太陽短波輻射的穿透率（Transmittance），降低地面與行人的輻射獲熱量 \citep{coccolo2018thermal}；第二為植物生理層面的「蒸散發作用（Evapotranspiration）」，透過將液態水轉化為氣態水蒸氣，將環境中的顯熱（Sensible Heat）轉化為潛熱（Latent Heat），進而實質降低周圍空氣溫度（$T_{\text{a}}$）並調節相對濕度（$\text{RH}$） \citep{coccolo2018thermal}。實證研究指出，在夏季極端熱壓力天氣條件下，具備高密度公共遮蔭的開放空間與毫無遮蔭的曝曬空間相比，其通用熱氣候指數（Universal Thermal Climate Index, UTCI）之差距可達 $8~^\circ\text{C}$ 至 $10~^\circ\text{C}$ 以上 \citep{coccolo2018thermal,lopezcabeza2025open}。

然而，上述 SVF、$H/W$ 高寬比、Albedo 等幾何形態指標，其共同的方法論預設是將微氣候表現「歸約」為單點或單一街道剖面之靜態純量描述——SVF 是觀測點自身的立體角積分，$H/W$ 是單一街道峽谷的截面比例，兩者皆未顯式表徵「該點與周遭其他實體之間的因果傳遞關係」。這造成兩項本研究認為關鍵的研究缺口：其一，SVF 與 $H/W$ 皆為時間不變（time-invariant）之幾何常數，無法捕捉遮蔭邊界隨太陽方位角逐時移動所產生的動態陰影效應，亦即同一 SVF 數值的兩個地點，可能在正午與傍晚承受截然不同的輻射負荷；其二，這些指標彼此獨立計算，未能表達第~\ref{sec:pathfinding_review}節所述「輻射遮蔽與風場阻滯相互競爭」的多機制耦合現象——高 $H/W$ 峽谷雖降低 SVF、增加遮蔭，卻可能同時因通風不良而升高對流阻力，兩股效應方向相反，單一靜態指標無法同時表徵。本研究以圖結構中的 \texttt{shadow}（遮蔽邊，隨太陽方位角動態建構）、\texttt{convective}（對流邊，依風向建構）與 \texttt{veg\_et}（植被蒸散邊）三種異質邊類型分別編碼這些機制，取代單一靜態幾何純量，其設計動機即在於填補既有形態指標無法表達的動態性與機制耦合性，此設計細節見第~\ref{sec:graph_construction}節。此一顯著的溫差，亦奠定了透過空間設計手段精準釋放遮蔭物質以改善行人感受度的理論基礎。


\subsection{戶外熱舒適指標之演進與選用比對}
\label{subsec:thermal_indices}

為了量化微氣候環境對人體的實質生理影響，環境科學界與熱力生理學界在過去數十年間陸續開發了多種熱舒適指標（Thermal Comfort Indices）。這些指標的演進趨勢，從「統計驗證」到「強化對底層熱平衡機制的深化」。

早期廣泛應用的預測平均票數（Predicted Mean Vote, PMV）是由 Fanger 於 1970 年代基於室內機械溫控環境建立的人體熱平衡方程 \citep{fanger1970thermal}。PMV 假設環境中的風速極低、輻射場均一，且人體處於新陳代謝的穩定平衡狀態。然而，當 PMV 被強行移植於複雜的戶外環境時，其理論假設便面臨嚴重的失效風險。戶外環境存在著劇烈波動的局部風速以及極度非均質的太陽直接、散射與地表長波輻射場，這使得 PMV 在評估戶外行人舒適度時常出現嚴重的預測失誤 \citep{blazejczyk2012comparison}。

為了解決室內指標的局限性，生理等效溫度（Physiological Equivalent Temperature, PET）應運而生 \citep{hoppe1999physiological}。PET 基於慕尼黑能量平衡模型（MEMI），將複雜的戶外氣象條件轉譯為一個等效的室內標準環境溫度，使評估結果具有直覺的溫度感知度，因而被廣泛應用於都市比較與建築戶外環境的實證研究中。然而，PET 在生理學層面仍存在著參數簡化過度。對於高濕度氣候（如台灣的高溫夏濕環境）中人體表面蒸散汗液的自適應調節機制模擬不足等缺陷 \citep{blazejczyk2012comparison}。

基於上述指標的局限，本研究明確選用通用熱氣候指數（UTCI）作為全域微氣候性能評估的核心指標。UTCI 是由國際生物氣候學學會（ISB）與歐洲科技合作委員會（COST Action 730）聯合研發的新興熱舒適體系 \citep{jendritzky2012utci}。其底層核心建構於改進的「Fiala 多節點人體熱傳遞與高溫調節模型」（Fiala et al., 2011） \citep{fiala2011utci}，該模型將人體細分為 12 個軀幹部位、共計 187 個組織節點，能夠高度動態地模擬人體在中樞神經系統調控下，面對多元氣象突變時發生的血管擴張/收縮、排汗（Sweating）以及顫抖等主動生理調適的過程 \citep{fiala2011utci,romaszko2022universal}。UTCI 透過全面整合空氣溫度（$T_{\text{a}}$）、平均輻射溫度（$T_{\text{mrt}}$）、風速（$v_{\text{a}}$）和相對濕度（$\text{RH}$）四大微氣候核心變數，輸出一個具有高度生理學解釋力的離散熱壓力分級（Discrete Thermal-Stress Classifications）\citep{lopezcabeza2025open}。近年來的都市數位孿生（Urban Digital Twins, UDT）與環境健康研究已充分證實，UTCI 在不同氣候帶與複雜都市幾何背景下，具有極高的通用性（Universal Capability）與跨地域的預測可靠性，是執行氣候調適型生成設計（Climate-Responsive Generative Design）最具科學信度的環境績效度量框架 \citep{lopezcabeza2025open,romaszko2022universal}。UTCI 所整合之四項微氣候變數與人體熱生理反應之關係，以及其於台灣高溫濕氣候之在地適用性，可分別以圖~\ref{fig:lit_utci_factors} 綜整。

\begin{figure}[htbp]
    \centering
    \includegraphics[width=0.92\textwidth]{../urban-thermal-gnn/04_training/figures/fig_lit_utci_factors.png}
    \caption[UTCI 生理學合成基礎與台灣在地熱壓力分級對照]{UTCI 之生理學合成基礎與台灣在地熱壓力分級對照：(a) 戶外熱舒適之四項微氣候因素——空氣溫度 $T_a$、平均輻射溫度 $T_{\text{mrt}}$、風速 $v_a$、相對濕度 RH——經 Fiala 多節點人體熱生理模型合成 UTCI 之關係示意 \citep{fiala2011utci}；(b) UTCI 六等熱壓力分級，並標註台灣夏至正午街谷之實測熱壓力常態範圍 \citep{su2023outdoor}。原創繪製。}
    \label{fig:lit_utci_factors}
\end{figure}


\section{建築形態之生成與多目標最佳化 (MOO)}
\label{sec:moo_review}

\subsection{基因演算法於參數化都市設計之應用}
\label{subsec:ga_design}

在都市與建築設計的早期階段（Early Design Stages），空間形態決策在本質上是一個多學科交叉點的複雜決策過程。設計師不僅需要滿足空間機能、美學與法規限制（如容積率、建蔽率與縮退距離），更必須積極最佳化節能減碳能源與戶外環境和最佳化等永續性能指標 \citep{aman2023ai,wu2023data}。由於這些設計目標之間往往存在著強烈的「非線性衝突」——例如，為了增加地面行人路徑的陰影遮蔽而增高建築物高度，可能同時導致鄰近建築的室內採光與空間日光自主性（Spatial Daylight Autonomy, sDA）大幅下降——因此，傳統仰賴人工試誤（Trial-and-Error）的設計範式已無法在龐大的設計解空間（Design Space）中找到最優的平衡點 \citep{daglier2025developing}。

為了破解此一困境，多目標最佳化（Multi-Objective Optimization, MOO）技術逐步被引入數位建築學領域 \citep{westermann2019surrogate}。其中，基於群體進化論適應機制的遺傳演算法（Genetic Algorithms, GA），因其具備強大的全局搜索能力與對複雜連續、非線性決策空間的適應性，成為建築性能最佳化（Building Performance Optimization, BPO）的核心演算法 \citep{wang2005object}。在現有學術文獻中，非支配排序遺傳演算法 II（Non-dominated Sorting Genetic Algorithm II, NSGA-II）是最具代表性與應用最廣泛的當代工作流 \citep{deb2002fast}。NSGA-II 透過快速的非支配排序機制，對種群中的個體依照性能優勢進行分層，並引入「擁擠距離（Crowding Distance）」指標來確保最優解在帕累托前沿（Pareto Front）上的分布均勻性與多樣性，有效避免了演算法過早落入局部最優解（Local Optima）的缺陷 \citep{deb2002fast,daglier2025developing}。透過參數化建模平台將建築物 envelope（外殼幾何）、容積率配置以及永續材質設定為搜索基因，NSGA-II 能夠自動進行數十代乃至數百代的個體代迭，最終輸出一個「帕累托最優解集（Pareto-optimal Solutions）」，讓建築師得以在科學資料的佐識下，彈性平衡各相互衝突的設計性能 \citep{cheng2023urbangenogan,daglier2025developing}。

值得注意的是，NSGA-II 演算法本身作為一種與問題無關（problem-agnostic）的通用搜索策略，其收斂品質與計算成本，完全取決於「適應度函數求值（Fitness Evaluation）」的效率——這是既有 GA/NSGA-II 建築應用文獻中鮮少被明確量化討論的隱性前提。當適應度函數涉及環境物理模擬（如日照分析、CFD 流場求解）時，演算法每代的種群評估時間即被模擬引擎的運算延遲所主導，使得 NSGA-II 理論上的全域搜索優勢，實務上受制於模擬求值次數的預算上限，被迫縮小種群規模或世代數以換取可行的運算時程，此一「演算法-模擬引擎」耦合瓶頸將於下節具體展開。此一限制正是本研究以 GNN 代理模型取代直接模擬引擎、將單次適應度求值時間由分鐘級大幅壓縮至低延遲尺度的核心動機，使 NSGA-II 得以在相同運算預算下執行遠多於既有文獻之世代數與種群規模，詳見第~\ref{sec:nsga2}節。


\subsection{Grasshopper 參數化生成與代理模型最佳化之工具權衡}
\label{subsec:gh_surrogate}

隨著運算建築學（Computational Architecture）的普及，McNeel 公司開發的 Rhinoceros/Grasshopper3D 視覺化參數環境已經成為當代參數化設計與性能模擬的行業標準平台。Grasshopper 的核心優勢在於其極為豐富且高度模組化的第三方元件生態系，使建築師能夠在統一的幾何立書框架內無縫串聯幾何生成與環境運算。

在最佳化工具層面，Grasshopper 內建的 Galapagos 模組是應用最廣的單目標遺傳演算法工具，適用於結構性地形成單一性能目標的極值求解。然而，面對複雜的都市微氣候設計，單一目標往往流於片面。為了實現真正的多目標演化，Wallacei 模組應運而生，該工具深度包裹了進階的 NSGA-II 模擬，不但能夠執行進階的 NSGA-II 模擬，更提供了極為強大的進化資料可視化面板（如平行坐標軸、帕累托前沿三維散點圖），幫助設計師更可視設計變數與性能指標之間的隱性視覺關係。進一步地，針對機器學習工作流，Wortmann 推出的 Opossum 套件引入了徑向基函數最佳化（Radial Basis Function Optimization, RBFOpt）等進階代理模型演算法，適用於處理計算成本高昂的黑盒目標函數最佳化，顯著縮短了演算法找尋帕累托最優解的時間 \citep{wortmann2017opossum,daglier2025developing}。

然而，儘管上述參數化工具的功能已相當成熟，但在實務上面臨的致命限制依然是「高效能數值模型時間成本」。在現行的最佳化架構中，每當遺傳演算法生成一組新的幾何個體，便必須調用後端的環境物理模擬引擎進行適應度（Fitness）評估。例如，調用基於 EnergyPlus 的 Ladybug Tools 進行全氣候熱力學計算，或是調用 OpenFOAM 運行計算流體力學（CFD）流場求解 \citep{sadeghipour2013ladybug}。CFD 為了求解偏微分方程，需要進行嚴格的網格剖分與耗費大量迭代，完成一次模擬往往需要數小時 \citep{xin2025urbangraph}。當 NSGA-II 演算法需要運行數百代、每代包含數百個個體時，總計算時間將達數週乃至數月之多。這種高延遲的計算瓶頸，直接導致現有的最佳化工具僅能在方案設計完成後進行 Retroactive Verification（事後認識），而難以在設計初期的選案迭代階段提供即時的 Performance-driven（性能驅動）決策引導 \citep{aman2023ai,daglier2025developing}。這正是驅使學術界轉向「低延遲、深度學習代理模型」的關鍵技術契機，亦即第~\ref{chapter:intro}章第~\ref{sec:problem}節所述本研究第一個研究缺口之文獻根源。


\section{深度學習應用於微氣候預測之代理模型}
\label{sec:surrogate_review}

\subsection{歐幾里得與非歐幾里得空間資料表徵之範式轉換}
\label{subsec:gnn_paradigm}

為了解決傳統數值模擬引擎的高運算延遲，利用資料驅動（Data-driven）的機器學習技術建構代理模型（Surrogate Model）已成為當代數位建築學與都市資訊學的主流趨勢 \citep{westermann2019surrogate,wu2023data}。代理模型的核心思想是透過機器學習演算法，去擬合高維幾何輸入與高擬真模擬模型輸出之間的複雜映射函數，從而將耗費大量計算的模擬模型時間壓縮至低延遲的推理尺度 \citep{aman2023ai,shan2025physics}。

在早期應用中，研究者多採用傳統的人工神經網路（ANN）或卷積神經網路（CNN）來進行環境性能預測。然而，CNN 的核心運算單元（如固定尺寸的卷積核）本質上是建立在「歐幾里得空間（Euclidean Space）」的資料表徵之上的。這意味著系統必須將都市環境強制簡化為規格化的 2D 像素矩陣或 3D 體素（Voxel）網格。這種像素型表徵在面對不規則、非均質且具有複雜幾何拓撲的都市實際空間時，暴露出顯著的局限性：

\begin{enumerate}
    \item 為了捕捉微小的幾何細節（如建築縮退、細長的人行步道），體素網格必須極度細化，進而引發嚴重的「維度災難（Curse of Dimensionality）」與記憶體儲量的爆炸。
    \item 都市實體（建築物、喬木、街道）在非歐幾里得幾何中分佈稀疏，規則化網格將包含大量無意義的空節點，造成極大的計算資源浪費。
\end{enumerate}

為了突破歐幾里得資料表徵的瓶頸，圖神經網路（Graph Neural Network, GNN）與圖卷積網路（Graph Convolutional Network, GCN）的引入引發了都市空間預測的範式轉換 \citep{hu2022times,oskarsson2023modeling}。GCN 之奠基性架構由 Kipf 與 Welling 提出，其核心思想是將卷積運算之局部感受野概念由規則網格推廣至任意圖結構，透過鄰接矩陣之正規化聚合，使每個節點能以其一階鄰域之特徵加權平均更新自身表示 \citep{kipf2017semi}。GNN 能夠自然而無損地將都市實體抽象化為「非歐幾里得空間」中的圖結構 $\mathcal{G} = (\mathcal{V}, \mathcal{E})$。在此框架下，都市組件（如建築體量、不同高程的空氣節點、喬木樹冠）被編碼為圖中的異質節點（Nodes, $\mathcal{V}$），而實體間的幾何毗鄰、流體聯通性、熱輻射覆蓋等複雜拓撲交互關係，則被精確編碼為圖中的邊（Edges, $\mathcal{E}$）\citep{wu2023data,shan2025physics}。然而，標準 GCN 假設所有邊共享單一種類之聚合權重，無法區分不同物理因果機制（如陰影遮蔽與對流傳導）之訊號貢獻；為此，Schlichtkrull 等人提出關係圖卷積網路（Relational GCN, R-GCN），為每種關係類型維護獨立的可學習投影矩陣，使模型得以在單次訊息傳遞中同時處理多種異質關係語義 \citep{schlichtkrull2018modeling}，此一機制正是本研究第~\ref{subsec:rgcn}節 RGCN 模組之理論基礎。透過圖神經網路的「訊息傳遞機制（Message Passing）」，每個空間節點能夠匯聚周遭鄰近節點的幾何與物理特徵 \citep{shan2025physics}。多項最新研究指出，GNN 在處理具有不規則幾何與高度空間變異的都市街道熱舒適性與建築群能耗預測任務時，其預測精度、一般化能力與對幾何形變的穩健性，都呈現出顯著優於傳統 CNN 與隨機森林等歐幾里得幾何模型的技術優勢 \citep{hu2022times,wu2023data,shan2025physics}。異質圖神經網路與傳統機器學習於都市微氣候建模上之表徵差異，可整合為圖~\ref{fig:lit_gnn_vs_ml} 之對比示意。

\begin{figure}[htbp]
    \centering
    \includegraphics[width=0.92\textwidth]{../urban-thermal-gnn/04_training/figures/fig_lit_gnn_vs_ml.png}
    \caption[異質圖神經網路與傳統歐幾里得機器學習之表徵對比]{異質圖神經網路與傳統歐幾里得機器學習於都市微氣候建模之資料表徵對比：(a) 傳統歐幾里得表徵（規則體素／像素網格 CNN、或扁平特徵向量輸入之隨機森林），多數網格為無意義之空節點；(b) 異質圖表徵，將建物、空氣、喬木編碼為相異節點類型，並以遮蔽、對流、蒸散、語義等異質邊編碼物理因果關係。原創繪製，理論基礎參考 \citet{kipf2017semi}、\citet{schlichtkrull2018modeling} 與 \citet{xin2025urbangraph}。}
    \label{fig:lit_gnn_vs_ml}
\end{figure}


\subsection{物理資訊機器學習：損失函數軟約束與結構歸納偏置之對比}
\label{subsec:piml_review}

儘管純資料驅動的深度學習模型在推理速度上表現優異，但其本質上屬於統計性黑盒模型，極易在未見過的測試資料上輸出違反物理守恆定律的預測結果（例如，局部空氣溫度無故自發升溫、能量不守恆）。為了解決這一致命缺陷，物理資訊機器學習（Physics-Informed Machine Learning, PIML）成為近年來人工智慧與工程科學結合的研究熱點。

物理資訊神經網路（Physics-Informed Neural Networks, PINN）是此領域的代表性架構。傳統 PINN 的實作邏輯是透過對控制物理規律的偏微分方程（Partial Differential Equations, PDEs）——如流體力學中的納維-斯托克斯方程、熱力學中的熱傳導方程——轉化為數學殘差，並作為正則化項（Regularization Terms）直接嵌入到神經網路的損失函數（Loss Function）中。這使得網路在訓練過程中，不僅要最小化資料點的預測誤差，同時也要最小化物理方程的殘差。然而，這種透過損失函數施壓的約束本質上屬於「軟約束（Soft Constraints）」。在多目標多任務最佳化訓練時，軟約束損失項與資料損失項之間往往存在著潛在的梯度衝突（Gradient Conflict），導致模型難以收斂，或在資料稀疏的隅角上依然可能發生違漏 \citep{xin2025urbangraph}。

為了同時兼顧軟約束 PINN 的收斂困境，\cite{xin2025urbangraph} 提出的 UrbanGraph 架構代表了物理資訊深度學習的最新演進方向 \citep{xin2025urbangraph}。UrbanGraph 放棄了傳統依賴損失函數控制物理殘差的作法，改為將物理第一性原理（First Principles）直接轉化為神經網路圖結構中的「結構歸納偏置（Hard Constraints / Structural Inductive Bias）」：

\begin{itemize}
    \item \textbf{純資料驅動 GNN：}空間鄰近矩陣（無固有物理資訊）
    \item \textbf{傳統 PINN：}損失函數 = 資料損失 + $\lambda \times$ 偏微分方程殘差損失（軟約束）
    \item \textbf{UrbanGraph 框架：}顯式因果拓撲 + 硬性因果剪枝（硬性歸納偏置）
\end{itemize}

在該架構中，時間序列與空間流體、輻射的物理因果關係被直接編碼為異質圖結構中的因果拓撲。例如，模型利用「幾何光線追跡（Ray Tracing）」演算法，依據太陽方位角建立「遮蔽邊（Shading Edges）」，將建築節點與其幾何陰影遮蔽的空氣節點進行實質連接；同時，依據風向因數建立「對流邊（Convection Edges）」，精確模擬上游風向的遞傳行為，並 programmatically 執行顯式因果剪枝（Explicit Causal Pruning），將無物理聯繫的無意義拓撲連結予以汰除 \citep{xin2025urbangraph}。

這種將物理因果關係直接嵌入於圖網路拓撲結構中的「硬性約束」範式，相較於 PINN 的軟性損失約束，具有突破性的技術優勢：它大幅降低了神經網路需要隱式學習的決策空間，使模型需要浮點運算量（FLOPs）接近縮減了高達 73.8\%，同時訓練時收斂速度提升了 21\% \citep{xin2025urbangraph}。更為關鍵的是，結構歸納偏置能夠保護代理模型在面對極端未知幾何形態的推理時，其輸出結果仍能維持高度的物理一致性（Physical Consistency），為建築設計的迭代提供了前所實的工程科學保證 \citep{oskarsson2023modeling,xin2025urbangraph}。


\section{都市路徑規劃文獻之侷限與本研究之範式轉移}
\label{sec:pathfinding_review}

\subsection{傳統尋路演算法於都市步行環境之抽象化}
\label{subsec:pathfinding_abstract}

都市公共步行空間不僅是城市交通網絡的蔓延神經，更是行人感知都市微氣候環境、進行實體社交活動的核心物理載體。在都市交通學與圖論（Graph Theory）的視角下，複雜的都市步行環境可以被高度抽象化為由節點與邊構成的數學拓撲圖模型 $\mathcal{G} = (\mathcal{V}, \mathcal{E})$，街道交叉口、建物出入口、廣場節點對應圖中頂點，街道中線、人行步道對應圖中的邊 \citep{hart1968formal}。Dijkstra 演算法作為廣度優先貪婪搜索策略，能穩定求解單源最短路徑問題；A* 演算法在此基礎上引入啟發值估算函數 $h(n)$，以當前節點至終點之空間直線距離加速收斂 \citep{hart1968formal}。然而，此類傳統尋路演算法之權重矩陣設定極為單一，幾乎完全基於空間幾何物理長度，這種以車行效率為核心導向的設計邏輯，在「以人為本都市設計」語境下已顯現顯著的研究空白：行人之路徑選擇行為並非忠實地以幾何最短路徑求解，而是受到沿途環境品質與熱舒適感受閾值等多維度因素之強烈制約。


\subsection{熱舒適導向尋路研究之發展與其結構性限制}
\label{subsec:thermal_comfort_path}

為解決此一落差，近年學術研究嘗試將動態環境物理因素引入尋路演算法之成本函數，建構「熱舒適導向綠疏步行網絡」\citep{coccolo2018thermal,lopezcabeza2025open}：圖中相鄰節點 $i$、$j$ 間之邊權重不再是恆定幾何常數，而是隨太陽軌跡、氣象條件即時波動之動態變數，將路徑段之熱壓力轉譯為行人之「電阻成本」\citep{lopezcabeza2025open}，使 $A^*$ 演算法得以在極端高溫時段自動迴避高熱暴露路段，規劃出幾何距離可能較長、但整體熱舒適表現較佳的「涼適路徑」。

然而，此一單次起訖點路徑求解典範存在一項本研究認為關鍵的結構性限制：既有文獻之評估單元始終是「單一行人的單次行程」，其成本函數僅回應該次查詢的起點與終點，而都市開放空間之設計決策——例如廣場動線配置、建物量體排列、喬木佈設策略——所服務的對象是整體人行動線網絡在一整日各時段的集體熱舒適表現，而非某一特定行程。換言之，既有熱舒適尋路研究解決的是「使用者查詢時的路徑選擇」問題，而非「設計者決策時的動線品質評估」問題，兩者評估尺度並不一致：對單次查詢而言最優的路徑，並不能直接回答「這個街廓的人行動線系統整體是否宜行」此一設計層級問題。此一評估尺度落差，即第~\ref{chapter:intro}章第~\ref{sec:problem}節所述本研究\textbf{第二個研究缺口}之文獻根源，亦構成本研究第~\ref{subsec:from_astar_to_walkway}節之核心立論——本研究並非在既有 $A^*$ 熱舒適尋路框架上做參數精進，而是將評估單元由「單次路徑查詢」上移至「人行動線網絡之時空平均熱暴露」，使其直接可作為第~\ref{sec:nsga2}節多目標最佳化之適應度函數，而非置於最佳化迴圈之外的下游查詢工具。

\chapter{系統分析與設計}
\label{chapter:methodology}

\section{整合運算框架}
\label{sec:framework}

既有的都市微氣候生成設計工作流面臨一個結構性矛盾：高擬真度模擬引擎能提供物理可信的熱壓力場，卻難以在設計迭代所需的時間內提供立即性回饋；而純資料驅動的機器學習代理模型雖可即時推論，卻常缺乏可追溯的地理資訊基礎與物理一致性保證。本研究之整合運算框架回應此一矛盾，由五個依序銜接的核心模組組成，形成一個從地理資訊建構、都市形態模擬、異質圖轉譯、物理資訊深度學習預測，到多目標最佳化與熱舒適人行動線評估的完整閉環流程：

\begin{enumerate}
    \item \textbf{模組一：基本 GIS 資訊建構（真實 OpenStreetMap 建物足跡）}（GIS Foundation Construction）\\
    以環境部智慧城鄉 IoT 微氣象感測網路建立研究場域之座標參考系統與感測器空間分布；並整合 OpenStreetMap 全域建物足跡與 ETH 全球樹冠高度資料（第~\ref{subsec:v4_real_dataset}節），取代早期單一依賴 NLSC 圖資來源之作法，支援 300 個真實場域之大規模自動化篩選與擷取，作為後續模擬與驗證之地理錨點。
    \item \textbf{模組二：都市形態模擬}（Urban Morphology Simulation）\\
    以 Honeybee／Radiance／EnergyPlus 於 Rhino Grasshopper 環境中執行全場域微氣候逐時模擬，模擬對象為真實擷取之都市量體幾何輸入，取得作為機器學習訓練監督信號之熱壓力目標值（V5：官方光線追蹤與熱傳求解引擎，見第~\ref{sec:ladybug_workflow}節）。
    \item \textbf{模組三：運算框架 —— 真實都市建物足跡與樓高擷取}（Real-World Building Footprint \& Height Extraction）\\
    自 NLSC 多重圖資服務融合擷取單一場域之高精度真實建物輪廓與樓層資訊，取代傳統仰賴蒙地卡羅隨機採樣之單一隨機幾何生成方式，與模組一之 OSM 大規模擷取互補，共同使模型驗證得以錨定於實際都市紋理。
    \item \textbf{模組四：異質圖建構與 PI-ST-GNN 推理}（Heterogeneous Graph Construction \& Inference）\\
    將場景幾何轉譯為包含建物節點與空氣評估節點之異質圖，並以物理資訊時空圖神經網路——Physics-Informed Spatio-Temporal Graph Neural Network（PI-ST-GNN，全名於此首次完整標示，後續章節均沿用此縮寫）——低延遲預測全場域 UTCI 熱壓力場。
    \item \textbf{模組五：Grasshopper 介面整合與多目標最佳化}（Multi-Objective Optimization \& Thermal Comfort Walkway Design）\\
    以訓練完成的 GNN 作為代理模型嵌入 NSGA-II 演算法，同步評估都市形態之熱舒適人行動線品質，透過 FastAPI／WebSocket 後端與 Rhino Grasshopper 介面進行即時通訊，支援互動式閉環設計決策。
\end{enumerate}

圖~\ref{fig:module_flow} 以模組方塊與資訊流箭頭呈現上述五模組之整合架構：模組一（真實 OSM 建物足跡＋IoT）與模組三（NLSC 高精度單場域擷取）為兩條互補之地理資訊輸入路徑，共同匯入模組二之微氣候模擬；模組二輸出之熱壓力標籤驅動模組四之異質圖建構與 PI-ST-GNN 訓練／推理；模組四訓練完成之代理模型最終嵌入模組五，於 Grasshopper 介面內支援 NSGA-II 多目標最佳化與熱舒適人行動線設計之即時回饋。

\begin{figure}[htbp]
    \centering
    \begin{tikzpicture}[
        node distance=8mm and 10mm,
        box/.style={rectangle, rounded corners=2pt, draw=black, thin,
            minimum width=3.05cm, minimum height=1.35cm, align=center,
            font=\small, text=black, fill=white},
        m1/.style={box},
        m2/.style={box},
        m3/.style={box},
        m4/.style={box},
        m5/.style={box},
        lbl/.style={font=\scriptsize\itshape, text=black!55},
        arr/.style={-{Stealth[length=2.6mm]}, thick, draw=black!70},
        darr/.style={-{Stealth[length=2.6mm]}, thick, dashed, draw=black!55}
    ]
        % row 1: two complementary GIS inputs
        \node[m1] (m1) {模組一\\基本 GIS 資訊建構\\{\scriptsize 真實 OSM＋IoT}};
        \node[m3, right=of m1] (m3) {模組三\\真實建物擷取\\{\scriptsize NLSC 三服務融合}};

        % row 2: simulation, converges from both
        \node[m2] (m2) at ($(m1)!0.5!(m3) + (0,-2.1)$) {模組二\\都市形態模擬\\{\scriptsize Honeybee／Radiance／EnergyPlus}};

        % row 3: graph + model
        \node[m4, below=14mm of m2] (m4) {模組四\\異質圖建構與\\PI-ST-GNN 推理};

        % row 4: optimization / design
        \node[m5, below=14mm of m4] (m5) {模組五\\Grasshopper 介面整合\\多目標最佳化};

        \draw[arr] (m1) -- (m2) node[midway, left=1mm, lbl] {座標／感測校準};
        \draw[arr] (m3) -- (m2) node[midway, right=1mm, lbl] {建物幾何};
        \draw[arr] (m2) -- (m4) node[midway, right=1mm, lbl] {UTCI 監督標籤};
        \draw[arr] (m4) -- (m5) node[midway, right=1mm, lbl] {訓練完成之代理模型};
        \draw[darr] (m5.west) .. controls +(-2.6,0) and +(-2.6,0) .. (m1.west)
            node[midway, left=1mm, lbl, align=center] {設計迭代\\回饋閉環};
    \end{tikzpicture}
    \caption[整合運算框架五模組架構與資訊流]{整合運算框架五模組架構與資訊流：模組一、三為互補之地理資訊輸入路徑，模組二至五依序構成模擬、建模、最佳化之閉環流程（虛線為設計迭代回饋路徑）}
    \label{fig:module_flow}
\end{figure}

圖~\ref{fig:data_pipeline_overview} 進一步以資料架構圖展示上述五個模組間之完整資料流向，並標示各處理階段與本章對應小節之關係。

\begin{figure}[htbp]
    \centering
    \includegraphics[width=\textwidth]{../urban-thermal-gnn/04_training/figures/fig_data_pipeline_overview.png}
    \caption[整合運算框架資料架構圖]{整合運算框架資料架構圖}
    \label{fig:data_pipeline_overview}
\end{figure}

上圖由左上四項資料來源（GIS 建物擷取、IoT 感測器、CWA 氣象站、隨機生成幾何）出發，經都市形態模擬、異質圖建構，匯入 PI-ST-GNN 推理引擎，最終驅動下方三項設計決策輸出（NSGA-II 最佳化、熱舒適人行動線設計、Grasshopper 介面回饋），色塊對應本章各小節。

\subsection{設計動機與限制條件}
\label{subsec:framework_motivation}

上述五模組之劃分並非任意的工程切割，而是直接對應第~\ref{chapter:intro}章第~\ref{sec:problem}節所揭示之三項研究缺口與其各自之限制條件。模組一、二、三共同回應「資料是否錨定於真實都市紋理」此一限制條件——若訓練資料仍停留於程序化隨機生成之封閉解空間，模型即便在自身測試集上取得高決定係數，仍難以回答「這個模型是否適用於真實案址」之根本問題，此即為何本研究於真實場景資料集（第~\ref{sec:v4_validation}節）中將建物幾何、樹冠、氣象輸入與微氣候模擬引擎全面替換為 OpenStreetMap、ETH 全球樹冠高度、環境部 IoT 逐時實測資料，以及 Radiance／EnergyPlus 光線追蹤與熱傳求解（V5），而非僅止於程序化模擬或近似式引擎的自我驗證。模組四回應「單次推理延遲是否落在設計迭代可承受之量級」此一限制條件，這是決定代理模型能否被實際嵌入 NSGA-II 演算法內迴圈、而非僅作為離線批次分析工具的關鍵前提；若代理模型之單次推理仍需以分鐘計，則第~\ref{sec:nsga2}節所述之數十代演化搜尋將重新落入與高擬真度模擬引擎相同的運算瓶頸，代理模型之存在意義將大打折扣。模組五回應「運算結果是否能實際回饋至建築設計決策」此一限制條件——一個精度再高的預測模型，若缺乏與 Rhino/Grasshopper 等建築師慣用設計環境的介面整合，其產出便難以直接進入方案發展流程；此一限制條件正是本研究選擇以 FastAPI／WebSocket 即時通訊、而非離線批次匯出報表作為系統輸出形式的核心理由。

\subsection{模組化之必要性論證}
\label{subsec:framework_modularity}

本研究刻意採用模組化架構，而非將地理資訊擷取、模擬取樣與模型訓練整合為單一不可拆解的端到端流程，原因有三。其一，\textbf{資料來源之可替換性}：模組一至三所處理之地理資訊來源本質上具有高度異質性與時效性——政府開放圖資之圖層版本會定期更新、OSM 標籤完整度因地區而異、IoT 測站布建範圍逐年擴增——若將資料擷取邏輯與下游模型訓練緊密耦合，任何資料來源之變動都將牽動整個系統的重新設計；模組化設計使模組三之建物擷取方法（OSM 批次擷取或三服務融合向量化）能被獨立替換或並行使用，而不影響模組四、五之圖神經網路架構與最佳化邏輯。其二，\textbf{驗證責任之可分離性}：模組二（都市形態模擬）之輸出品質決定了模組四訓練標籤之上限，若二者未被清楚劃分為獨立模組，模型預測誤差將難以歸因於「模擬引擎本身之物理近似限制」或「圖神經網路之學習能力不足」，此一可歸因性正是第~\ref{sec:limitations}節局限性討論得以逐條追溯原因、而非籠統歸咎於「模型精度不足」的方法論基礎。其三，\textbf{部署環境之可攜性}：模組五之 FastAPI／WebSocket 後端與 Grasshopper 前端刻意與模組一至四之訓練流程解耦，使得同一套已訓練完成之代理模型檔案，能在不同硬體環境（訓練用 GPU 工作站與設計端筆記型電腦）之間直接搬遷部署，訓練與推理兩端無需共用相同的資料前處理程式碼路徑。

\subsection{PI-ST-GNN 全名與可解釋性定位}
\label{subsec:framework_interpretability}

PI-ST-GNN（Physics-Informed Spatio-Temporal Graph Neural Network，物理資訊時空圖神經網路）此一命名之三個組成部分，分別對應模型設計中三項刻意保留的可解釋性錨點，而非僅為技術堆疊之標籤。「Physics-Informed」對應第~\ref{sec:physics_loss}節之三項物理資訊軟約束損失（輻射一致性、時序平滑性、風場阻滯），其設計動機並非單純追求數值精度之提升，而是確保模型輸出可被人類專家以物理常識直接檢核——例如陰影區之預測溫度不應高於日照區，此一可檢核性使建築師與口試委員得以在不理解模型內部權重的前提下，仍能對模型輸出之合理性進行判斷。「Spatio-Temporal」對應模型同時處理空間關係型圖卷積（第~\ref{subsec:rgcn}節 RGCN）與時序演化（第~\ref{subsec:lstm}節 LSTM 暖啟動機制）雙重結構，使模型輸出之全日 11 時段熱壓力場彼此之間具有物理上合理的連續性，而非各時步獨立、可能出現突兀跳動的離散預測。「Graph Neural Network」則對應第~\ref{sec:graph_construction}節以異質圖顯式編碼建物與空氣節點間之遮蔽、對流、蒸散、語義、毗鄰五類因果關係邊，使模型之訊息傳遞路徑本身即具備可追溯性——第~\ref{subsec:trustworthiness}節之 GNNExplainer 分析，即是沿著這些顯式邊類型回溯特定預測值之主要貢獻來源，此一可解釋性分析路徑，在純像素卷積或全連接黑盒模型中並不存在對應的結構基礎。

\subsection{與 UrbanGraph 之關係}
\label{subsec:framework_urbangraph_relation}

本研究之 RGCN 關係型聚合公式與 LSTM 暖啟動機制，係採用並延伸自 \citet{xin2025urbangraph} 提出之 UrbanGraph 架構（詳見第~\ref{subsec:rgcn}、\ref{subsec:lstm}節之公式推導），其異質圖之因果拓撲建構邏輯（第~\ref{sec:graph_construction}節之遮蔽邊、對流邊建構規則）亦直接沿用 UrbanGraph 之硬性因果剪枝設計。本研究在此一既有架構基礎上之延伸具體有三：其一，將原架構單一節點類型之設計顯式區分為建物與空氣兩種異質節點類型，並以獨立編碼器分別處理其靜態幾何與動態氣象特徵；其二，於原有圖拓撲硬約束之外，另行引入三項物理資訊軟約束損失作為訓練過程之補充監督機制；其三，將原本可能侷限於單一時步預測之架構，擴展為一次輸出全日 11 時段 UTCI 場之多步驟時序預測，並整合至第~\ref{chapter:results}章之端到端低延遲推理與設計迴圈。上述沿用與延伸之分界，於第~\ref{sec:contributions}節研究貢獻中將以逐項對照方式再次明確標示，避免讀者將既有架構之既有能力誤認為本研究之原創貢獻。

\subsection{設計原則小結}
\label{subsec:framework_principles}

綜合上述四小節，本研究整合運算框架之設計可歸納為三項一貫原則：（一）\textbf{真實性優先於規模}——寧可縮小訓練場景數量（V5 之 300 個真實場景），也優先確保幾何、氣象與物理模擬引擎均錨定於真實世界，而非以程序化生成手段換取更大但脫離真實紋理的資料規模；（二）\textbf{延遲預算驅動架構選型}——模組四之圖神經網路架構選型（RGCN 而非全連接網路、代理模型而非直接呼叫模擬引擎）之根本考量始終是能否落在模組五 NSGA-II 演化搜尋所能負擔之延遲預算內，而非單純追求排行榜式的精度極大化；（三）\textbf{可解釋性作為架構約束、而非事後補充}——第~\ref{subsec:framework_interpretability}節所述之物理資訊損失與顯式異質邊設計，於架構設計階段即已納入考量，使第~\ref{subsec:trustworthiness}節之可解釋性分析得以有明確的結構基礎可供回溯，而非在模型訓練完成後才回頭嘗試以事後方法（post-hoc）解釋一個純黑盒模型的行為。


\section{模組一：基本 GIS 資訊建構}
\label{sec:gis_foundation}

本研究以政府開放圖資作為建構研究場域地理資訊基礎之資料來源，在不依賴商用測繪資料授權的前提下，取得具備樓層屬性的建物量體資料，並建立一套可重現的 GIS 資訊擷取流程。圖~\ref{fig:site_context} 由國家尺度逐層放大至研究場域尺度，說明研究場域在真實地理脈絡中的定位。

\begin{figure}[htbp]
    \centering
    \includegraphics[width=\textwidth]{../urban-thermal-gnn/04_training/figures/fig_geo1_site_context.png}
    \caption[研究場域地理脈絡]{研究場域地理脈絡：由國家尺度逐層放大至場域尺度}
    \label{fig:site_context}
\end{figure}

上圖 (a) 為台灣全島尺度，紅框標示真實場景資料集之擴大研究區域、紅星標示陽明交大光復校區錨點；(b) 為擴大研究區域之真實 IoT 感測器（藍點，2,018 站）與 300 個選定真實訓練場景（綠點）之空間分布；(c) 為一個真實 $80\times80$ 公尺場景之幾何：真實 OSM 建物量體、ETH 樹冠與空氣評估節點格點——取代早期版本之程序化示意場景，直接以真實地理幾何呈現場域尺度。

\subsection{建物圖資取用流程與場景篩選演算法}
\label{subsec:crs_sources}

場景幾何統一以 WGS84 地理座標（EPSG:4326）為基準，並於各場景局部運算時轉換為以場景中心為原點之公尺制平面座標。訓練資料所需之建物圖資面臨一核心約束：須在數百個場景之區域尺度上批次、可重現地取得真實建物足跡與樓高，而互動式測繪服務（如國土測繪中心 WMTS／WMS，須逐圖磚人工套疊判讀）難以支撐此一取樣規模。為此，本研究之真實場景資料集改以 OpenStreetMap（OSM）為建物與地表圖資來源，其離線區域擷取機制與樓高判讀邏輯詳見第~\ref{sec:building_extraction}節；本節聚焦於場景候選之篩選演算法。

場景候選之產生並非於區域內任意佈點，而是以真實環境部 IoT 測站之座標為中心，使每一候選場景中心恰為一真實測站；候選經最小間距 150 公尺之貪婪去重，確保 $80\times80$ 公尺場景互不重疊。篩選演算法對每一候選逐一檢核四項條件（真實建物、地表異質性、真實 IoT 站點、真實樹冠，詳見第~\ref{subsec:v4_real_dataset}節），全數通過者再依綜合品質評分（真實樓高建物數、IoT 密度、樹冠像元數、建物密度加權）排序，取前 300 名構成訓練場景集。此一「以真實測站為中心、以四項條件過濾、以品質評分排序」之演算法，使取樣同時滿足感測器校準之近接性與真實幾何之完整性，其結果之空間分布見第~\ref{subsec:v4_real_dataset}節。


\subsection{IoT 微氣象感測器空間資料}
\label{subsec:iot_spatial}

除靜態建物圖資外，本研究整合環境部智慧城鄉空品微型感測器網路之即時觀測資料，作為模擬邊界條件校準、氣象驅動與研究場域選定之依據。為使真實場景資料集（第~\ref{subsec:v4_real_dataset}節）達科學可信之取樣數，感測器資料涵蓋範圍由早期版本之新竹市 15 公里圈擴大至以陽明交大光復校區為錨點、涵蓋新竹市／新竹縣／桃園市／苗栗縣／臺中市之研究區域，區域內共納入 2,018 站真實感測器，其空間分布與逐時量測特性如圖~\ref{fig:iot_sensor_map} 所示。

\begin{figure}[htbp]
    \centering
    \includegraphics[width=\textwidth]{../urban-thermal-gnn/04_training/figures/fig_iot_sensor_map.png}
    \caption[IoT 微型感測器空間分布實測圖]{真實場景資料集擴大研究區域之環境部 IoT 微型感測器空間分布實測圖（區域內逐時觀測，2025 年夏季正午時段）}
    \label{fig:iot_sensor_map}
\end{figure}

上圖 (a) 為空氣溫度分布，色階由綠（低溫）至紅（高溫），可見都市化核心與丘陵地帶顯著溫差；(b) 為相對濕度分布，沿海與內陸呈現梯度差異；紅色星號標示研究區域錨點（陽明交大光復校區）。

在感測器密度分析階段，本研究以感測器站點之空間座標本身為對象，並以核密度估計（Kernel Density Estimation, KDE）評估其空間分布密度，用以界定真實場景之候選取樣區域。此一感測器密度分析與最終選定之 300 個真實場景之對照，如圖~\ref{fig:studyarea} 所示。

\begin{figure}[htbp]
    \centering
    \includegraphics[width=0.95\textwidth]{../urban-thermal-gnn/04_training/figures/fig_studyarea_sensor_density.png}
    \caption[研究範圍選定]{真實場景資料集研究範圍選定：擴大區域之 IoT 感測器密度曲面與 300 個選定真實場景對照}
    \label{fig:studyarea}
\end{figure}

上圖 (a) 為擴大研究區域內之核密度估計曲面，青點為 2,018 站真實感測器、綠點為 300 個選定之真實訓練場景、紅星為錨點；(b) 為錨點核心區之局部放大，呈現選定場景與周邊真實感測器之空間對應關係。


\subsection{全球樹冠高度資料}
\label{subsec:canopy_data}

建物與感測器圖資之外，喬木樹冠高度是異質圖中 \texttt{veg\_et}（植被蒸散）邊緣與空氣節點 $T_h$（最近樹木高度影響）特徵的真實資料來源。本研究採用 ETH Zurich 發布之全球樹冠高度資料集 ETH GlobalCanopyHeight（10 公尺解析度、2020 年版本，WGS84 地理座標）\citep{lang2023high}，該資料集以 Sentinel-2 多光譜影像結合 GEDI 星載光達（Spaceborne LiDAR）訓練深度學習迴歸模型反演逐像元樹冠高度，覆蓋全球陸地表面。研究場域所在圖磚（N24E120）之空間範圍與研究場域錨點之對照如圖~\ref{fig:canopy_map} 所示。

\begin{figure}[htbp]
    \centering
    \includegraphics[width=\textwidth]{../urban-thermal-gnn/04_training/figures/fig_canopy_height_map.png}
    \caption[樹冠高度資料與研究場域對照]{ETH GlobalCanopyHeight 樹冠高度資料與研究場域對照}
    \label{fig:canopy_map}
\end{figure}

上圖 (a) 為研究場域周邊 15 公里範圍內之樹冠高度分布，疊繪於真實衛星影像底圖之上，色階採與 ETH 官方線上展示工具一致之 inferno 色彩尺度（深紫至亮橙代表樹冠高度遞增），虛線圓圈與圖~\ref{fig:iot_sensor_map} 採相同之 15 公里範圍，供跨圖對照；灰階（無色階疊加）區域代表該處無樹冠訊號（NoData），對應道路、建物等不透水鋪面。可見新竹沿海丘陵地帶具有大面積連續天然林覆蓋（樹冠高度可達 30 公尺以上），與都市化核心區之低樹冠覆蓋形成鮮明對比。(b) 為場域局部放大，青框標示 $80\times80$ 公尺研究場域之實際地理位置與尺度；可見研究場域本身座落於高密度建成核心，該 $80\times80$ 公尺範圍內之像元均為無樹冠訊號，但周邊數十至數百公尺外仍可見明顯樹冠覆蓋。

在資料整合流程上，本研究依循 sensing\_integration/canopy\_loader.py 之取樣邏輯：對每一棵樹木節點之場域局部座標 $(x, y)$，先轉換為 WGS84 經緯度，再以最近像元法（Nearest-neighbor Sampling）自 GeoTIFF 讀取對應樹冠高度值；若該位置存在有效觀測值（高度 $\ge 1$ 公尺），則以此真實觀測值覆蓋原先由隨機幾何採樣器生成之樹木高度與推算樹冠半徑，並標記資料來源為 \texttt{canopy\_map}；若該位置無有效觀測，則保留隨機生成之高度值，標記為 \texttt{random\_fallback}。此一「真實資料優先、程序化生成為輔」的分層取樣策略，使異質圖之樹木節點特徵在有真實觀測依據之處錨定於實測數據。


\subsection{真實場景訓練資料集}
\label{subsec:v4_real_dataset}

為使模型訓練直接立基於真實都市地理資訊而非全然程序化生成之幾何，本研究於前述 GIS 圖資基礎上，建立一套完全以真實資料驅動之 300 場景訓練資料集。此一真實場景之選址、地理資料與氣象邊界條件，同時是本研究最終採用之 V5 主要模型與其前身（V1--V4，見第~\ref{sec:v4_validation}節）共用之基礎——僅下游之微氣候模擬引擎（第~\ref{sec:ladybug_workflow}節）與異質圖動態邊建構方法（第~\ref{subsec:edges}節）有所差異，場址、建物、樹冠與 IoT 氣象等地理資料本身並無版本間差異。其選址以陽明交通大學光復校區為錨點，為達科學可信之取樣數而將搜尋範圍擴及新竹市、新竹縣、桃園市、苗栗縣與臺中市；候選場址直接取自環境部 IoT 測站之真實座標（去重至最小間距 150 公尺，避免 $80\times80$ 公尺場域重疊），並逐一檢核以下四項條件，僅保留同時滿足全部條件者：

\begin{enumerate}
    \item \textbf{真實建物幾何（陰影訊息）：}場域 120 公尺陰影脈絡範圍內至少含 3 棟真實 OSM 建物足跡，且至少 1 棟位於場域中心 70 公尺內，確保場景具真實建物量體與陰影結構；
    \item \textbf{地表覆蓋異質性：}場域周邊具建成（built）以外之至少一類地面材質（草地、水體、裸土、鋪面等，取自 OSM \texttt{landuse}／\texttt{natural}／\texttt{leisure}），合計 $\ge 2$ 類物理表面，提供反射率與蓄熱異質性；
    \item \textbf{高密度真實 IoT：}場域中心即為一真實環境部微型感測器站點（最近測站距離 0 公尺），可直接以該站真實觀測驅動與校準；
    \item \textbf{真實植栽陰影：}場域 $80\times80$ 公尺範圍內具 ETH GlobalCanopyHeight 有效樹冠像元，確保含真實植栽遮蔽訊息。
\end{enumerate}

全區共 304 個候選場址通過全部條件，依綜合品質評分（真實樓高建物數、IoT 密度、樹冠像元、建物密度加權）取前 300 名，其於真實地圖上之分布如圖~\ref{fig:selected_sites} 所示。300 場景依序輪派至 6、7、8 三個月份（各 100 場景），使資料集涵蓋整個實測夏季。各項資料真實度說明如下：建物\textbf{足跡}為 100\% 真實（每場景平均 11.3 棟，直接來自 OSM）；建物\textbf{樓高}則因臺灣 OSM 高度標籤稀疏，僅約 14\% 建物具真實 \texttt{height}／\texttt{building:levels} 標籤，其餘以三層（10.8 公尺）之樓層預設值回填，屬本資料集之已知限制；喬木高度取自 ETH 樹冠實測；空氣溫度與相對濕度則為\textbf{全部 300 場景}各自所在 IoT 測站之真實逐時觀測（溫、濕度覆蓋率均為 300/300）。

\begin{figure}[htbp]
    \centering
    \includegraphics[width=\textwidth]{../urban-thermal-gnn/04_training/figures/fig_selected_real_sites.png}
    \caption[真實訓練場景選取結果]{真實訓練場景選取結果：(a) 300 個場景於擴大搜尋範圍（新竹／桃園／苗栗／臺中，錨定陽明交大光復校區）之區域分布，點色對應綜合選址評分；(b) 錨點核心區之場景，$80\times80$ 公尺場域邊界按真實比例繪製於 OpenStreetMap 底圖之上。每一場景同時滿足真實建物、地表異質性、真實 IoT 站點與真實樹冠四項條件。}
    \label{fig:selected_sites}
\end{figure}

為呈現訓練資料所涵蓋之真實都市幾何多樣性，圖~\ref{fig:scene_array} 自 300 個場景中取 25 個清晰易讀者以 $5\times5$ 陣列並列，各場景之建物量體均為真實 OSM 足跡（色階對應真實簷高），綠點為 ETH 樹冠實測植栽；可見場景涵蓋高密度建成、街廓邊緣、疏落郊區等不同都市紋理與棟距分佈。

\begin{figure}[htbp]
    \centering
    \includegraphics[width=0.96\textwidth]{../urban-thermal-gnn/04_training/figures/fig_scene_array_25.png}
    \caption[真實訓練場景代表性樣本]{300 個真實訓練場景中之 25 個代表性場景（$5\times5$ 陣列）：真實 OSM 建物足跡（色階＝簷高）與 ETH 樹冠。各子圖標題示場景編號、建物棟數與所屬月份。}
    \label{fig:scene_array}
\end{figure}

圖~\ref{fig:scene_array} 所示之靜態幾何多樣性，於逐時模擬後進一步轉譯為隨時間演變之熱壓力場。圖~\ref{fig:heatmap_grid} 取其中 9 個場景（依喬木數量遞增排序），以 V5 Radiance／EnergyPlus 模擬引擎（第~\ref{sec:ladybug_workflow}節）呈現 08:00 至 18:00 六個代表時步之 UTCI 熱力圖並列結果，各時步之數值皆為該場景所屬月份之模擬代表值，非跨月平均。可見隨喬木數量增加，場域內低溫（藍綠）區域範圍隨之擴大；各場景之熱壓力峰值時段普遍集中於太陽仰角最高之 10:00--14:00，午後（16:00--18:00）則隨太陽輻射減弱而回落——此一逐時演變之空間熱壓力型態，即為第~\ref{sec:graph_construction}節異質圖動態邊（\texttt{shadow}、\texttt{veg\_et}、\texttt{convective}）所需學習之核心監督訊號來源。

\begin{figure}[htbp]
    \centering
    \includegraphics[width=\textwidth]{../urban-thermal-gnn/04_training/figures/fig_v5_heatmap_grid_9x6.png}
    \caption[V5 模擬熱力圖陣列]{V5（Radiance／EnergyPlus）模擬之 9 場景 $\times$ 6 時步 UTCI 熱力圖陣列：場景由上至下依喬木數量遞增排序，各列標示場景編號與建物／喬木棟數；灰色為真實 OSM 建物足跡，綠色圓為 ETH 樹冠實測植栽；各時步之數值皆為該場景所屬月份之模擬代表值，非跨月平均。}
    \label{fig:heatmap_grid}
\end{figure}


\section{模組二：都市形態模擬——物理式微氣候逐時模擬}
\label{sec:ladybug_workflow}

微氣候逐時模擬以 Ladybug Tools 生態系之 Honeybee／Radiance／EnergyPlus 官方 \texttt{utci\_comfort\_map} recipe 實作，逐一評估場景內每個空氣節點之逐時熱壓力目標值 $T_a, \text{MRT}, V_a, \text{RH}, \text{UTCI}$。其運算核心包含四個相互耦合之真實物理求解模組：

\begin{enumerate}
    \item \textbf{天空可視因子（Sky View Factor, SVF）}：以 Radiance \texttt{sky\_view} recipe 執行光線追蹤，自各節點向天空半球發射大量取樣射線，依建物與喬木之三維量體幾何計算實際被遮擋之立體角比例。
    \item \textbf{短波輻射與陰影}：以 Radiance 二階段（Two-phase）方法求解直達、漫射與地面反射三分量短波輻射通量，取代單純的幾何陰影布林判定，逐時反映太陽位置、大氣條件與周邊量體之真實光線遮蔽關係。
    \item \textbf{戶外平均輻射溫度（MRT）}：短波分量取自上述 Radiance 求解結果，長波分量則以 EnergyPlus 求解建物表面溫度後，透過 \texttt{ladybug-comfort} 之 \texttt{OutdoorSolarCal} 合成，為真實建物外殼熱傳導與蓄熱行為之直接輸出，而非經驗式近似。
    \item \textbf{風場}：現行版本以每一場景單一之參考風速代表全場域風況（見下段殘餘限制說明），尚未逐點求解街谷尺度風場遮蔽。
\end{enumerate}

UTCI 直接取自 \texttt{utci\_comfort\_map} recipe 之原生多項式輸出（不再另行呼叫 \texttt{pythermalcomfort} 套件重新合成）。模擬時段涵蓋 08:00 至 18:00 共 11 個整點時步（$T = 11$），捕捉日間熱壓力最嚴峻之逐時變化軌跡。

需說明者：本研究之微氣候模擬引擎歷經 V1--V4 之自製 Python 物理近似模型（幾何射線近似 SVF 與陰影、地表能量平衡經驗式 MRT，\textbf{未}呼叫任何外部模擬套件），至 V5 起正式改採上述 Radiance 6.2.1（光線追蹤）與 EnergyPlus 25.1.0（建物熱傳求解）之官方模擬流程——此一轉換之完整比較與 V1--V4 近似引擎之瑕疵揭露，見第~\ref{sec:v4_validation}節與附錄~\ref{appx:dynamic_edges_fix}。V5 現行版本仍有以下已知殘餘限制：（一）風場以每場景單一參考風速代表，未逐點求解街谷尺度繞流；（二）建物外殼構造採 ASHRAE 標準構造假設，而非逐棟真實材料實測；（三）陰影判定之來源索引（哪一建物／喬木造成特定節點陰影）為幾何後處理之推導代理值，非 recipe 原生輸出。上述限制均不影響 UTCI 本身之物理求解性質，但為後續研究可精進之方向，詳見第~\ref{sec:limitations}節。

氣象邊界條件於第~\ref{subsec:v4_real_dataset}節之真實場景資料集中，直接採用各場景所在地之環境部 IoT 測站真實逐時氣溫與相對濕度（經合理性驗證，故障測站退回中央氣象署區域觀測）、中央氣象署測站真實風速，以及區域 TMY（Typical Meteorological Year）EPW 之實測日射剖面（因 CWA 該測站未提供日射觀測，故日射採區域典型氣象年之實測輻射，而非場景當日觀測），並由 \texttt{13\_build\_scenario\_epw\_v5.py} 逐場景合成對應之客製化 EPW 檔案供 Radiance／EnergyPlus 讀取；此一設計使模擬之熱壓力目標值錨定於真實觀測氣候。

此一模擬流程對每個場景輸出逐時之全場域 UTCI，即為機器學習之監督標籤。圖~\ref{fig:timestep_utci} 以四個真實場景示範此一熱舒適訓練資料之生成：各列為一場景，各行為 08:00 至 18:00 之代表時步，節點顏色為 Radiance／EnergyPlus 模擬之逐時 UTCI。可清楚觀察到熱壓力於日間之累積與消退（清晨低、正午高、傍晚回落），以及真實建物與植栽在各時步所投射之陰影所形成之持續性空間結構——此即時空圖神經網路所需學習之「幾何—時序—熱壓力」映射關係之資料來源。

\begin{figure}[htbp]
    \centering
    \includegraphics[width=\textwidth]{../urban-thermal-gnn/04_training/figures/fig_v5_timestep_utci.png}
    \caption[熱舒適訓練資料之逐時生成]{熱舒適訓練資料之逐時生成（V5，Radiance／EnergyPlus 模擬）：真實場景於 08:00--18:00 各代表時步之模擬 UTCI 場。展現熱壓力之日間動態與真實幾何遮蔽所致之空間梯度。}
    \label{fig:timestep_utci}
\end{figure}

此一模擬工作流之輸出，除作為機器學習訓練之監督信號外，亦是第~\ref{sec:nsga2}節熱舒適人行動線評估與多目標最佳化之基準比對依據。


\section{模組三：真實都市建物足跡與樓高擷取——開放街圖（OSM）真實建物與地表鋪面圖資擷取}
\label{sec:building_extraction}

\subsection{既有隨機生成方法之侷限}
\label{subsec:parametric_limitation}

模型訓練所需的大規模場景資料集，傳統上由隨機幾何採樣器（Stochastic Geometry Sampler）以蒙地卡羅法在容積率（FAR）、建蔽率（BCR）與縮退規範之法規約束下生成：每個場景包含 2--5 棟建築（矩形或 L 形平面、面寬與進深各 5--25m、樓層數 2--12 層）與 2--5 棵喬木（高度 4--12m），共產生 300 個合法場景。此方法能有效覆蓋法規允許之形態解空間，但其產出之建物量體、棟數關係、街廓紋理與地表覆蓋均為程序化生成，與真實都市紋理之棟距、街廓尺度、樓高分布與鋪面異質性統計特性未必一致，使模型於真實場地應用時的一般化可信度難以驗證。因此，本研究以開放街圖（OpenStreetMap, OSM）之真實建物足跡、樓高與地表覆蓋圖資，全面取代訓練場景之程序化幾何輸入，詳見以下兩小節；既有隨機生成方法則保留作為第~\ref{sec:v4_validation}節之對照資料集。圖~\ref{fig:scenarios} 示例六個典型隨機生成場景，可見其建築量體雖滿足法規約束，但棟距、朝向與棟數關係呈現均質化之程序化特徵，與圖~\ref{fig:scene_array} 之真實場景紋理形成對照。

\begin{figure}[htbp]
    \centering
    \includegraphics[width=\textwidth]{../urban-thermal-gnn/04_training/figures/fig_geo4_scenarios.png}
    \caption[隨機生成場景示意]{六個典型隨機生成場景之建築量體與喬木配置示意}
    \label{fig:scenarios}
\end{figure}


\subsection{OSM 離線區域建物足跡與樓高擷取}
\label{subsec:osm_building_extraction}

以線上 Overpass API 逐場景查詢建物圖資，在 300 個場景之取樣規模下面臨兩項限制：公開服務之流量限制使查詢頻繁逾時或遭拒，且每次查詢之網路延遲使批次取樣耗時不可控。本研究因此改採離線區域擷取：一次性下載 Geofabrik 釋出之臺灣區域 \texttt{.osm.pbf} 封裝檔，以 pyosmium 函式庫解析全部建物與地表覆蓋要素於研究區域邊界框內者，建立含 STRtree 空間索引之常駐記憶體資料結構；此後任一場景之建物與材質查詢即為純記憶體之空間索引查詢，無需再次連線。

建物幾何之萃取僅處理封閉方式（Closed Way）——即端點座標相同、由 OSM 節點座標依序連接而成之多邊形——此涵蓋研究區域內絕大多數建物；由關係（Relation）組成之少數大型建物複合體多邊形未納入處理範圍，如實計入未擷取項目，而非以近似輪廓填補。建物樓高之判讀依序嘗試 \texttt{height}、\texttt{building:height} 標籤，其次以 \texttt{building:levels} 樓層數乘以標準層高（3.6 公尺）推算；當標籤含分號分隔之多值（如合併建物之 \texttt{building:levels="11;10;12"}，反映該建物實際由不同樓高之區塊合併登錄）時，取各值之平均。當上述來源皆缺失時，以三層樓（10.8 公尺）之預設值回填，並以布林旗標標記該筆高度之真實性——僅當高度直接來自 OSM 標籤時計為真實，回填值不計入；此一標記機制使第~\ref{subsec:v4_real_dataset}節得以逐場景統計真實樓高之比例。


\subsection{OSM 地表鋪面與材質圖資擷取}
\label{subsec:osm_material_extraction}

建物幾何之外，地表材質之空間異質性直接決定各節點之反射率（Albedo）與蓄熱特性，是都市微氣候模擬不可或缺的邊界條件。本研究以 OSM 之 \texttt{landuse}、\texttt{natural}、\texttt{leisure} 三類標籤，將地表覆蓋歸類為草地／植被、裸土、鋪面／混凝土、瀝青／工業、水域五種物理材質：\texttt{landuse=grass／meadow／park／forest} 等歸為草地，\texttt{landuse=residential／commercial／retail} 等歸為鋪面，\texttt{landuse=industrial} 歸為瀝青，\texttt{natural=water} 與 \texttt{leisure=swimming\_pool} 等歸為水域，\texttt{landuse=farmland／construction} 等歸為裸土；此一標籤對應表直接沿用第~\ref{subsec:crs_sources}節之 OSM 離線區域索引，與建物擷取共用同一次區域資料載入。

每一材質多邊形依此對應表轉換為地表覆蓋代碼，並以 1 公尺網格柵格化為場景之 \texttt{land\_cover\_map}：地面材質先行填入，建物足跡範圍隨後覆蓋為「建物（built）」代碼，反映屋頂材質優先於其下方地面材質之物理邏輯。此一代碼直接對應第~\ref{sec:ladybug_workflow}節模擬引擎之反射率查表（建物／鋪面 0.30、裸土 0.40、水域 0.00、草地 0.20），取代既有版本中僅以單一場地平均反射率參數近似地表異質性之作法。圖~\ref{fig:material_zones} 以一個真實場景示範此一取用流程：(a) 為 OSM 原始材質多邊形與建物足跡疊圖；(b) 為柵格化後實際餵入模擬引擎之 \texttt{land\_cover\_map}。

\begin{figure}[htbp]
    \centering
    \includegraphics[width=\textwidth]{../urban-thermal-gnn/04_training/figures/fig_material_zones.png}
    \caption[OSM 地表鋪面與材質圖資擷取]{OSM 地表鋪面與材質圖資擷取：(a) 原始材質多邊形與建物足跡；(b) 柵格化 land\_cover\_map。}
    \label{fig:material_zones}
\end{figure}

須說明者：OSM 之 \texttt{landuse} 標註並非對地表覆蓋之窮盡標記，如圖~\ref{fig:material_zones}(b) 所示，場景內相當比例之網格未落於任何材質多邊形範圍內（多對應道路、未標註之鋪面或私有地面）；此類未標註網格於模擬時退回以場域平均反射率參數近似，而非誤植為特定材質，此一取捨保守地避免了對未知地表覆蓋之武斷分類。


\subsection{三服務融合建物向量化方法（單一場地案例驗證用途）}
\label{subsec:vectorize_method}

OSM 離線區域擷取雖能支撐 300 場景規模之批次訓練資料建構，其建物樓高標籤覆蓋率有限；對於需要高精度樓高驗證之單一場地案例研究（第~\ref{chapter:results}章），本研究另建立一套以政府開放圖資為來源之三服務融合建物向量化方法，作為 OSM 之互補驗證來源。真實建物量體之擷取面臨一項核心限制：單一圖資來源難以同時提供精確幾何輪廓與樓層屬性。TOPO01K 圖層雖標註樓高文字，但其柵格化 hatch 填充在高縮放層級下受制圖線寬影響而破碎為大量細小連通元件，直接以連通元件分析取得的建物遮罩並不可靠；B5000 與 EMAP 雖幾何品質較佳，卻不含樓層資訊。本研究因此設計一套三服務交叉驗證方法，讓三個各有所長、各有所缺的圖資來源互補：

\begin{enumerate}
    \item \textbf{建物範圍重建：}將 TOPO01K 之粉紅色填色像素與黑色線繪像素（外框、內部 hatch、樓高文字）聯集為單一「建物範圍遮罩」，而非僅取粉紅色填色本身 —— 因為 hatch 線條與填色像素在製圖時緊密交疊，聯集遮罩可還原被 hatch 切割破碎的填色區域至建物尺度之完整連通元件，且無需透過形態學膨脹等會扭曲邊界的操作。
    \item \textbf{粗細線分離：}以形態學開運算（Morphological Opening）對建物範圍內之黑色線像素進行處理，去除短而細之斜向 hatch 筆劃，保留較粗、較連續之樓高分區邊界線，藉此將單一建物範圍細分為對應各自獨立樓高之子區塊。
    \item \textbf{子區塊輪廓萃取：}對每一子區塊以 OpenCV \citep{bradski2000opencv} 實作之邊界追蹤演算法（cv2.\allowbreak findContours，基於 Suzuki \& Abe 之拓樸邊界追蹤法 \citep{suzuki1985topological}）取得次像素封閉輪廓，再以 Douglas--Peucker 演算法 \citep{douglas1973algorithms} 化簡為平滑多邊形，作為該樓高分區之向量化幾何輸出。
    \item \textbf{多字元群集樓高辨識：}掃描每一子區塊內部之黑色文字像素群集（依尺寸與長寬比篩除細長之殘留 hatch 筆劃），對每一候選群集個別以 Tesseract OCR 引擎 \citep{smith2007overview} 執行光學文字辨識，取代僅於整棟建物幾何中心點單次辨識之作法 —— 大型建物複合體內常包含多個並列之不同樓高分區（如 \texttt{6R}、\texttt{5R}、\texttt{4R} 並列），單次中心點辨識幾乎必然錯過實際文字位置。
    \item \textbf{跨服務信心度評估：}將每一子區塊之向量化輪廓分別與 B5000 黑線封閉區域、EMAP 淡紫灰填色區域進行像素重疊率比對，依兩項驗證來源之通過情形標記高／中／低信心度，供後續篩選使用。
\end{enumerate}

圖~\ref{fig:gis_fusion_pipeline} 完整呈現上述三服務融合方法之系統流程，包含三項圖資之像素對齊機制、TOPO01K 主處理鏈之四道步驟，以及 B5000／EMAP 驗證分支如何匯聚為最終信心度分級輸出。

\begin{figure}[htbp]
    \centering
    \includegraphics[width=\textwidth]{../urban-thermal-gnn/04_training/figures/fig_gis_fusion_pipeline.png}
    \caption[三服務融合建物向量化流程]{三服務融合建物向量化系統流程圖}
    \label{fig:gis_fusion_pipeline}
\end{figure}

上圖橘色為 TOPO01K 主幾何與樓高擷取鏈（建物範圍重建→粗細線分離→子區塊輪廓萃取→多字元群集樓高辨識）；深灰、紫色分別為 B5000、EMAP 之獨立驗證分支；綠色為最終輸出，附帶信心度分級與完整屬性表。

樓層文字之判讀依循新竹地區都市計畫現況圖之標示慣例：\texttt{NR} 表示地上第 N 層、\texttt{T} 表示地面層、\texttt{NB} 表示地下第 N 層、\texttt{NM} 表示夾層。當同一子區塊經多字元群集辨識比對後仍檢出一種以上互斥樓層標籤時，判定為分區尚未完全切割成功之案例，如實標記為待覆核而非任意保留其中一筆結果，以維持資料集之可追溯性。


\subsection{三類幾何來源之分工關係}
\label{subsec:extraction_vs_parametric}

本研究之三類幾何來源並非互斥替代，而依各自之取樣規模與精度優勢分工：OSM 離線區域擷取（第~\ref{subsec:osm_building_extraction}節）支撐 300 場景大規模真實訓練資料集，以真實足跡與地表覆蓋換取模型訓練分佈與真實都市紋理之統計一致性；三服務融合方法（本節）提供單一場地之高精度樓高驗證，作為第~\ref{chapter:results}章案例研究與 NSGA-II 最佳化之現況基準（Baseline）；既有隨機生成場景則保留作為第~\ref{sec:v4_validation}節之對照資料集，用以檢驗真實地理環境訓練相對於程序化幾何訓練之外部效度差異。第~\ref{chapter:results}章之案例研究以三服務融合方法擷取之真實建物足跡作為場地現況輸入，評估模型於未見過真實幾何配置下之一般化表現。


\section{模組四：異質圖建構與 PI-ST-GNN 推理——異質圖建構流程}
\label{sec:graph_construction}

\subsection{節點類型與特徵設計}
\label{subsec:nodes}

都市微氣候系統中，空間元素的物理特性差異顯著——建築物為不透氣剛體，主導陰影投射與風場阻滯；空氣評估點為連續流體場採樣，承載氣象時序資訊。本研究因此採用\textbf{異質圖}（Heterogeneous Graph）結構，將兩類節點分別賦予不同語義維度的特徵向量：

\noindent\textbf{建物節點（Object Nodes，$\mathcal{V}_{\text{obj}}$）}

建物節點共 $N_{\text{obj}}$ 個，對應場景內所有建築量體，每個節點攜帶 7 維靜態幾何特徵向量（見表~\ref{tab:obj_features}）；喬木則不構成獨立之建物節點，其高度資訊改以空氣節點第 7 維特徵（樹木高度正規化，見表~\ref{tab:air_features}）之形式，隨鄰近空氣評估點一併編碼。

\begin{table}[htbp]
    \centering
    \caption{建物節點特徵定義（7維靜態特徵）}
    \label{tab:obj_features}
    \small
    \begin{tabular}{lllc}
        \toprule
        \textbf{特徵名稱} & \textbf{物理意義} & \textbf{正規化方式} & \textbf{範圍} \\
        \midrule
        \texttt{height / 50}        & 建築簷高（m）                  & 除以 50m 上限     & $[0, 1]$ \\
        \texttt{floors / 12}        & 樓層數                        & 除以 12 層上限    & $[0, 1]$ \\
        \texttt{footprint / 2000}   & 底層投影面積（m$^2$）           & 除以 2000m$^2$    & $[0, 1]$ \\
        \texttt{cx / 80}            & 質心 X 座標（m）               & 除以場地寬 80m    & $[0, 1]$ \\
        \texttt{cy / 80}            & 質心 Y 座標（m）               & 除以場地高 80m    & $[0, 1]$ \\
        \texttt{gfa / 20000}        & 總樓地板面積（m$^2$）           & 除以 20000m$^2$   & $[0, 1]$ \\
        \texttt{is\_L\_shape}       & 平面形態旗標（L形=1，矩形=0）  & 二值              & $\{0, 1\}$ \\
        \bottomrule
    \end{tabular}
\end{table}

\noindent\textbf{空氣評估節點（Air Nodes，$\mathcal{V}_{\text{air}}$）}

空氣評估節點共 $N_{\text{air}} = 6{,}241$ 個，佈設於場地 $79 \times 79$ 個 1.0m 間距格點上，每個節點在每個時步攜帶 9 維動態特徵向量（見表~\ref{tab:air_features}）：

\begin{table}[htbp]
    \centering
    \caption{空氣評估節點特徵定義（每時步 9 維動態特徵）}
    \label{tab:air_features}
    \small
    \begin{tabular}{lllc}
        \toprule
        \textbf{特徵名稱} & \textbf{物理意義} & \textbf{正規化方式} & \textbf{範圍} \\
        \midrule
        \texttt{Ta\_norm}   & 空氣溫度（°C）             & z-score 標準化  & $\sim\mathcal{N}(0,1)$ \\
        \texttt{MRT\_norm}  & 平均輻射溫度（°C）         & z-score 標準化  & $\sim\mathcal{N}(0,1)$ \\
        \texttt{Va\_norm}   & 水平風速（m/s）            & z-score 標準化  & $\sim\mathcal{N}(0,1)$ \\
        \texttt{RH\_norm}   & 相對濕度（\%）             & z-score 標準化  & $\sim\mathcal{N}(0,1)$ \\
        \texttt{SVF}        & 天空視角因子              & 原始值 $[0,1]$ & $[0, 1]$ \\
        \texttt{shadow}     & 陰影遮蔽旗標              & 二值            & $\{0, 1\}$ \\
        \texttt{Bh / 50}    & 最近建築高度影響（m）      & 除以 50m 上限   & $[0, 1]$ \\
        \texttt{Th / 12}    & 最近樹木高度影響（m）      & 除以 12m 上限   & $[0, 1]$ \\
        \texttt{Ts\_norm}   & 地表溫度估算值（°C）       & z-score 標準化  & $\sim\mathcal{N}(0,1)$ \\
        \bottomrule
    \end{tabular}
\end{table}

其中地表溫度 $T_s$ 並非直接量測值，而由下式之對流熱平衡估算計算得出：

\begin{equation}
    T_s = T_a + \frac{\text{GHI}}{h_c}
    \label{eq:surface_temp}
\end{equation}

式中 $\text{GHI}$（W/m$^2$）為太陽總輻照度（來自 EPW 逐時資料），$h_c$（W/m$^2$K）為對流換熱係數（本研究設定為 10 W/m$^2$K）。此估算方法在忽略長波輻射再輻射與儲熱效應的一階近似下，能有效代表都市鋪面之熱輻射貢獻 \citep{oke1988street}。

表~\ref{tab:obj_features}、\ref{tab:air_features} 所定義之特徵向量彼此並非孤立的數字欄位，而是共同描述同一個都市街道剖面中的物理幾何關係：太陽仰角決定建物之陰影投射範圍（對應 \texttt{shadow} 旗標）與喬木冠層之樹蔭覆蓋（對應 \texttt{Th} 最近樹木高度影響），兩者共同決定行人所在位置之熱壓力高低。圖~\ref{fig:street_crosssection} 以側視剖面示意此一關係：同一街谷寬度 $W$、建物高度 $H$ 之幾何條件下，曝曬於陰影與樹蔭範圍外之行人位置可能承受 UTCI $> 38^\circ\text{C}$ 之強烈熱壓力，而位於樹蔭遮蔽下之行人位置則可能低於 $32^\circ\text{C}$；空氣節點沿垂直方向多高度取樣（圖右側藍點），使異質圖得以捕捉此一熱壓力隨高度與遮蔽狀態變化之空間梯度，而非僅以單一地面高度之純量值代表整條街道剖面之熱環境。

\begin{figure}[htbp]
    \centering
    \includegraphics[width=0.92\textwidth]{../urban-thermal-gnn/04_training/figures/fig_street_crosssection.png}
    \caption[都市街道剖面示意：太陽角度、建築陰影、喬木樹蔭與行人熱壓力]{都市街道剖面示意：太陽角度、建築陰影、喬木樹蔭與行人熱壓力之幾何關係。原創繪製，示意性呈現表~\ref{tab:obj_features}、\ref{tab:air_features}所定義特徵向量之物理對應關係，非特定場景之精確幾何。}
    \label{fig:street_crosssection}
\end{figure}


\subsection{邊緣類型與因果拓樸}
\label{subsec:edges}

本研究定義五種關係型邊緣（Relational Edges），各自對應一種明確的物理因果機制（見表~\ref{tab:edge_types}）。與建物節點不同，物件節點（Object Nodes）實際上同時涵蓋建築與喬木兩類實體（第~\ref{subsec:nodes}節），\texttt{shadow} 與 \texttt{veg\_et} 兩者之來源端因此分別限定為建築類與喬木類物件節點，而非泛稱之「建物」：

\begin{table}[htbp]
    \centering
    \caption{異質圖五類關係型邊緣定義}
    \label{tab:edge_types}
    \small
    \begin{tabularx}{\textwidth}{lllX}
        \toprule
        \textbf{邊緣類型} & \textbf{連接方向} & \textbf{建立方法} & \textbf{物理因果意義} \\
        \midrule
        \texttt{shadow}      & 建築物件 $\to$ 空氣 & 幾何陰影投射計算 & 建築遮擋導致輻射熱降低 \\
        \texttt{veg\_et}     & 喬木物件 $\to$ 空氣 & 樹冠覆蓋半徑判斷 & 植被蒸散作用形成冷卻效應 \\
        \texttt{convective}  & 空氣 $\to$ 空氣 & 對流方向性權重   & 風場驅動的對流熱交換 \\
        \texttt{semantic}    & 物件 $\leftrightarrow$ 物件 & 全連接（完整關係） & 建物與喬木之語義形態相互關係 \\
        \texttt{contiguity}  & 空氣 $\leftrightarrow$ 空氣 & KNN（$k = 8$）  & 空間毗鄰性熱梯度傳導 \\
        \bottomrule
    \end{tabularx}
\end{table}

其中 \texttt{semantic} 與 \texttt{contiguity} 為靜態邊，於資料載入階段一次建立並全程共用；\texttt{shadow}、\texttt{veg\_et}、\texttt{convective} 三者則隨太陽方位角、風向逐時步變化，於每個時步重新計算，構成動態邊集合 \texttt{dynamic\_edges[t]}——此三類動態邊之來源索引（哪一個建築造成該空氣節點陰影、哪一棵喬木覆蓋該空氣節點）由幾何後處理腳本重建取得：程序化主線資料集使用 \texttt{11\_build\_dynamic\_edge\_cache.py}，真實場景資料集（V5）則使用 \texttt{16\_build\_dynamic\_edge\_cache\_v5.py}（採真實逐場景太陽方位角與真實逐月中央氣象署風向），完整修正過程記錄於附錄~\ref{appx:dynamic_edges_fix}。圖~\ref{fig:graph_construction} 以真實資料集之場景 0（$N_{\text{obj}}=3$，$N_{\text{air}}=6{,}241$）具體呈現此一異質圖之節點分布、真實建立之 KNN 毗鄰邊局部結構，以及完整節點／邊型別綱要。

\begin{figure}[htbp]
    \centering
    \includegraphics[width=\textwidth]{../urban-thermal-gnn/04_training/figures/fig_graph_construction.png}
    \caption[異質圖建構實例]{異質圖建構實例（取自真實訓練資料場景 0）}
    \label{fig:graph_construction}
\end{figure}

上圖 (a) 為建物節點（紅色量體）與空氣評估節點（依時均 UTCI 著色）之空間分布；(b) 為局部 $18\times18$ 公尺窗口內 KNN($k$=8) 毗鄰邊之真實連接結構，以及該範圍內建物節點向場外其他建物延伸之全連接語義邊；(c) 為異質圖節點／邊型別綱要，對應 9 維空氣節點特徵（含地表溫度）與五類關係型邊緣定義。

圖~\ref{fig:graph_topology} 進一步以另一真實場景（場景 10）之全域視角，呈現 \texttt{semantic} 邊如何在場域內所有建物節點間建立全連接拓樸——不同於圖~\ref{fig:graph_construction}(b) 僅呈現單一建物之局部窗口，此處可清楚看見 $N_{\text{obj}}=4$ 個建物節點兩兩相連之完整語義邊網絡，與其下方 $6{,}241$ 個空氣節點所構成之密集 KNN 毗鄰網格形成尺度對比。圖中另以綠色圓標示場景內之喬木冠層位置作為空間對照，惟依第~\ref{subsec:nodes}節之節點定義，喬木高度僅作為空氣節點特徵（樹木高度正規化）之來源，本身不構成獨立圖節點，故未參與語義邊之連接。

\begin{figure}[htbp]
    \centering
    \includegraphics[width=\textwidth]{../urban-thermal-gnn/04_training/figures/fig_geo5_graph_topology.png}
    \caption[異質圖拓樸全域視角]{異質圖拓樸全域視角（真實訓練資料場景 10）}
    \label{fig:graph_topology}
\end{figure}

上圖 (a) 為完整 $80\times80$ 公尺場域，橘點為 6,241 個空氣節點（KNN 毗鄰邊為求可讀性僅隨機顯示 3,000／49,928 條），橘色量體為建物節點、綠色圓為喬木冠層之空間位置（僅供視覺對照，非圖節點），紫色虛線為建物節點間之全連接語義邊；(b) 為場域左下 $40\times40$ 公尺象限放大，標示邊緣類型與 9 維空氣節點特徵之完整定義。

作為對照，圖~\ref{fig:real_scene_graph} 以第~\ref{subsec:v4_real_dataset}節真實場景資料集中之一個實際場景，展示完全以真實 GIS 圖資建構之異質圖：(a) 之建物量體均為 OSM 真實足跡（色階為真實簷高，實線外框標示具真實樓高標籤者），綠點為 ETH 樹冠實測植栽；(b) 之每個空氣節點依 4 公尺格點佈設，節點顏色為以該場景所在真實 IoT 測站逐時氣溫驅動之模擬尖峰 UTCI，可見建物遮蔽與植栽降溫在空間上形成之熱舒適梯度，即為 PI-ST-GNN 所學習之真實空間結構。

\begin{figure}[htbp]
    \centering
    \includegraphics[width=\textwidth]{../urban-thermal-gnn/04_training/figures/fig_real_scene_graph.png}
    \caption[真實 GIS 圖資建構之異質圖範例]{以真實 GIS 圖資建構之異質圖範例（真實場景資料集）：(a) 真實 OSM 建物與 ETH 樹冠幾何；(b) 空氣節點感測網格與 k-NN 鄰接邊，節點色為真實 IoT 氣溫驅動之模擬尖峰 UTCI。}
    \label{fig:real_scene_graph}
\end{figure}


\subsection{圖索引對齊機制（$+N_{\text{obj}}$ 偏移量）}
\label{subsec:index_offset}

異質圖在送入 RGCN 之前，本研究採用\textbf{節點串接}（Node Concatenation）策略，將建物節點特徵向量 $\mathbf{H}_{\text{obj}} \in \mathbb{R}^{N_{\text{obj}} \times d_{\text{obj}}}$ 與空氣節點特徵 $\mathbf{H}_{\text{air}} \in \mathbb{R}^{N_{\text{air}} \times d_{\text{air}}}$ 沿節點維度串接，形成單一混合節點張量：

\begin{equation}
    \mathbf{H} = \left[ \mathbf{H}_{\text{obj}} ; \mathbf{H}_{\text{air}} \right] \in \mathbb{R}^{(N_{\text{obj}} + N_{\text{air}}) \times d}
    \label{eq:node_concat}
\end{equation}

此機制意味著空氣節點在串接後的全域索引空間中，其索引應整體偏移 $N_{\text{obj}}$。因此，KNN 毗鄰邊（\texttt{contiguity}）在建立時，所有空氣節點的 \texttt{edge\_index} 均需加上偏移量 $N_{\text{obj}}$：

\begin{equation}
    \texttt{edge\_index}_{\text{contiguity}} \mathrel{+}= N_{\text{obj}}
    \label{eq:offset}
\end{equation}

此對齊機制確保跨節點類型之邊緣（如 \texttt{shadow}: 建物 $\to$ 空氣）在統一索引空間中仍能正確尋址目標節點，是整個異質圖能以單一 RGCN 前向傳播處理的關鍵實作細節，完整偏移量計算流程另見附錄~\ref{appx:index_offset_detail}。


\section{物理資訊時空圖神經網路 (PI-ST-GNN) 模型架構}
\label{sec:gnn_model}

\subsection{整體架構概述}
\label{subsec:model_overview}

PI-ST-GNN 模型由以下五個主要子模組組成，依前向傳播順序排列（詳見圖~\ref{fig:arch_pipeline}）：

\begin{enumerate}
    \item \textbf{物件節點編碼器（Object Encoder）：}以兩層 MLP 將 7 維靜態幾何特徵投影至 128 維隱藏表示 $\mathbf{h}_{\text{obj}} \in \mathbb{R}^{N_{\text{obj}} \times 128}$。
    \item \textbf{空氣節點編碼器（Air Encoder）：}以兩層 MLP 將每時步 9 維動態氣象特徵投影至 128 維隱藏表示 $\mathbf{h}_{\text{air},t} \in \mathbb{R}^{N_{\text{air}} \times 128}$，對 $T = 11$ 個時步分別計算。
    \item \textbf{全域情境編碼器（Global Context Encoder）：}將場景整體氣象背景（8 維逐時統計特徵：全場域平均 $T_a, V_a, \text{RH}, \text{GHI}$ 及其差分）以 MLP 編碼為時序情境向量 $\mathbf{c} \in \mathbb{R}^{T \times 192}$，與每時步之空氣節點表示融合，注入全域氣候感知能力。
    \item \textbf{關係圖卷積網路（RGCN）：}在每個時步 $t$，將建物與空氣節點之串接表示送入三層堆疊的 RGCN 塊（RGCNBlock），在五種關係型邊緣上進行異質訊息聚合，輸出包含空間因果資訊的混合節點表示。
    \item \textbf{LSTM 時序建模與輸出頭：}以 Fusion MLP 融合 RGCN 輸出與全域情境向量後，送入單層 LSTM，其初始隱藏狀態 $\mathbf{h}_0$ 由 $t = 0$ 時刻的 RGCN 輸出投影得出，最終由輸出 MLP 預測 11 時步之 UTCI 值 $\hat{\mathbf{u}} \in \mathbb{R}^{N_{\text{air}} \times T}$。
\end{enumerate}

\begin{figure}[htbp]
    \centering
    \includegraphics[width=0.92\textwidth]{../urban-thermal-gnn/04_training/figures/fig_arch_clean.png}
    \caption[整合運算框架示意圖]{PI-ST-GNN 整合運算框架示意圖}
    \label{fig:arch_pipeline}
\end{figure}

上圖左側輸入包含建物幾何特徵（7 維）與多時步空氣節點氣象特徵（9 維 $\times$ 11 步），中部為 RGCN 空間建模與 LSTM 時序融合，右側輸出全場域 UTCI 熱壓力場（11 時段），並以物理資訊損失函數進行約束訓練。

上述五個子模組彼此之銜接關係，以及第~\ref{sec:graph_construction}節之異質圖建構如何作為此模型之前置輸入，完整呈現於圖~\ref{fig:overview_pistgnn}。此圖並將第~\ref{subsec:lstm}節所述「LSTM 初始隱藏狀態由 RGCN 輸出投影得出、而非零向量初始化」之時序展開細節，以下方獨立面板逐時步具體標示——每個時步 $t$ 之異質圖先經 RGCN 完成當時步之空間訊息擴散，其輸出再與前一時步之 LSTM 隱藏狀態 $h_{t-1}$ 共同送入 LSTM 單元更新 $h_t$，如此逐時步展開直至 $t=10$（對應 18:00），最終輸出張量 $(N_{\text{air}}, T{=}11, 256)$ 送入輸出 MLP 完成多時步解碼。此一「空間擴散在前、時序遞迴在後」之逐時步交替結構，是 PI-ST-GNN 得以同時捕捉建物陰影瞬時空間效應與都市熱慣性累積效應兩種不同時間尺度物理機制的架構關鍵。

\begin{figure}[htbp]
    \centering
    \includegraphics[width=\textwidth]{../urban-thermal-gnn/04_training/figures/fig_overview_pistgnn_white.png}
    \caption[PI-ST-GNN 模型架構總覽]{PI-ST-GNN 模型架構總覽}
    \label{fig:overview_pistgnn}
\end{figure}

上圖呈現異質圖建構、編碼、解碼三階段，及 RGCN 擴散＋LSTM 遞迴之逐時步展開細節。


\subsection{關係圖卷積網路（RGCN）}
\label{subsec:rgcn}

本研究的核心空間建模模組為\textbf{關係圖卷積網路}（Relational Graph Convolutional Network, RGCN） \citep{wu2023data, shan2025physics}。不同於同質圖卷積（GCN）以單一共享權重聚合所有鄰域，RGCN 為每種關係類型 $r \in \mathcal{R}$ 維護獨立的可學習投影矩陣 $\mathbf{W}_r$，得以在訊息傳遞過程中區分不同物理因果關係之訊號貢獻。本研究之關係型聚合公式與殘差＋LayerNorm 更新規則，採用並延伸 \citet{xin2025urbangraph} 於 UrbanGraph 架構中提出之關係圖卷積形式，於此基礎上進一步將節點集合區分為建物與空氣兩種異質節點類型（見第~\ref{subsec:nodes}節），並以獨立編碼器分別處理其靜態／動態特徵。

對於節點 $i$，第 $l$ 層 RGCN 塊（RGCNBlock）的更新規則定義為：

\begin{equation}
    \mathbf{h}_i^{(l+1)} = \sigma\!\left( \mathbf{W}_{\text{self}} \mathbf{h}_i^{(l)}
    + \sum_{r \in \mathcal{R}} \frac{1}{\left|\mathcal{N}_r(i)\right|}
    \sum_{j \in \mathcal{N}_r(i)} \mathbf{W}_r \mathbf{h}_j^{(l)} \right) + \mathbf{h}_i^{(l)}
    \label{eq:rgcn}
\end{equation}

其中 $\mathbf{W}_{\text{self}}$ 為自迴路權重矩陣、$\mathbf{W}_r$ 為關係 $r$ 專屬投影矩陣（本研究 $|\mathcal{R}| = 5$）、$\mathcal{N}_r(i)$ 為節點 $i$ 在關係 $r$ 下的鄰居集合、$\sigma(\cdot)$ 為 PReLU 激活函數，末項為殘差連接並於激活後接 LayerNorm 與 Dropout（$p=0.1$）。模型採用三層堆疊（$L=3$）的 RGCNBlock，允許訊息在三跳鄰域內擴散，有效涵蓋 $80\times80$ 公尺場地內建築陰影的最大投射距離；各關係矩陣之完整維度設定見附錄~\ref{appx:rgcn_detail}。此一訊息聚合機制之圖示見圖~\ref{fig:rgcn_detail}。

\begin{figure}[htbp]
    \centering
    \includegraphics[width=0.85\textwidth]{../urban-thermal-gnn/04_training/figures/fig_rgcn_detail.png}
    \caption[RGCN 訊息聚合機制示意圖]{RGCN 模組訊息聚合機制示意圖}
    \label{fig:rgcn_detail}
\end{figure}

上圖中央目標節點 $i$ 接收來自五類關係型鄰居節點的加權訊息，每類關係使用獨立投影矩陣 $\mathbf{W}_r$ 轉換後進行均值聚合，最終與自迴路項疊加並通過殘差＋LayerNorm 輸出；三層堆疊使訊息傳遞最遠可達三跳距離之鄰域。


\subsection{全域情境融合與 LSTM 時序建模}
\label{subsec:lstm}

場景整體氣候狀態（如晴天 GHI 強度、背景風速）是影響全場域 UTCI 分佈之關鍵驅動因素，但此資訊難以從局部節點特徵中推導，因此需要一個獨立於局部訊息傳遞之外的全域情境通道。全域情境編碼器對每個時步 $t$ 提取 8 維全場域統計特徵，以三層 MLP 映射至 192 維情境向量：

\begin{equation}
    \mathbf{c}_t = \text{MLP}_{\text{ctx}}\!\left( \left[\mu_{T_a}^t,\ \mu_{V_a}^t,\ \mu_{\text{RH}}^t,\ \text{GHI}^t,\ \Delta\mu_{T_a}^t,\ \Delta\mu_{V_a}^t,\ \Delta\text{GHI}^t,\ \text{sol\_alt}^t \right] \right)
    \label{eq:ctx}
\end{equation}

對每個空氣節點，Fusion MLP 融合 RGCN 輸出之節點特徵（256維）與全域情境向量（192維），投影至 256 維融合表示 $\mathbf{f}_{i,t}$，再送入共享 LSTM 捕捉 UTCI 之日間熱積累動態：

\begin{equation}
    \mathbf{o}_{i,t},\ \mathbf{h}_{i,t}^{\text{LSTM}} = \text{LSTM}\!\left( \mathbf{f}_{i,t},\ \mathbf{h}_{i,t-1}^{\text{LSTM}} \right)
    \label{eq:lstm}
\end{equation}

LSTM 之初始隱藏狀態由 $t=0$ 時刻的 RGCN 輸出投影得出而非零向量，使模型能以場景幾何資訊初始化時序動態的起始狀態，避免冷啟動問題；此一暖啟動機制沿用 \citet{xin2025urbangraph} 之設計。最終，輸出 MLP 將 LSTM 輸出映射至 UTCI 預測值 $\hat{u}_{i,t}$（z-score 正規化空間中預測，推理後去正規化還原至實際溫度單位 °C）。完整超參數設定見附錄~\ref{appx:model_hparams}。


\section{物理引導損失函數}
\label{sec:physics_loss}

純資料驅動之損失函數僅要求預測值逼近模擬標籤，並不保證預測結果服從基本熱力學原則——模型仍可能在資料稀疏區域輸出陰影節點溫度高於日照節點的物理不可能結果。本研究因此於資料損失之外，引入三項\textbf{物理資訊損失函數}（Physics-Informed Loss）作為訓練過程中的軟約束監督機制 \citep{shan2025physics}，總損失函數定義為：

\begin{equation}
    \mathcal{L}_{\text{total}} = \mathcal{L}_{\text{data}} + \alpha_1 \mathcal{L}_{\text{rad}} + \alpha_2 \mathcal{L}_{\text{temp}} + \alpha_3 \mathcal{L}_{\text{wind}}
    \label{eq:total_loss}
\end{equation}

其中 $\mathcal{L}_{\text{data}}$ 為資料驅動項（與物理式模擬標籤之 MSE 損失），$\alpha_1, \alpha_2, \alpha_3$ 為物理損失權重係數。三項物理約束各自對應之空間、時間、風場作用機制與懲罰觸發條件，統整於圖~\ref{fig:physics_loss}；各懲罰項之完整符號定義與數值設定說明如下三小節，其嚴謹解析式推導則置於附錄~\ref{appx:physics_loss_derivation}。

\begin{figure}[htbp]
    \centering
    \includegraphics[width=0.9\textwidth]{../urban-thermal-gnn/04_training/figures/fig_physics_loss.png}
    \caption[物理引導損失函數示意圖]{物理引導損失函數示意圖}
    \label{fig:physics_loss}
\end{figure}

上圖三項約束分別作用於不同物理維度：輻射空間一致性約束（左）比較日照／陰影節點群組之預測值相對關係；都市熱慣性時序平滑約束（中）限制相鄰時步 UTCI 之最大變化率；風場阻滯尾流約束（右）抑制高層建築背風側之異常熱累積。灰色區域代表各約束實際產生懲罰梯度之觸發範圍，其餘區域損失為零，僅在物理不可能之預測情形下介入訓練。


\subsection{輻射空間一致性約束（$\mathcal{L}_{\text{rad}}$）}
\label{subsec:radiation_penalty}

建築陰影投射是影響行人 UTCI 最關鍵的幾何因子之一：日照節點之熱壓力應顯著高於陰影節點，而非兩者相近或反轉 \citep{coccolo2018thermal}。當太陽高度角 $\theta_{\text{sol}} > 10^\circ$ 時，依 SVF 與即時陰影旗標將節點分類為日照節點（$\text{SVF}_i > 0.5$ 且未受遮蔽）與陰影節點（$\text{SVF}_i < 0.5$ 且受遮蔽），懲罰項要求場域陰影節點均值與日照節點均值之間至少維持最小間隔 $m_{\text{rad}}$，對未達此間隔（含反轉）之部分施以平方懲罰：

\begin{equation}
    \mathcal{L}_{\text{rad}} =
    \left[ \text{ReLU}\!\left( \bar{u}_{\text{shade}} - \bar{u}_{\text{sun}} + m_{\text{rad}} \right) \right]^2
    \label{eq:lrad}
\end{equation}

其中 $\bar{u}_{\text{shade}}$、$\bar{u}_{\text{sun}}$ 分別為當前時步場域陰影節點與日照節點之預測均值，最小間隔 $m_{\text{rad}} = 0.5$（對應正規化 UTCI 空間中約半個標準差）；式~\eqref{eq:lrad}於每個滿足 $\theta_{\text{sol}}>10^\circ$ 之時步各計算一次，訓練時取其對有效時步之平均。此設計要求日照節點至少比陰影節點高出 $m_{\text{rad}}$，而非僅禁止陰影反超日照——理由是完全日照與完全遮蔭之間應存在物理上有意義的溫差，僅約束「不反轉」仍可能容許模型學習出兩者幾乎相等的退化解。


\subsection{都市熱慣性時序平滑約束（$\mathcal{L}_{\text{temp}}$）}
\label{subsec:temporal_penalty}

都市物質（混凝土、瀝青鋪面）具有顯著的熱儲存效應，使 UTCI 在相鄰時步間的變化率受到自然物理限制——熱壓力不應在一小時內劇烈跳躍。本研究將允許之最大 UTCI 變化率設定為 $5^\circ\text{C/hr}$（正規化空間中 $\delta_{\text{temp}} = 0.625$，其換算依據見附錄~\ref{appx:physics_loss_derivation}）：

\begin{equation}
    \mathcal{L}_{\text{temp}} = \text{mean}_{i,t}\!\left[ \text{ReLU}\!\left( \left| \hat{u}_{i,t+1} - \hat{u}_{i,t} \right| - \delta_{\text{temp}} \right) \right]
    \label{eq:ltemp}
\end{equation}

式中對超出容許變化率之部分採\textbf{線性}鉸鏈懲罰（而非平方），使梯度大小不隨違反程度二次放大，讓訓練初期時序尚不穩定時，此一約束不至於因梯度過大而主導總損失。此懲罰項能有效防止模型在學習日出／日落高輻射梯度時產生非物理的時間振盪，提升整體逐時預測軌跡的熱力學自洽性 \citep{oskarsson2023modeling}。


\subsection{風場阻滯尾流約束（$\mathcal{L}_{\text{wind}}$）}
\label{subsec:wind_penalty}

高層建築（$B_h > 20\text{m}$，即正規化特徵 $B_h/50 > 0.4$）在背風側產生顯著的渦旋死角（Leeward Wake），使局部熱能累積形成「熱死角」現象 \citep{oke1987boundary}。對於位於高層建築鄰近範圍之空氣節點集合 $\mathcal{B}_{\text{tall}}$，風場阻滯懲罰定義為：

\begin{equation}
    \mathcal{L}_{\text{wind}} = \text{mean}_{i \in \mathcal{B}_{\text{tall}}}
    \left[ \text{ReLU}\!\left( \hat{u}_i - \mu_{\hat{u}} - 1.5\,\sigma_{\hat{u}} \right) \right]
    \label{eq:lwind}
\end{equation}

其中 $\mu_{\hat{u}}$、$\sigma_{\hat{u}}$ 分別為當前時步全場域預測 UTCI（所有節點，非僅 $\mathcal{B}_{\text{tall}}$）之均值與標準差，懲罰同樣採線性鉸鏈形式。閾值 $\mu_{\hat{u}} + 1.5\sigma_{\hat{u}}$ 之選定依循統計離群偵測之常見慣例——在常態分布假設下，超出均值 $1.5$ 個標準差約對應最熱的 6.7\% 尾端節點，將「熱死角」操作型定義為場域中溫度顯著偏離常態分布中心之高層建築鄰近節點。此約束強制模型學習高層建築背風側之熱累積模式，使設計師能可視化熱死角位置，進而最佳化建築間距與通風廊道配置。


\section{模組五：Grasshopper 介面整合與多目標最佳化——熱舒適人行動線設計}
\label{sec:nsga2}

\subsection{從最短路徑到熱舒適動線評估}
\label{subsec:from_astar_to_walkway}

早期構想以 $A^*$ 最短路徑演算法將 UTCI 熱壓力場轉譯為單一起訖點間之熱迴避路徑：以節點間歐氏距離與熱壓力懲罰之加權組合作為邊成本，求解一條在行走效率與熱舒適間取捨之路徑。然而，都市開放空間設計決策的對象通常並非單一使用者的單次行程，而是整體人行動線網絡在一整日各時段的熱舒適表現；單點對單點的最短路徑求解難以回答「這個街廓的人行動線系統整體而言是否宜行」此一設計層級問題。本研究因此將原先的路徑搜尋機制，重新定位為\textbf{熱舒適人行動線設計}（Thermal Comfort-Based Walkway Design）評估指標，直接整合進 NSGA-II 多目標最佳化之適應度函數，而非作為獨立之單次路徑查詢工具。

具體而言，設計者於 Grasshopper 中定義之人行動線網絡幾何（步行軸線、廣場動線、建物出入口連通路徑）作為固定拓樸輸入，GNN 代理模型即時預測之全場域 UTCI 熱壓力場則沿此網絡逐段取樣，計算全動線之\textbf{時空平均熱暴露指標}：

\begin{equation}
    \bar{\Phi}_{\text{walkway}} = \frac{1}{T}\sum_{t=1}^{T} \frac{1}{|E_{\text{walk}}|}
    \sum_{(i,j) \in E_{\text{walk}}} D_{ij}\cdot \Phi\!\left(\text{UTCI}_{j,t}\right)
    \label{eq:walkway_exposure}
\end{equation}

其中 $E_{\text{walk}}$ 為人行動線網絡之邊集合、$D_{ij}$ 為邊之實際步行距離、$\Phi(\cdot)$ 為正規化熱壓力懲罰函數，其定義與完整推導見附錄~\ref{appx:walkway_derivation}。此一時空平均指標取代原先僅評估單一路徑之作法，使評估對象與 NSGA-II 之全域最佳化目標在尺度上一致。


\subsection{Grasshopper 介面與非同步通訊架構}
\label{subsec:gh_integration}

系統採用\textbf{非同步用戶端-伺服器架構}，讓建築師能在熟悉的設計環境中獲得 AI 的即時反饋：在 Grasshopper 中開發之專用套件（\texttt{UTCIPredictor}、\texttt{UTCIOptimizer}）負責將 Rhino 3D 幾何物件序列化為異質圖相容張量；後端透過 IETF RFC 6455 \textbf{WebSocket} 通訊協定 \citep{rfc6455} 與 FastAPI 通訊，載入預先訓練完成的 PI-ST-GNN 模型權重回傳預測結果與最佳化進度；設計師於 Rhino 中移動建築時，系統更新場地 1.0m 間距之高解析度熱舒適度網格，使原本須另行呼叫模擬引擎之熱舒適評估得以整合入設計操作流程 \citep{sadeghipour2013ladybug}。

圖~\ref{fig:gh_websocket} 以時序圖具體呈現此一通訊架構之兩種訊息流：(a) 單次幾何更新之同步請求-回應循環，直接對應每次推論之低延遲特性；(b) NSGA-II 最佳化工作之非阻塞式執行——最佳化迴圈於獨立執行緒（\texttt{ThreadPoolExecutor}）中運行，避免長時間運算阻塞 WebSocket 事件迴圈，並透過 \texttt{asyncio.Queue} 將逐代進度（世代數、當前最佳解、可行解數量）以非同步訊息即時串流回 Grasshopper 前端，使建築師無需等待完整最佳化流程結束即可觀察收斂趨勢。

\begin{figure}[htbp]
    \centering
    \includegraphics[width=\textwidth]{../urban-thermal-gnn/04_training/figures/fig_gh_websocket_sequence.png}
    \caption[非同步通訊架構時序圖]{Grasshopper--FastAPI／WebSocket 非同步通訊架構時序圖}
    \label{fig:gh_websocket}
\end{figure}

上圖實線為同步 WebSocket 訊息，虛線為經 \texttt{asyncio.Queue} 橋接執行緒之非同步進度回呼。


\subsection{染色體編碼機制與多目標函數}
\label{subsec:design_variables}

最佳化引擎將都市配置參數化為連續染色體基因組，所有變數於運算前歸一至 $[0,1]$ 區間：每棟建築由質心座標 $(c_x, c_y)$、面寬 $w$、進深 $d$、平面旋轉角度與樓層數共六項參數定義；每棵喬木由座標 $(x,y)$、樹冠半徑與高度四項參數組成，直接決定遮蔭效果與綠化率貢獻。

系統透過 NSGA-II \citep{deb2002fast} 同時處理三個相互競爭的適應度指標：

\begin{enumerate}
    \item \textbf{$f_1$（全域 UTCI 極小化）：}最小化全場地所有開放空間空氣節點在 11 時段內的時空平均熱壓力。
    \item \textbf{$f_2$（人行動線熱暴露極小化）：}最小化式~\eqref{eq:walkway_exposure}定義之人行動線時空平均熱暴露指標 $\bar{\Phi}_{\text{walkway}}$。
    \item \textbf{$f_3$（綠化率極大化）：}最大化場地綠化率，兼顧生態價值。
\end{enumerate}

法規合規性透過 5 維違規向量（FAR、BCR、法定縮退、場地容納性、建築碰撞）加以監控，採「可行性優先」選擇啟發式：符合規範之可行解絕對優先於違規解排序。最終輸出之帕累托前緣代表一組最優設計集，讓建築師在整體熱舒適性能、人行動線品質與綠化率之間進行數據驅動的決策。

本研究之最佳化引擎（\texttt{07\_optimization/nsga2\_engine.py}）已完整實作 $f_1$（全域 UTCI）與 $f_3$（綠化率）兩項目標之適應度函數，並以此二維目標空間驗證第~\ref{sec:pareto_analysis}節之收斂表現；$f_2$（人行動線熱暴露）之數學定義與推導已於式~\eqref{eq:walkway_exposure}及附錄~\ref{appx:walkway_derivation}完整給出，其整合入適應度評估器（\texttt{FitnessEvaluator}）之工程實作——即將染色體所定義之人行動線網絡幾何與 GNN 預測場逐段取樣、回傳第三項目標值——列為第~\ref{sec:future_work}節之近期後續工作，目前尚未併入單次最佳化執行流程，特此說明，避免與第~\ref{sec:pareto_analysis}節實際呈現之收斂結果所涵蓋目標數產生混淆。


\subsection{與 Wallacei 都市型態最佳化案例之比較}
\label{subsec:wallacei_comparison}

Grasshopper 生態系中，Wallacei 是目前應用最廣泛之進階 NSGA-II 視覺化最佳化外掛，已見於多項都市型態多目標最佳化實務研究 \citep{wortmann2017opossum, daglier2025developing}。既有以 Wallacei 為核心之都市形態最佳化案例，其適應度函數之環境績效評估多直接呼叫 Ladybug Tools／EnergyPlus 進行逐代個體之全氣候模擬，因此在演算法收斂所需的數百代、每代數百個體之規模下，單次最佳化任務之總運算時間常達數週甚至數月，實務上多僅能於方案設計完成後進行事後驗證，而非設計初期之即時決策引導 \citep{aman2023ai, daglier2025developing}。

本研究之最佳化引擎與 Wallacei 案例共用相同的 NSGA-II 核心演算法與 Grasshopper 介面整合模式，惟以第~\ref{sec:gnn_model}節訓練完成之 PI-ST-GNN 取代 Wallacei 案例中直接呼叫之高擬真度模擬引擎作為適應度函數之環境績效來源。此一取代之目的並非單純追求運算速度提升，而在於使多目標最佳化之搜尋迴圈得以在設計對話的時間尺度內完成，讓建築師於方案概念階段即可即時比較不同人行動線配置與量體組合之熱舒適表現，而非僅能於方案定案後進行單次驗證。兩者於相同法規約束與設計變數空間下之收斂表現比較，見第~\ref{chapter:results}章。

\chapter{實驗成果分析與驗證}
\label{chapter:experiments}

第三章建立了 PI-ST-GNN 之完整方法論架構。本章之目的並非展示運算速度的極致最佳化，而是回答一個對建築資訊學實務更關鍵的問題：這套物理資訊圖神經網路，是否真正學到了都市微氣候的空間因果結構，其預測結果是否足以支撐建築師的設計決策。本章依序呈現訓練配置、測試集量化驗證、消融實驗之設計選擇分析，以及運算效能對互動式設計工作流的實質意義。


\section{訓練流程與超參數配置}
\label{sec:training}

\subsection{訓練環境與硬體配置}
\label{subsec:hardware}

本研究所有訓練與推理實驗均於搭載下列硬體之工作站執行：

\begin{itemize}
    \item \textbf{GPU：}NVIDIA GeForce RTX 5070 Ti Laptop GPU（視訊記憶體 12.8 GB）
    \item \textbf{CPU：}AMD Ryzen AI 9 HX（架構：Zen 5）
    \item \textbf{框架：}PyTorch 2.x + 自製輕量圖引擎（無 PyG 依賴）
\end{itemize}


\subsection{訓練超參數與學習率排程}
\label{subsec:training_params}

PI-ST-GNN 之訓練遵循下列配置，詳見表~\ref{tab:training_config}：

\begin{table}[htbp]
    \centering
    \caption{模型訓練超參數配置}
    \label{tab:training_config}
    \small
    \begin{tabular}{llc}
        \toprule
        \textbf{超參數} & \textbf{說明} & \textbf{設定值} \\
        \midrule
        批次大小（Batch Size）         & 每批場景數（逐場景）            & 1 \\
        最大訓練週期（Max Epochs）     & 允許最長訓練回合數              & 250 \\
        早停耐心（Early Stopping）     & 驗證損失無改善之容忍週期數      & 20 \\
        初始學習率（LR）               & Adam 最佳化器初始步長             & $1 \times 10^{-3}$ \\
        學習率排程（LR Scheduler）     & ReduceLROnPlateau，觸發後對半   & factor = 0.5 \\
        學習率下限（Min LR）           & 排程器觸發之最小學習率下限      & $\approx 7.8 \times 10^{-6}$ \\
        資料損失函數（$\mathcal{L}_{\text{data}}$）& 預測值與標籤之 MSE  & --- \\
        物理損失權重（$\alpha_1, \alpha_2, \alpha_3$）& 輻射、時序、風場    & 各自調校 \\
        Dropout                        & RGCN 訓練期間 Dropout 率        & 0.1 \\
        \bottomrule
    \end{tabular}
\end{table}

學習率排程採用 ReduceLROnPlateau 策略：當驗證損失在連續數個 Epoch 未見改善時，學習率縮減為原來的 0.5 倍。本研究最終採用之 V5-300 模型（真實場景資料集，Radiance／EnergyPlus 模擬引擎，見第~\ref{sec:v4_validation}節）訓練中，學習率於第 23、111、130 Epoch 觸發三次衰減（自 $10^{-3}$ 逐步降至 $1.25 \times 10^{-4}$），模型於第 139 Epoch 因連續 20 個 Epoch 無改善觸發早停（Early Stopping）而結束訓練，最佳權重取自第 \textbf{119} Epoch，驗證 $R^2$ 達到 \textbf{0.9837}。

\subsection{訓練曲線與收斂特性}
\label{subsec:convergence}

圖~\ref{fig:training_curves} 展示訓練損失、驗證損失與驗證 $R^2$ 之逐 Epoch 演化曲線，其分段特性揭示模型學習都市微氣候映射關係的三個階段：初期（前約 10 個 Epoch）驗證 $R^2$ 自暖身階段急速上升，顯示模型迅速掌握異質圖節點特徵與熱壓力間的基本映射；中期（Epoch 10--100）$R^2$ 持續波動精調，對應真實建物遮蔽、植栽蒸散、對流與真實地表熱傳導等物理管道之關係權重逐步分化學習；後期（Epoch 100 後）訓練與驗證損失同步趨於平穩收斂，驗證 $R^2$ 穩定於 0.98 附近，未見發散，確認物理資訊損失函數對解空間提供了有效的正則化，模型未發生過擬合。

\begin{figure}[htbp]
    \centering
    \includegraphics[width=\textwidth]{../urban-thermal-gnn/04_training/viz_output/training_v5_300/figA_training_curves.png}
    \caption[模型訓練收斂曲線]{模型訓練收斂曲線（V5-300，共 139 Epoch，早停於最佳 Epoch 119）}
    \label{fig:training_curves}
\end{figure}

上圖上半部為訓練損失（藍）與驗證損失（橘）之逐 Epoch MSE 演化，灰色垂直線標示學習率衰減時間點；下半部為驗證 $R^2$ 逐 Epoch 提升曲線，最終收斂至最佳 $R^2 = 0.9837$（Epoch 119，紅點標示）。


\section{測試集量化評估結果}
\label{sec:evaluation}

\subsection{整體預測精確度指標}
\label{subsec:overall_metrics}

在 45 個真實場景測試集、共計 166,188 個預測點次（各場景空氣節點數依真實建物量體與場地邊界而異，介於 136--400 節點間，逐場景 $\times$ 11 時步加總）上，PI-ST-GNN（V5-300：真實 OSM／ETH／IoT 地理資料 $+$ Radiance／EnergyPlus 模擬引擎，見第~\ref{sec:v4_validation}節）之最終測試集評估結果如表~\ref{tab:main_results} 所示。

\begin{table}[htbp]
    \centering
    \caption[測試集評估結果]{PI-ST-GNN（V5-300）測試集評估結果（共 166,188 預測點次，45 個真實測試場景）}
    \label{tab:main_results}
    \small
    \begin{tabular}{lcccc}
        \toprule
        \textbf{指標} & \textbf{$R^2$} & \textbf{RMSE（°C）} & \textbf{MAE（°C）} & \textbf{等級準確率} \\
        \midrule
        PI-ST-GNN（本研究） & \textbf{0.9753} & \textbf{0.833} & \textbf{0.524} & \textbf{93.9\%} \\
        \bottomrule
    \end{tabular}
\end{table}

四項指標對應四個不同的設計輔助情境，缺一不可：決定係數 $R^2 = 0.9753$ 代表模型能解釋 UTCI 變異量的 97.53\%，超過一般代理模型部署所需之 $R^2 > 0.97$ 門檻 \citep{westermann2019surrogate}；均方根誤差 RMSE $= 0.833\,^\circ\text{C}$（去正規化後之實際溫度單位）遠小於 UTCI 熱壓力等級間距（最小間距 $6\,^\circ\text{C}$），確保等級判斷之邊界不因模型誤差而錯置 \citep{fiala2011utci}；平均絕對誤差 MAE $= 0.524\,^\circ\text{C}$ 低於 RMSE，反映多數評估點具有極高逐點精確度、少數離群值未顯著拉高整體偏差；熱壓力等級分類準確率 93.9\% 則是實際設計輔助情境下最直接相關的指標——建築師關心的通常不是溫度小數點後第二位，而是某一開放空間是否落在「強烈」或「極端」熱壓力等級。此組數字係於 Radiance／EnergyPlus 模擬之真實場景上取得，較純程序化幾何之數值指標略低，其原因與方法論意義詳見第~\ref{sec:v4_validation}節。

\begin{figure}[htbp]
    \centering
    \includegraphics[width=0.85\textwidth]{../urban-thermal-gnn/04_training/viz_output/training_v5_300/figB_scatter_test.png}
    \caption[測試集預測值散點圖]{測試集預測值與模擬標籤之散點圖（V5-300，共 166,188 個樣本）}
    \label{fig:scatter_test}
\end{figure}

上圖點雲緊密分佈於對角線 $y=x$ 附近，色彩代表密度（黃色為高密度區域），圖中所示分布形態與表~\ref{tab:main_results}之 $R^2=0.9753$ 一致，顯示模型在全 UTCI 值域內均具優秀的預測一致性。


\subsection{逐時段預測精確度分析}
\label{subsec:hourly_r2}

圖~\ref{fig:hourly_r2} 展示 11 個時步（08:00--18:00）之個別 $R^2$ 值。所有時步之 $R^2$ 均維持在良好水準，充分驗證模型對日間熱壓力動態的全面捕捉能力；高峰熱壓力時段（10:00--14:00）因太陽高度角最大、陰影梯度最為顯著，RGCN 的空間幾何建模能力得以充分發揮，$R^2$ 表現因而相對穩定。此一時段性差異本身即是模型確實學到「陰影邊界＝關鍵訊號」此一物理直覺的間接證據，而非單純記憶訓練分佈。

\begin{figure}[htbp]
    \centering
    \includegraphics[width=0.8\textwidth]{../urban-thermal-gnn/04_training/viz_output/training_v5_300/figC_hourly_r2.png}
    \caption[逐時段驗證 $R^2$ 分佈]{11 個時步（08:00--18:00）之逐時段驗證 $R^2$ 分佈（V5-300）}
    \label{fig:hourly_r2}
\end{figure}


\subsection{熱壓力等級分類混淆矩陣}
\label{subsec:confusion_matrix}

圖~\ref{fig:confusion} 以 UTCI 熱壓力六等分類體系之混淆矩陣呈現分類品質。高頻等級（中等熱壓力 26--32°C、強烈熱壓力 32--38°C）之分類準確率最高，反映模型在訓練資料密集區域的優異表現；較稀有等級（極端熱壓力 $>46°C$）之準確率相對略低，源於真實場景中該等級樣本數天然較少，而非模型結構性缺陷；絕大多數誤分案例均發生於相鄰等級之臨界點，不存在跨兩個等級以上的嚴重誤判，顯示模型誤差具連續性而非系統性偏移，此點對於建築師解讀預測結果的可信度至關重要。

\begin{figure}[htbp]
    \centering
    \includegraphics[width=0.65\textwidth]{../urban-thermal-gnn/04_training/viz_output/training_v5_300/figD_confusion_matrix.png}
    \caption[熱壓力等級分類混淆矩陣]{測試集 UTCI 熱壓力六等級分類混淆矩陣（V5-300，正規化百分比）}
    \label{fig:confusion}
\end{figure}

上圖對角線元素代表各等級之正確分類率；非對角線元素代表鄰等級間之混淆。各等級準確率分別為：中等熱壓力 88.97\%、強烈熱壓力 92.90\%、極強熱壓力 95.35\%、極端熱壓力 91.50\%。


\subsection{空間熱壓力場視覺化比較}
\label{subsec:spatial_comparison}

圖~\ref{fig:spatial_comparison} 以選定測試場景展示模型預測之空間 UTCI 熱壓力場與 Radiance／EnergyPlus 模擬標籤的視覺比較。預測場與標籤場在建築陰影輪廓、高熱區域分佈及梯度走向上均高度吻合，殘差集中於建築邊緣之陰影過渡帶——此一誤差分布特性本身具設計意義：模型在「陰影是否存在」的判斷上高度可靠，誤差主要來自陰影邊界的精確位置這類次要幾何細節，對建築師進行量體配置之初期決策影響有限。

\begin{figure}[htbp]
    \centering
    \includegraphics[width=\textwidth]{../urban-thermal-gnn/04_training/viz_output/training_v5_300/figE_spatial_comparison.png}
    \caption[空間 UTCI 熱壓力場比較]{測試場景空間 UTCI 熱壓力場比較（V5-300，正午 12:00 時步）}
    \label{fig:spatial_comparison}
\end{figure}

上圖左列為物理式模擬標籤，右列為 PI-ST-GNN 預測值，下方色條代表 UTCI 值域（°C），殘差圖（右下）顯示誤差主要集中在建築邊緣之陰影過渡帶。

圖~\ref{fig:feature_fields} 進一步展示輸入特徵場與 UTCI 輸出場之空間對應關係，並沿時間軸呈現其逐時演變：天空視角因子（SVF）對 UTCI 分佈具顯著驅動作用，高 SVF（開放地面）對應高 UTCI、低 SVF（建築陰影）對應低 UTCI，此一單調對應關係之忠實還原，正是第~\ref{sec:physics_loss}節輻射一致性約束發揮作用的直接證據；而此一對應關係隨時間仍保持穩定（各列間色彩分布之相對格局一致），顯示模型並未僅學到單一時步之表面統計巧合。

\begin{figure}[htbp]
    \centering
    \includegraphics[width=\textwidth]{../urban-thermal-gnn/04_training/figures/fig_geo6_feature_fields.png}
    \caption[輸入特徵場與輸出場空間分佈]{真實場景（V5）輸入特徵場與 UTCI 輸出場空間分佈：四個代表時步（08:00、11:00、14:00、18:00）對照}
    \label{fig:feature_fields}
\end{figure}

上圖由左至右分別為天空視角因子（SVF，靜態）、空氣溫度（Ta）、平均輻射溫度（MRT）及模型預測 UTCI 值，由上至下為四個代表時步；灰色多邊形為真實建物量體，黑點標示該時步位於陰影中之空氣節點。


\section{設計選擇之消融實驗}
\label{sec:ablation}

高整體精度本身不足以說明模型「為何」有效——若不拆解各設計選擇之獨立貢獻，便無從判斷 RGCN 的異質圖結構、抑或物理引導損失，何者才是模型可信賴性的真正來源。本節以五組消融變體（A0--A6，編號沿用各變體所加入之核心模組層次，與第~\ref{sec:v4_validation}節之資料集版本編號 V1--V5 為互不相關之獨立編號體系）系統性地移除或加入各核心模組，於真實場景資料集（V5-300 相同之訓練／驗證／測試切分）上評估其獨立貢獻。

\subsection{消融實驗設計}
\label{subsec:ablation_design}

五個消融變體之設計邏輯如下：

\begin{itemize}
    \item \textbf{A0（MLP 基準線）：}以純多層感知器（MLP）替換整個 GNN 架構，直接從節點座標與氣象統計特徵映射至 UTCI 輸出，不具空間圖建模能力，作為「無空間感知」基準。
    \item \textbf{A3（GNN 無物理損失）：}保留完整 RGCN + LSTM 架構，但訓練時移除所有物理資訊損失項（$\alpha_1 = \alpha_2 = \alpha_3 = 0$），僅以資料損失 $\mathcal{L}_{\text{data}}$ 訓練。
    \item \textbf{A4（GNN + $\mathcal{L}_{\text{rad}}$）：}在 A3 基礎上加入輻射一致性損失 $\mathcal{L}_{\text{rad}}$（$\alpha_1 > 0$），驗證陰影約束的獨立貢獻。
    \item \textbf{A5（GNN + $\mathcal{L}_{\text{rad}}$ + $\mathcal{L}_{\text{temp}}$）：}在 A4 基礎上加入時序平滑性損失 $\mathcal{L}_{\text{temp}}$（$\alpha_1, \alpha_2 > 0$），驗證時序約束的額外效益。
    \item \textbf{A6（完整 PI-ST-GNN）：}在 A5 基礎上加入風場阻滯損失 $\mathcal{L}_{\text{wind}}$（$\alpha_1, \alpha_2, \alpha_3 > 0$），為本研究之完整提案模型。
\end{itemize}


\subsection{消融實驗量化結果}
\label{subsec:ablation_results}

表~\ref{tab:ablation} 彙整五個消融變體與完整模型在測試集上的量化結果：

\begin{table}[htbp]
    \centering
    \caption[消融實驗量化結果比較]{消融實驗量化結果比較（真實場景測試集，166,188 預測點次）}
    \label{tab:ablation}
    \footnotesize
    \begin{tabular}{lccccp{2.8cm}}
        \toprule
        \textbf{模型變體} & \textbf{$R^2$} & \textbf{RMSE（°C）} & \textbf{MAE（°C）} & \textbf{準確率} & \textbf{組成說明} \\
        \midrule
        A0（MLP 基準線）    & 0.9926 & 0.455 & 0.230 & 97.4\% & 無 GNN，無物理約束 \\
        A3（GNN, 無物理）   & 0.9821 & 0.708 & 0.453 & 95.1\% & RGCN+LSTM \\
        A4（$+\mathcal{L}_{\text{rad}}$） & 0.9805 & 0.740 & 0.449 & 94.9\% & A3 + 輻射約束 \\
        A5（$+\mathcal{L}_{\text{temp}}$）& 0.9777 & 0.792 & 0.496 & 94.6\% & A4 + 時序約束 \\
        A6（完整模型）      & 0.9655 & 0.984 & 0.564 & 94.0\% & A5 + 風場約束 \\
        \midrule
        \textbf{最終（V5-300）} & \textbf{0.9753} & \textbf{0.833} & \textbf{0.524} & \textbf{93.9\%} & 完整 PI-ST-GNN，主要報告模型 \\
        \bottomrule
    \end{tabular}
\end{table}

\begin{figure}[htbp]
    \centering
    \includegraphics[width=0.85\textwidth]{../urban-thermal-gnn/04_training/figures/fig8_ablation.png}
    \caption[消融實驗各變體比較圖]{消融實驗各變體 $R^2$（左）與 RMSE（右）比較圖。}
    \label{fig:ablation}
\end{figure}


\subsection{消融結果討論：數值精度與物理可信賴性的取捨}
\label{subsec:ablation_discussion}

圖~\ref{fig:ablation} 以長條圖並列呈現各變體之 $R^2$（左）與 RMSE（右）。純 MLP 基準線（A0）於兩項數值指標上取得各變體中最佳表現（$R^2=0.9926$, RMSE$=0.455$°C）；此一數值優勢之訓練細節（樣本規模、空間平滑效應等）屬模型工程層面之技術分析，完整說明見附錄~\ref{appx:ablation_numeric_detail}。就本研究之設計輔助宗旨而言，更關鍵的觀察在於：A0 缺乏物理約束，其預測可能違反輻射陰影之空間一致性（即輸出「陰影區比日照區更熱」之物理悖論），此類違反難以被 $R^2$ 或 RMSE 這類彙總指標捕捉，卻直接構成建築設計決策之誤導風險。

從 A3 至 A6 逐步加入物理損失，RMSE 呈現單調上升（A3 $0.708\to$A4 $0.740\to$A5 $0.792\to$A6 $0.984$°C），但每項約束各自針對特定的熱力學一致性問題發揮作用：$\mathcal{L}_{\text{rad}}$ 之加入（A3$\to$A4）強制模型學習陰影梯度之方向性，改善局部輻射一致性；$\mathcal{L}_{\text{temp}}$ 之加入（A4$\to$A5）防止逐時預測出現非物理跳躍，改善時序軌跡的熱慣性自洽性；$\mathcal{L}_{\text{wind}}$ 之加入（A5$\to$A6）約束高層建築背風側之熱累積極端值，改善風場阻滯情境的物理合理性。這些一致性提升難以被純數值指標捕捉，卻是建築師採信模型輸出、將其用於實際設計決策之前提。換言之，完整 PI-ST-GNN 以較程序化資料更明顯的數值精度取捨，換取設計輔助工具最不可或缺的特質：預測結果在物理上站得住腳。

須說明者：本節五個消融變體之訓練，係於真實場景資料集（V5-300 相同切分，動態邊快取已修正，\texttt{shadow}／\texttt{veg\_et}／\texttt{convective} 三類關係均透過真實邊傳遞訊息）上重新執行，故其異質圖訊息傳遞已包含全部五類關係型邊緣，不受 V1--V4 期間動態邊缺陷之影響（該缺陷之完整記錄見附錄~\ref{appx:dynamic_edges_fix}）。五個變體彼此僅在物理損失項組合上有差異，圖結構部分完全一致，故變體間之相對比較仍屬內部一致；本節 A6 變體與表~\ref{tab:main_results}之 V5-300 官方主要報告模型間之數值差異說明，見附錄~\ref{appx:ablation_numeric_detail}。


\section{推理效能對互動式設計工作流之意義}
\label{sec:performance}

\subsection{推理延遲基準測試}
\label{subsec:inference_timing}

本研究對 PI-ST-GNN 進行 200 次重複推理計時（單場景，程序化主線之固定 6,241 個空氣節點網格，11 時步），結果如表~\ref{tab:timing} 所示；此一節點數為程序化主線之固定取樣網格，V5-300 所用之真實場景節點數（136--400，第~\ref{subsec:overall_metrics}節）較此更小，故下列延遲數字可視為 V5-300 實際推理延遲之保守上界，而非低估：

\begin{table}[htbp]
    \centering
    \caption[推理延遲基準測試]{GNN 推理延遲基準測試（RTX 5070 Ti, 200 次重複）}
    \label{tab:timing}
    \small
    \begin{tabular}{lccc}
        \toprule
        \textbf{評估方法} & \textbf{平均延遲（ms）} & \textbf{最小延遲（ms）} & \textbf{最大延遲（ms）} \\
        \midrule
        PI-ST-GNN（本研究） & \textbf{61.5 ± 24.4} & 37.1 & 185.6 \\
        傳統 CFD 模擬（參考值）& $\sim$2--8 小時 & --- & --- \\
        Ladybug Tools 微氣候模擬 & 數分鐘至數十分鐘 & --- & --- \\
        \bottomrule
    \end{tabular}
\end{table}

此處關注的重點並非推理延遲的絕對數字，而在於代理模型的評估時間尺度是否落在設計迭代與最佳化搜尋可承受的範圍內。相較於傳統 CFD（單次求解約 2--8 小時）與 Ladybug Tools 微氣候模擬（數分鐘至數十分鐘）需以離線批次方式執行，代理模型將單次適應度求值壓縮至可即時互動的量級，使原本受限於模擬求值成本的兩種工作流程得以成立：其一，設計師於 Rhino／Grasshopper 環境中調整建築量體後，可在設計操作的時間尺度內取得全場域更新的 UTCI 熱圖；其二，將代理模型嵌入 NSGA-II 之適應度函數，使數百代、每代數十至數百個體之最佳化搜尋得以在可行的運算時程內完成，而非傳統模擬驅動流程所需之數週時間。本研究之貢獻在於方法工作流程之整合可行性，而非單純追求運算速度之極大化。


\subsection{訓練資料集覆蓋範圍驗證}
\label{subsec:coverage}

為驗證 300 個真實訓練場景對都市形態空間的覆蓋充分性，本研究分析訓練集在主要幾何特徵（FAR、BCR、建築棟數、樓高）上的統計分佈。由於場景直接取自真實 OSM 建物量體而非法規約束下之程序化生成，其分布較程序化幾何遠為寬廣且不均勻：FAR 中間 90\% 範圍（第 5--95 百分位）為 0.54--12.53（全距 0.23--40.83，涵蓋自低密度郊區至高密度都市核心之真實建物量體），BCR 中間 90\% 範圍為 0.18--1.00（全距 0.08--1.00），每場景建物棟數 3--208 棟（平均 11.3 棟），建物簷高 3--162 公尺（平均 12.6 公尺）。此一遠較程序化幾何寬廣之真實形態覆蓋範圍，是第~\ref{chapter:results}章以真實都市建物量體進行案例驗證、以及第~\ref{sec:nsga2}節 NSGA-II 最佳化搜尋時，判斷結果是否落入分布外推（Out-of-Distribution）風險的重要參照基準——第~\ref{subsec:pareto_morphology}節之最佳化結果所落之 FAR 區間即落於此訓練分布之中間範圍內。


\section{真實場景資料集：版本演化與模擬引擎驗證}
\label{sec:v4_validation}

本研究之真實場景資料集與模擬引擎歷經五個版本之演化，於此簡述以說明實驗歷程與本研究最終採用 V5-300 為主要報告模型之原因。V1 以八維空氣節點特徵於程序化幾何上建立基礎預測能力；V2 增列地表溫度為第九維特徵；V3 進一步納入街谷高寬比（$H/W$）為第十維特徵；上述三版皆以程序化蒙地卡羅幾何與典型氣象年為訓練資料。V4 為資料版本之根本轉變——由程序化幾何全面改採真實 OSM 建物、ETH 樹冠幾何，並以真實環境部 IoT 測站之逐時實測氣溫濕度驅動模擬（場址、建物、樹冠與 IoT 氣象等地理資料之完整說明見第~\ref{subsec:v4_real_dataset}節），惟其微氣候模擬仍採自製 Python 物理近似引擎。V5 在完全承襲 V4 真實地理資料與氣象邊界條件之基礎上，將模擬引擎全面替換為 Radiance／EnergyPlus 官方模擬流程（第~\ref{sec:ladybug_workflow}節），並同時修正下述動態邊缺陷，是本研究最終採用之主要報告模型，其完整訓練與測試結果已於第~\ref{sec:evaluation}節（表~\ref{tab:main_results}、圖~\ref{fig:training_curves}--\ref{fig:spatial_comparison}）報告，本節不再重複列出。

須揭露者：V1--V4 訓練時所用之異質圖動態邊快取（\texttt{shadow}／\texttt{veg\_et}／\texttt{convective} 三類關係）與程序化主線曾發生之缺陷同屬一類——快取檔案從未於真實場景資料集上被實際建立，故 V1--V4 之異質圖訊息傳遞實際上僅透過 \texttt{semantic} 與 \texttt{contiguity} 兩類靜態邊進行；V4 當時表面報告之測試集數值（$R^2 \approx 0.99$ 量級）即在此一未修正狀態下取得，不具方法完整性，不應作為有效驗證結果引用。此一缺陷直到 V5-300 以 \texttt{16\_build\_dynamic\_edge\_cache\_v5.py}（真實逐場景太陽方位角、真實逐月中央氣象署風向）重新建構動態邊快取後方才修正，完整技術記錄見附錄~\ref{appx:dynamic_edges_fix}。

圖~\ref{fig:v5_comparison} 就數個未見過之真實測試場景，並列呈現其真實 OSM 幾何、Radiance／EnergyPlus 模擬之真值 UTCI 場、V5-300 模型預測場與逐點誤差，作為表~\ref{tab:main_results}量化指標之空間視覺化補充。

\begin{figure}[htbp]
    \centering
    \includegraphics[width=\textwidth]{../urban-thermal-gnn/04_training/figures/fig_v5_scene_comparison.png}
    \caption[真實測試場景之預測比較]{以真實地理環境與 Radiance／EnergyPlus 引擎訓練之 PI-ST-GNN（V5-300）於未見過真實測試場景之預測比較（14:00）：由左至右為真實 OSM 幾何、真值 UTCI（Radiance／EnergyPlus 模擬）、模型預測、逐點誤差。幾何欄之虛線框標示 $80\times80$ 公尺基地邊界，實色為基地內建物與喬木，淡色為基地外仍參與遮蔭運算之周邊紋理脈絡。}
    \label{fig:v5_comparison}
\end{figure}

須對照說明者：V5-300 之 $R^2 = 0.9753$ 略低於本研究另一條程序化主線（第~\ref{sec:evaluation}節前述、以合成幾何訓練）之 $0.9956$，此差異為合理且可預期——Radiance／EnergyPlus 模擬之空間變異本質上較自製近似式更為劇烈與寫實（局部 MRT 極端值可達 77\,$^\circ$C 以上），且建物外殼構造採 ASHRAE 標準假設而非逐棟實測材料，真實 OSM 幾何與真實 IoT 觀測亦含更高之雜訊與異質性（如樓高標籤稀疏、單一氣象站空間代表性等）。然而模型於\textbf{未見過之真實場景}上仍維持 RMSE 僅 $0.83\,^\circ\text{C}$、遠小於 UTCI 熱壓力等級間距，此為模型於真實地理環境與物理模擬雙重條件下之外部效度直接證據，亦是本研究方法論上最嚴謹之驗證結果，故雖數值指標略低於程序化主線，仍以 V5-300 作為全文主要報告模型。


\section{小結}
\label{sec:chapter4_summary}

本章以測試集量化評估、消融實驗與推理效能基準測試，驗證了 PI-ST-GNN 作為都市微氣候代理模型之有效性與可信賴性。本研究最終採用之主要報告模型（V5-300：真實地理資料 $+$ Radiance／EnergyPlus 模擬引擎 $+$ 已修正動態邊快取，附錄~\ref{appx:dynamic_edges_fix}）在 166,188 個真實場景測試樣本上達到 $R^2=0.9753$、RMSE$=0.833$°C、熱壓力等級分類準確率 93.9\%，為本研究方法論上最嚴謹之驗證結果（第~\ref{sec:v4_validation}節）；消融實驗進一步揭示，純數值指標上表現最佳的 MLP 基準線，缺乏圖結構與物理約束所提供的空間因果建模能力與預測物理一致性，而完整 PI-ST-GNN 以輕微數值犧牲換取了設計輔助工具不可或缺的熱力學可信賴性；代理模型將單次適應度求值壓縮至可即時互動之時間量級，則為第~\ref{chapter:results}章之多目標最佳化與熱舒適人行動線評估，提供了得以整合入設計操作流程之核心運算引擎。

\chapter{結果分析與設計應用討論}
\label{chapter:results}

\section{整合系統設計流程回顧}
\label{sec:system_review}

第三、四章分別建立了 PI-ST-GNN 之方法論架構與訓練驗證結果。本章聚焦於整合系統在\textbf{實際設計情境}下的應用操作成果，包含：（1）多目標最佳化帕累托前緣之設計方案解讀，涵蓋熱舒適人行動線設計目標；（2）第~\ref{sec:building_extraction}節真實建物擷取流程於既有場地案例之驗證表現；以及（3）系統整合架構之效能瓶頸與最佳化方向討論。


\section{多目標最佳化設計方案分析}
\label{sec:pareto_analysis}

\subsection{NSGA-II 收斂過程與帕累托前緣形成}
\label{subsec:nsga2_convergence}

以訓練完成之 PI-ST-GNN 作為評估代理，NSGA-II 在 $80 \times 80$ 公尺場地尺度下同時最佳化第~\ref{subsec:design_variables}節定義之三項目標：$f_1$（最小化全場域時空平均 UTCI）、$f_2$（最小化人行動線時空平均熱暴露 $\bar{\Phi}_{\text{walkway}}$）與 $f_3$（最大化綠化率）。演算法設定如下：

\begin{itemize}
    \item \textbf{族群大小：}50 個個體
    \item \textbf{演化代數：}100 代
    \item \textbf{交叉算子：}模擬二進位交叉（SBX，交叉概率 0.9，$\eta_c = 15$）
    \item \textbf{突變算子：}多項式突變（突變概率 $1/n_{\text{var}}$，$\eta_m = 20$）
    \item \textbf{選擇策略：}可行性優先 + 非支配排序 + 擁擠距離選擇
\end{itemize}

由於每次個體評估之 GNN 推理已能提供立即性回饋（詳細延遲數值見附錄~\ref{appx:latency_breakdown}），100 代 $\times$ 50 個體規模的完整最佳化過程理論上可在數分鐘內完成，相較第~\ref{subsec:wallacei_comparison}節所述、直接呼叫高擬真度模擬引擎之 Wallacei 案例常需數週至數月的收斂時間，此一運算時間尺度的壓縮，是本研究得以在方案概念階段即進行多目標搜尋、而非僅能於方案定案後事後驗證的關鍵前提。

在正式導入第三項目標 $f_2$ 之前，本研究首先以一次僅含 $f_1$、$f_3$ 二目標之基準實驗（族群 40、演化 50 代、亂數種子 42，V5-300 checkpoint，驗證 $R^2=0.9837$），逐代真實呼叫 GNN 推理，於 GPU（NVIDIA RTX 5070 Ti）上共 404.2 秒完成，結果如圖~\ref{fig:nsga2_real}(a) 所示：族群最佳全域平均 UTCI 由第 1 代之 $36.29^\circ\text{C}$ 逐步下降至第 50 代之 $36.25^\circ\text{C}$，族群最佳綠化率則由 $0.009$ 上升至 $0.048$，可行解比例由 38/40 提升並穩定於 40/40。此一基準實驗之帕累托前緣形態，將於第~\ref{subsec:pareto_morphology}節詳述——與第~\ref{sec:v4_validation}節所述動態邊缺陷尚未修正之早期版本模型（V1--V4）之基準實驗結果不同，本次以 V5-300 執行之結果並未出現退化解現象，其原因與意義一併於下節討論。

\begin{figure}[htbp]
    \centering
    \includegraphics[width=\textwidth]{../urban-thermal-gnn/04_training/figures/fig_nsga2_convergence.png}
    \caption[基準二目標 NSGA-II 收斂結果]{基準二目標 NSGA-II 收斂結果（PI-ST-GNN V5-300 checkpoint 作為適應度求值器，族群 40、演化 50 代）}
    \label{fig:nsga2_real}
\end{figure}

上圖 (a) 為逐代族群最佳全域 UTCI（紅）與最佳綠化率（綠）之收斂曲線；(b) 為第 50 代最終帕累托前緣（$n=8$ 非支配解），色階代表各解之容積率（FAR）。


\subsection{帕累托前緣之真實形態與現行目標函數之侷限}
\label{subsec:pareto_morphology}

圖~\ref{fig:nsga2_real}(b) 之最終帕累托前緣揭示一項與本研究先前（V1--V4，動態邊缺陷尚未修正）基準實驗明顯不同之發現：6 個非支配解之 FAR 值收斂於 $2.03$--$2.05$ 之合理中高密度區間，落於第~\ref{subsec:coverage}節真實場景訓練集 FAR 分布之中段範圍內，\textbf{並未}出現先前版本模型所觀察到之「幾乎不興建任何建築量體」退化解（Degenerate Solution）現象。

此一對比之根本原因，最可能與第~\ref{sec:v4_validation}節、附錄~\ref{appx:dynamic_edges_fix}所揭露之動態邊缺陷有關：V1--V4 因動態邊快取從未真正建立，模型從未透過 \texttt{shadow}、\texttt{veg\_et} 兩類關係型邊實際學習建物遮蔭與植栽降溫之因果效應——即便 $f_1$（全域平均 UTCI）目標理論上應誘發模型辨識建物與喬木之降溫貢獻，一個從未見過真實遮蔭／降溫訊號傳遞路徑的模型，自然難以在其預測輸出中忠實反映此一效應，使最佳化引擎在缺乏正向誘因之情況下傾向「不蓋房子、盡量種樹」之退化解；V5-300 修正此一缺陷後，模型得以透過真實建立之動態邊學習建物與喬木對周邊空氣節點之實際降溫效果，NSGA-II 因而在同一二目標設定下即能辨識「適度建築量體可提供遮蔭、降低全域平均 UTCI」之正向訊號，收斂至合理密度區間而非退化解。換言之，先前版本所觀察到之退化解，其成因較可能主要並非二目標函數設計本身之結構性缺陷，而是動態邊缺陷導致模型未能學習建物／植栽之真實降溫因果關係；此一發現同時佐證了異質圖顯式邊結構（相對於純黑盒模型）於此類應用之關鍵價值——邊結構缺陷可被直接追溯並修正，且修正後之因果學習效果可於下游最佳化任務之行為變化中被觀察與驗證，此為本研究可解釋性設計（第~\ref{subsec:framework_interpretability}節）之具體效益展示。

儘管如此，現有 5 維違規向量（FAR、BCR、縮退、容納性、碰撞）仍僅檢核上限，未設定最低開發強度之下限約束，此一結構性缺口並未因動態邊修正而消失，僅是在 V5-300 修正後之模型行為下未被觸發；為避免此一風險於其他場地條件、氣候情境或搜索設定下重新浮現，第~\ref{sec:limitations}節與第~\ref{sec:future_work}節仍建議將最低容積率或建築棟數下限之可行性約束，列為後續強化措施。

此一發現對後續研究仍有直接啟示：第~\ref{subsec:from_astar_to_walkway}節與附錄~\ref{appx:walkway_derivation}所定義之第三項目標 $f_2$（人行動線熱暴露 $\bar{\Phi}_{\text{walkway}}$）之數學形式雖已完整推導，其研究動機——全空間平均指標難以反映行人沿路徑之動態熱暴露歷程——獨立於退化解問題而成立（第~\ref{sec:problem}節第二項研究缺口）；第~\ref{subsec:abc_comparison}節即以真實三目標實驗，驗證 $f_2$ 整合後之路徑感知評估能力與其對帕累托前緣量體分布之影響，並非作為退化解之修正手段，而是本研究對「路徑尺度熱舒適評估」此一研究缺口之直接回應。


\subsection{三目標實驗：路徑配置之三方案比較}
\label{subsec:abc_comparison}

為將第三項目標 $f_2$（人行動線熱暴露）實際整合進最佳化流程並驗證其路徑感知評估能力，本研究將 $f_2$ 正式接入 \texttt{FitnessEvaluator}（第~\ref{sec:nsga2}節），並針對同一 $80\times80$ 公尺場地與建築演化搜索空間，設計 A、B、C 三種固定人行路線配置，各自執行一次完整之三目標 NSGA-II 最佳化（族群 40、演化 50 代、亂數種子 42，V5-300 checkpoint，驗證 $R^2=0.9837$）：路線 A 沿建築臨街面配置、路線 B 貫穿開放廣場／中庭、路線 C 對角線橫越基地。此一設計直接回應口試委員對「建築形態與人行路徑之互動」須有具體比較實驗、而非僅呈現單一最佳化結果之要求。

三次獨立實驗之結果如圖~\ref{fig:abc_comparison} 所示，可歸納三項具體發現。其一，\textbf{三目標帕累托前緣呈現更寬廣之量體權衡區間}：三組實驗最終帕累托前緣之 FAR 範圍分別為路線 A（1.82--2.14）、路線 B（0.62--1.78）、路線 C（1.75--2.01），與第~\ref{subsec:pareto_morphology}節二目標基準實驗收斂於單一 FAR 區間（2.03--2.05）相比，三目標版本之非支配解涵蓋明顯更寬之量體區間，尤其路線 B 之帕累托前緣向下延伸至 FAR 0.62，顯示人行動線熱暴露目標確實為最佳化引擎引入了額外的量體—路徑權衡維度：部分非支配解以降低建築量體換取路徑熱舒適之改善，此一權衡在僅有 $f_1$、$f_3$ 之二目標設定下並不存在。三組實驗之 FAR 範圍均落於第~\ref{subsec:coverage}節真實場景訓練集 FAR 分布之涵蓋範圍內（中間 90\% 為 0.54--12.53），故最佳化搜尋未落入分布外推風險。

其二，\textbf{不同路線配置產生可辨識之最佳化差異}：路線 B（開放廣場路線）之最終最佳人行動線熱暴露（$\bar{\Phi}_{\text{walkway}}=0.336$）優於路線 C（$0.341$）與路線 A（$0.342$），且路線 B、C 之最終帕累托前緣規模（均為 $n=40$）亦大於路線 A（$n=20$），顯示當人行路線通過場地中央開放空間時，最佳化引擎有更大的建築配置自由度可同時兼顧陰影遮蔽與路徑熱舒適，而路線 A 之臨街面路線因路徑貼近場地邊界，可調整之建築退縮與遮蔭配置選項相對受限。此一差異證實了本研究之核心論點：人行路線之選址本身即是影響都市設計最佳化結果的關鍵變數，而非與建築量體配置無關之獨立參數。

其三，\textbf{全域 UTCI 目標與路徑選擇之交互作用}：三路線之最佳全域平均 UTCI 差異相對較小（$36.25$--$36.33^\circ\text{C}$），顯示 $f_1$ 對路線選擇之敏感度低於 $f_2$，此一現象合乎預期——$f_1$ 為全場域空間平均，對局部路徑之幾何差異本不敏感，此正是第~\ref{sec:problem}節第二個研究缺口所述「全空間平均最佳化難以反映行人沿路徑之動態熱暴露歷程」之量化實證：若僅以 $f_1$ 作為評估指標，三條路線之設計方案將被判定為近乎等價，唯有透過 $f_2$ 之路徑感知評估，才能區辨出路線 B 之顯著優勢。

\begin{figure}[htbp]
    \centering
    \includegraphics[width=\textwidth]{../urban-thermal-gnn/04_training/figures/fig_abc_comparison.png}
    \caption[A/B/C 三種人行路線配置之三目標 NSGA-II 比較實驗]{A/B/C 三種人行路線配置之三目標 NSGA-II 比較實驗：上排為三條路線於場地上之幾何配置與各自最佳全域 UTCI 解之建築／喬木佈局；下排 (a)(b) 為全域 UTCI 與人行動線熱暴露之收斂曲線，(c) 為三組實驗最終代帕累托前緣之散點對照。}
    \label{fig:abc_comparison}
\end{figure}

需說明之範疇限制：本實驗之路線為固定輸入（非最佳化搜尋對象），驗證的是「已推導之 $\bar{\Phi}_{\text{walkway}}$ 公式與三目標形式化在不同路線幾何下之有效性與區辨力」，而非完整之多路徑 $A^*$ 動線圖同時最佳化；後者仍列為第~\ref{sec:future_work}節之後續工作。


\subsection{設計疊代示範：從帕累托前緣到方案選擇}
\label{subsec:design_iteration_demo}

第~\ref{subsec:abc_comparison}節之三目標最佳化並不輸出單一「最優方案」，而是輸出一組非支配之候選方案集合，實際的設計決策仍須由設計者依當下案址之優先考量從中選擇。圖~\ref{fig:design_iteration} 以路線 B（開放廣場路線）最終帕累托前緣中兩個真實非支配解為例，示範此一比較與選擇過程：方案 1 為前緣上全域平均 UTCI 最低之解（$36.26^\circ\text{C}$），其建築量體集中且 FAR 較高（$1.78$）；方案 2 為前緣上人行動線熱暴露最低之解（$\bar{\Phi}_{\text{walkway}}=0.336$），其建築量體配置退離路線兩側、FAR 較低（$0.71$），以犧牲部分全域平均表現（UTCI 上升至 $36.33^\circ\text{C}$）換取路線熱舒適度之改善。此一比較呈現了本系統作為設計輔助工具之核心互動模式：設計者輸入建築輪廓與人行動線後，工具並非提供單一線性答案，而是提供一組可供權衡比較的候選方案，設計者可依據案址之實際優先權重（例如行人流量較高的路段應優先考慮方案 2 之配置邏輯）進行選擇與後續之量體調整，形成第~\ref{sec:framework}節所述「輸入幾何 → 運算輸出 → 設計回饋調整」之閉環決策過程。

\begin{figure}[htbp]
    \centering
    \includegraphics[width=\textwidth]{../urban-thermal-gnn/04_training/figures/fig_design_iteration.png}
    \caption[設計疊代示範：從建築輪廓輸入到計算輸出、再到設計回饋調整之閉環]{設計疊代示範：以路線 B 帕累托前緣中兩個真實非支配解為例，示範建築量體配置在全域熱舒適與人行動線熱舒適兩種優先考量下之具體權衡差異，模擬設計者於候選方案間比較、選擇並回饋調整之閉環決策情境。}
    \label{fig:design_iteration}
\end{figure}


\section{真實建物擷取流程之案例驗證}
\label{sec:real_case_validation}

第~\ref{subsec:extraction_vs_parametric}節說明了真實建物擷取流程與隨機生成場景之互補分工：前者用於真實都市紋理下之案例驗證，後者支撐模型訓練所需之大規模樣本集。本節以第~\ref{subsec:vectorize_method}節三服務融合方法擷取之真實建物量體，作為 PI-ST-GNN 之推理輸入，檢視模型在未見過真實幾何配置下的表現。

由於本案例擷取之建物幾何（棟距、樓高分布、平面形態）與訓練資料集之 300 個真實場景分布未必完全重疊，此一驗證本質上是對模型一般化能力的分布外（Out-of-Distribution）壓力測試，而非對訓練集內插能力的重複驗證。第~\ref{subsec:coverage}節所述之訓練集幾何特徵覆蓋範圍（FAR 中間 90\% 為 0.54--12.53、BCR 中間 90\% 為 0.18--1.00、簷高 3--162m），可作為判斷真實案例是否落於模型可靠一般化區間的參照基準：若真實場地之量體統計特徵落於此範圍內，可對模型輸出有較高信心；若明顯超出，則預測結果應視為初步參考，建議搭配傳統模擬工具進行交叉驗證。


\section{系統整合架構效能討論}
\label{sec:system_discussion}

\subsection{FastAPI/WebSocket 通訊延遲}
\label{subsec:api_latency}

整合系統的端到端延遲由幾何序列化（Grasshopper → Python）、GNN 推理與結果回傳三個環節組成，合計落在低延遲等級（詳細拆解與各環節毫秒數見附錄~\ref{appx:latency_breakdown}），遠低於人眼察覺延遲（$\sim$1 秒）的門檻，能提供令設計師滿意的「近即時」互動體驗，並支撐第~\ref{subsec:abc_comparison}節所述、涉及數十代族群演化的 NSGA-II 搜尋於可行時程內完成。


\subsection{使用者操作流程概覽}
\label{subsec:user_workflow_overview}

前述效能數字所服務的最終對象，是實際操作 Grasshopper 的設計者；系統是否真正達成即時閉環設計，最終仍須以使用者從「周遭環境模型建好」到「拿到熱舒適評估」之間實際要走的步驟來檢驗，而非僅以延遲毫秒數衡量。圖~\ref{fig:user_workflow} 將本研究實作之 \texttt{UTCIPredictor}（單次即時預測）與 \texttt{UTCIOptimizer}（多目標最佳化搜尋）兩個 Grasshopper 元件之操作流程，統整為單一之 1--8 步驟流程圖：步驟 1--2 為一次性前置準備（啟動背景運算服務、安裝通訊套件）；步驟 3--5 為場地幾何準備與元件連接（匯入基地與建物幾何、接上輸入端口、送出計算請求）；步驟 5 之後依應用情境分岔為路徑 A（單次即時評估，提供立即性回饋、可邊調整設計邊觀察熱力圖）與路徑 B（多目標最佳化搜尋，背景執行緒進行、不阻塞 Rhino 操作、可隨時取消）；步驟 8 為結果讀取與設計回饋，兩路徑之輸出（熱力圖網格、數值指標、帕累托候選方案）皆直接呈現於設計者當下操作之 Rhino 場景中，不另外產生 PDF 報告。

\begin{figure}[htbp]
    \centering
    \includegraphics[width=\textwidth]{../urban-thermal-gnn/04_training/figures/fig_user_workflow.png}
    \caption[Grasshopper 元件操作流程（1--8 步驟）]{Grasshopper 元件操作流程（1--8 步驟）：資料準備（步驟 1--3）→ 匯入元件（步驟 4--5）→ 計算（路徑 A／B，步驟 6--7）→ 結果讀取與設計回饋（步驟 8）。}
    \label{fig:user_workflow}
\end{figure}

完整之逐步操作手冊、輸入輸出端口對照表與常見狀況排除，詳見附錄~\ref{appx:user_manual}。


\subsection{物理一致性與建築設計輔助決策模型的可解釋性}
\label{subsec:trustworthiness}

本系統在建築設計輔助工具的脈絡中，預測結果的\textbf{物理合理性（Physical Plausibility）}比純粹的數值精度更為重要 \citep{shan2025physics}。一個預測「陰影區域比日照區域更熱」的 AI 模型，即便整體 $R^2$ 很高，也可能在關鍵設計決策點上誤導建築師。

物理資訊損失函數在以下具體情境中的功效：

\begin{enumerate}
    \item \textbf{陰影形態評估：}$\mathcal{L}_{\text{rad}}$ 確保模型始終正確反映陰影帶的降溫效果，使建築師在評估不同立面開口率或遮陽板配置時，獲得符合直覺的熱舒適反饋。
    \item \textbf{日間熱積累趨勢：}$\mathcal{L}_{\text{temp}}$ 確保 UTCI 預測曲線符合都市熱慣性的物理規律——早晨緩升、午後達峰、傍晚緩降，使設計師能正確評估不同朝向建築的日間熱積累差異。
    \item \textbf{高層建築風場影響：}$\mathcal{L}_{\text{wind}}$ 確保高層建築背風側的「熱死角」能被正確識別，這是建築設計中常被忽略但對行人熱舒適影響顯著的空氣動力學現象 \citep{oke1987boundary}。
\end{enumerate}

上述物理一致性的維持，使得 PI-ST-GNN 不僅是一個「黑箱預測器」，更是一個\textbf{具有熱力學解釋能力的設計輔助工具}，能在設計師的專業判斷基礎上提供可驗證的科學依據。

為進一步驗證此一物理一致性主張並非僅為訓練指標上的巧合，本研究以 GNNExplainer 對本研究最終採用之 V5-300 模型進行事後可解釋性分析 \citep{shan2025physics}，結果如圖~\ref{fig:gnnexplainer} 所示。以真實場景資料集中一處熱壓力峰值節點（$\text{UTCI}=41.7°\text{C}$）為分析目標，梯度顯著性分析顯示該節點之預測值主要由其空間鄰域中的空氣溫度 $T_a$（正規化歸因值 1.00）主導，平均輻射溫度 $T_{\text{mrt}}$ 居次（0.79），相對濕度 RH 居第三（0.53），陰影旗標 $f_{sh}$、風速 $v_a$ 與天空視角因子 SVF 則為次要貢獻，此一特徵重要性排序反映真實場景中氣溫與輻射環境對熱壓力之主導地位，與都市微氣候學之物理直覺一致。

RGCN 各層之關係權重范數分析（附錄~\ref{appx:dynamic_edges_fix}記錄之真實場景資料集動態邊修正，五類關係於訓練時皆具真實動態邊）顯示，五類關係之權重范數皆落於相近量級（見圖~\ref{fig:gnnexplainer}(c)），未呈現單一關係獨大之型態，且\textbf{不再出現任何「無真實訓練邊」之虛線填充樣式}——此為 V5-300 動態邊缺陷修正後之直接可視化證據：淺層（第 1--2 層）中各關係權重相近，反映模型在淺層同時整合多種物理管道之訊號；深層則呈現 \texttt{contiguity}（空間毗鄰）權重上升之分化趨勢，與空間毗鄰關係需經多跳擴散才能涵蓋建築陰影最大投射距離（第~\ref{subsec:rgcn}節）之設計動機相符。此一分佈型態較單一關係獨大更符合物理直覺：陰影、植被蒸散、對流、語義、毗鄰五種因果機制理應共同、而非由單一路徑主導地貢獻於熱壓力預測。此一可解釋性驗證，補強了單純數值指標難以呈現的「模型是否學到正確物理機制」此一問題的實證依據。

\begin{figure}[htbp]
    \centering
    \includegraphics[width=\textwidth]{../urban-thermal-gnn/04_training/figures/fig_gnnexplainer.png}
    \caption[GNNExplainer 可解釋性分析]{PI-ST-GNN 之 GNNExplainer 事後可解釋性分析（真實場景資料集場景 4，V5-300 checkpoint，驗證 $R^2=0.9837$）}
    \label{fig:gnnexplainer}
\end{figure}

上圖 (a) 為目標節點（黑圈標示，UTCI=41.7°C）周邊之空間梯度顯著性，顏色越深代表該節點對目標預測值之影響越大；(b) 為九維輸入特徵之歸因排序，空氣溫度為主導特徵；(c) 為三層 RGCN 各關係類型之權重 Frobenius 范數，五類關係皆具真實訓練邊、范數量級相近；(d) 為目標節點兩跳鄰域內前 40 條重要邊之子圖，邊寬與顏色深淺代表梯度歸因強度。


\subsection{系統局限性與誤差來源}
\label{subsec:limitations}

儘管本系統取得優異的整體預測性能，仍存在若干局限性值得說明：

\begin{enumerate}
    \item \textbf{氣候情境涵蓋範圍之限制：}本研究另一條程序化主線之訓練資料集僅涵蓋夏至（6月21日）之單一典型氣候日；第~\ref{sec:v4_validation}節之主要報告模型（V5-300）真實資料集雖已改以 2025 年 6、7、8 月之真實逐時觀測氣象驅動，將氣候情境由單日典型值擴展為真實夏季逐時變異，但仍未涵蓋冬季東北季風、梅雨季陰天等其他季節條件，模型對此類條件之一般化能力尚待驗證；若應用於其他季節，建議以對應季節之真實觀測或 EPW 重新微調。
    \item \textbf{平流與濕球效應未建模：}現有模型未明確建模都市熱島效應中的大尺度平流傳熱，以及濕地、噴泉等蒸發冷卻設施的局部降溫效果，此類效應在極端熱壓力情境下可能不可忽略。
    \item \textbf{極端熱壓力等級準確率偏低：}如消融實驗所示，極端熱壓力（UTCI $> 46°\text{C}$）之分類準確率明顯低於其他等級，主因是訓練資料中此等級樣本佔比極低（僅出現於正午直射、高 SVF 條件下），建議後續研究採用加權採樣或針對高溫情境進行額外資料擴增。
    \item \textbf{單一行區尺度限制：}本研究聚焦於單一 $80 \times 80$ 公尺行區，尚未考量相鄰行區間的微氣候互動（如建築群之間形成的峽谷效應或大尺度通風廊道），多行區聯合建模為未來重要研究方向。
    \item \textbf{真實建物擷取之樓高辨識限制：}第~\ref{subsec:vectorize_method}節之光學文字辨識於複合式大型建物之並列樓高分區仍偶有辨識失敗或分區未完全切割之情形（已於流程中如實標記為待覆核案例，而非以猜測值填補），此類案例於本章之真實案例驗證中須人工複核樓高屬性後方可採用。
\end{enumerate}

\chapter{結論與未來研究}
\label{chapter:conclusion}

\section{研究結論}
\label{sec:conclusions}

本研究成功開發並驗證了一套整合物理資訊機器學習、異質圖神經網路、真實都市 GIS 建物擷取與多目標進化演算法的都市設計優化框架，並以 Rhino/Grasshopper 介面實現了從幾何建模到微氣候評估的閉環即時決策系統。以下針對各核心模組之研究結論進行系統性總結：

\noindent\textbf{一、真實都市建物擷取流程之建立}

本研究建立一套三服務融合之 GIS 建物向量化方法，整合 NLSC TOPO01K（樓高標註）、B5000（幾何驗證）與 EMAP（面積驗證）三項政府開放圖資，取代單一仰賴蒙地卡羅隨機採樣之隨機幾何生成方式，使模型得以在真實都市紋理下進行案例驗證，而非僅於程序化生成之封閉解空間內自我評估。

\noindent\textbf{二、異質圖建構之有效性}

本研究提出以五類關係型邊緣（陰影、植被蒸散、對流、語義、毗鄰）建立的異質圖結構，能有效編碼都市行區中建築量體與開放空間之物理因果關係。$N_{\text{obj}}$ 索引偏移機制實現了建物節點與 6,241 個空氣評估節點在單一張量空間中的統一處理，使 RGCN 能以單次前向傳播完成跨節點類型的關係推理。

\noindent\textbf{三、PI-ST-GNN 模型性能}

在 300 個隨機生成都市場景、共計 3,707,154 個預測點次的測試集評估中，PI-ST-GNN 達到：
\begin{itemize}
    \item 決定係數 $R^2 = 0.9957$，遠超代理模型部署門檻 $R^2 > 0.99$；
    \item 均方根誤差 RMSE $= 0.194°\text{C}$（去正規化後之實際 UTCI 單位），遠小於 UTCI 熱壓力等級最小間距（$6°\text{C}$）；
    \item 熱壓力等級分類準確率 98.3\%，確保預測值能可靠地對應至設計決策所需的六等熱壓力分級體系。
\end{itemize}

\noindent\textbf{四、物理資訊損失函數之作用}

消融實驗證實，三項物理資訊損失（輻射一致性 $\mathcal{L}_{\text{rad}}$、時序平滑性 $\mathcal{L}_{\text{temp}}$、風場阻滯 $\mathcal{L}_{\text{wind}}$）在純數值指標上的貢獻並非單調遞增，但其核心價值在於確保預測結果的\textbf{熱力學合理性}，防止模型輸出違反物理常識的「陰影區比日照區更熱」等悖論，提升模型作為建築設計輔助工具的\textbf{可信賴性}。

\noindent\textbf{五、運算效能對即時設計工作流之解鎖效果}

PI-ST-GNN 的推理延遲僅 \textbf{61.5 ms}（RTX 5070 Ti 上測得，200 次平均），相較傳統 CFD 模擬（2--8 小時）效率提升達 $10^5$ 倍量級。此速度量級的突破直接使以下新型設計互動模式成為可能：在 Grasshopper 環境中即時更新全場域 UTCI 熱圖、NSGA-II 100 代演化在 5 分鐘內完成，相較直接呼叫高擬真度模擬引擎之既有 Wallacei 案例常需數週至數月，本研究之最佳化迴圈得以在設計對話的時間尺度內完成。

\noindent\textbf{六、多目標最佳化的設計決策支援與其現行侷限}

以訓練完成之 PI-ST-GNN 作為適應度求值器，NSGA-II 已實際完成一次真實最佳化運行（族群 40、演化 50 代），在全域 UTCI 最小化與綠化率最大化二項目標間求解帕累托前緣；人行動線熱暴露最小化（第三項目標 $f_2$）已於第~\ref{subsec:from_astar_to_walkway}節完成數學定義與推導，取代僅評估單一起訖點路徑的傳統尋路演算法，但其整合進適應度求值器之工程實作列為近期後續工作。真實運行結果同時誠實地揭露：現行二目標函數在缺乏最低開發強度約束下會收斂至近乎零建築量體之退化解（第~\ref{subsec:pareto_morphology}節），此一發現本身即是本研究對「以 GNN 代理模型驅動之多目標最佳化」在實務部署前必須解決之目標函數設計問題的具體貢獻。


\section{研究貢獻與創新點}
\label{sec:contributions}

本研究之學術貢獻與實務創新可歸納為以下四點：

\begin{enumerate}
    \item \textbf{真實圖資交叉驗證之都市建物擷取方法：}
    本研究提出三服務融合（TOPO01K、B5000、EMAP）之建物向量化方法，以粗細線分離、子區塊輪廓萃取與多字元群集光學文字辨識，解決單一圖資來源無法同時提供精確幾何輪廓與樓層屬性之限制，並以跨服務像素重疊率提供每筆建物幾何之信心度評估，使真實都市量體得以系統性地整合進微氣候預測工作流。

    \item \textbf{延伸物理一致之異質時空圖神經網路架構，並補以物理資訊軟約束損失：}
    本研究之 RGCN 關係型聚合公式與 LSTM 暖啟動機制，採用並延伸 \citet{xin2025urbangraph} 之 UrbanGraph 架構設計（詳見第~\ref{subsec:rgcn}、\ref{subsec:lstm}節）。在此基礎上，本研究之貢獻具體體現於三點：其一，將節點集合顯式區分為建物與空氣兩種異質節點類型，並以獨立編碼器分別處理其靜態幾何與動態氣象特徵，區別於原架構之單一節點類型設計；其二，於 UrbanGraph 原有之圖拓撲硬約束（Hard Constraints）之外，另引入三項\textbf{物理資訊軟約束損失}（輻射一致性、時序平滑、風場阻滯），作為訓練過程中的補充監督機制；其三，以\textbf{多步驟時序預測}一次輸出全日 11 時段 UTCI 場，並整合至第~\ref{chapter:results}章之端到端低延遲推理與設計迴圈中。不同於現有研究多以純資料驅動方式預測單一時步的熱舒適指標 \citep{wu2023data}，本研究之上述三項擴展在應用完整性上有所補強。

    \item \textbf{低延遲代理模型實現即時閉環設計：}
    本研究建立了首個推理延遲低至 61.5 ms、並直接嵌入 Rhino/Grasshopper 設計環境的城市微氣候預測系統。此一延遲水準使得 NSGA-II 進化演算法能在合理的設計時間內（$\approx$5 分鐘）完成百代優化，從根本上改變了性能導向設計（Performance-Driven Design）在微氣候評估環節的工作流程，相較既有以 Wallacei 直接呼叫高擬真度模擬引擎之作法，大幅壓縮了最佳化迴圈所需的等待時間。

    \item \textbf{熱舒適人行動線設計之多目標決策框架：}
    本研究將 GNN 預測之 UTCI 場沿設計者定義之人行動線網絡進行時空平均取樣，建立熱舒適人行動線設計評估指標 $\bar{\Phi}_{\text{walkway}}$（式~\eqref{eq:walkway_exposure}），取代僅評估單一起訖點間路徑之傳統尋路演算法，使評估尺度與都市設計決策一致；其作為 NSGA-II 第三項適應度函數之工程整合列為近期後續工作。搭配 5D 違規向量（FAR、BCR、縮退、容納性、碰撞）之法規感知最佳化，本研究以已完整實作之全域 UTCI 與綠化率二目標完成真實最佳化運行，在帕累托前緣上呈現二者間的設計權衡，並誠實揭露現行目標函數在缺乏開發強度下限約束時之退化解風險，為後續研究提供具體、以實證為基礎的改進方向。
\end{enumerate}


\section{研究局限性}
\label{sec:limitations}

儘管本研究在技術實現與應用驗證上取得顯著成果，仍存在以下局限性。各項局限依「限制描述 → 原因 → 改善方向」三段式逐條說明：

\begin{itemize}
    \item \textbf{氣候情境的單一性：}模型對於冬季強烈東北季風、梅雨季連續陰天或颱風侵襲等極端氣候情境之適應能力尚未系統驗證。原因在於訓練資料以台灣夏至（6月21日）單一氣候日為基礎，此一選擇雖能鎖定全年熱壓力最嚴峻之情境以聚焦訓練訊號，卻使模型未接觸過冬季低太陽高度角、高濕陰天等大幅不同之邊界條件組合。改善方向為系統性引入春、夏、秋、冬四季代表性氣候日之 EPW 資料重新訓練或微調（見第~\ref{sec:future_work}節第 3 項）。
    \item \textbf{場地尺度局限：}模型以 $80 \times 80$ 公尺行區為模擬單元，在街廓連棟型配置下可能存在邊界效應誤差。原因在於現有異質圖架構未建立跨行區之連接邊緣，鄰近行區的熱環境交互影響（如建築群峽谷效應引發的背景風速放大）無法透過訊息傳遞機制傳入模型；此一設計取捨是為了控制單一場景之圖規模與訓練成本。改善方向為將異質圖延伸至多行區聯合建模，以行區間界面作為跨行區連接邊緣（見第~\ref{sec:future_work}節第 2 項）。
    \item \textbf{極端熱壓力等級準確率瓶頸：}UTCI $> 46°\text{C}$ 之極端熱壓力等級分類準確率相對較低，此類高熱情境恰為公共健康最需關注之情境。原因為訓練資料中此等級樣本佔比極低（僅出現於正午直射、高 SVF 節點），使模型在此區間之梯度訊號稀疏，屬於典型的類別不平衡問題而非模型結構性缺陷。改善方向為針對高溫情境進行加權採樣或額外資料擴增，提升極端等級樣本於訓練集中之代表性。
    \item \textbf{訓練資料仍以隨機生成場景為主：}第~\ref{sec:building_extraction}節建立之真實建物擷取流程，目前僅用於既有場地案例驗證，尚未回饋為大規模訓練資料集之直接來源。原因在於真實建物擷取流程於本研究階段之產出規模（單一研究場域之案例集）尚不足以支撐訓練所需之場景多樣性，且部分子區塊之樓高辨識仍需人工複核（見第~\ref{subsec:limitations}節第 5 項），限制了自動化規模化的可行性。改善方向為擴大擷取流程之涵蓋範圍並建立真實建物資料庫（見第~\ref{sec:future_work}節第 1 項），使模型面對明顯超出訓練集幾何統計特徵範圍（如超高層建築群、極端窄巷紋理）之真實場地時，預測結果目前仍應視為初步參考。
    \item \textbf{缺乏真實感測器驗證：}本研究之預測精度評估以 EnergyPlus/Ladybug 模擬輸出為基準，尚未以現地 IoT 感測器之實測 UTCI 資料進行交叉驗證。原因在於現有 IoT 感測器網路（第~\ref{subsec:iot_spatial}節）雖提供溫濕度觀測值，用於氣象邊界條件之空間降尺度校準，但並未直接量測 UTCI 所需之平均輻射溫度與風速，無法作為逐點 UTCI 之直接驗證來源。改善方向為部署具備完整四項微氣候變數量測能力之密集感測器網路，收集連續逐時實測資料以進行物理驗證（見第~\ref{sec:future_work}節第 6 項）。
    \item \textbf{現行多目標函數之退化解風險：}第~\ref{subsec:pareto_morphology}節之真實 NSGA-II 運行結果顯示，現行已實作之二目標函數（全域 UTCI 最小化、綠化率最大化）在缺乏開發強度下限約束時，會收斂至 FAR 遠低於訓練分布範圍之近乎零建築量體退化解。原因在於兩項目標之數學形式皆未要求任何建築量體存在——空地本身之預測 UTCI 通常已低，綠化率亦不以建築存在為前提，而現有 5 維違規向量僅檢核上限，未設定最低開發強度。改善方向為引入最低容積率或建築棟數下限之可行性約束，並儘速完成人行動線熱暴露目標 $f_2$ 之工程整合，以三目標之相互制衡降低單一退化方向的可能性（見第~\ref{sec:future_work}節第 7 項）。
\end{itemize}


\section{未來研究建議}
\label{sec:future_work}

基於上述研究成果與局限性，本研究提出以下未來研究方向：

\begin{enumerate}
    \item \textbf{真實建物資料庫回饋訓練集：}
    將第~\ref{sec:building_extraction}節之三服務融合擷取流程擴大應用於新竹市及周邊都市行區，建立真實建物足跡與樓高之大規模資料庫，逐步取代或補充現有隨機生成場景，使模型訓練分佈與真實都市紋理之統計特性更為一致。

    \item \textbf{多行區與街廓尺度延伸：}
    將異質圖架構延伸至多行區聯合建模，以行區間界面（道路境界線、退縮空間）作為跨行區連接邊緣，研究街廓尺度通風廊道、熱島聚集效應等大尺度都市氣候現象。此延伸需要在圖分割（Graph Partitioning）與批次訓練策略上進行適配。

    \item \textbf{季節性與多氣候場景泛化：}
    系統性地引入春、夏、秋、冬四季代表性氣候日之 EPW 訓練資料，以及不同台灣氣候分區（如花東縱谷、高屏熱帶氣候區）之場地條件，建立具備跨時間、跨地域泛化能力的「通用都市微氣候代理模型（Universal Urban Microclimate Surrogate）」。

    \item \textbf{行人流量資料整合熱舒適人行動線評估：}
    結合真實世界之行人流量資料（來自手機信令、共享單車或行人計數器），對熱舒適人行動線設計目標 $f_2$ 進行人流加權修正。當前研究以均勻權重評估動線網絡各邊，未來可融入「高人流量路段優先改善熱舒適度」的需求導向最佳化。

    \item \textbf{材質與介面精細化建模：}
    現有模型以地表溫度之一階估算（式~\eqref{eq:surface_temp}）代替精確材質建模。未來可引入不同鋪面材質（瀝青、透水磚、綠化鋪面、白色反射鋪面）之熱物性資料庫，以材質節點擴充異質圖之節點類型，進一步提升模型對都市設計精細材料決策的響應能力。

    \item \textbf{現地實測驗證與感測器網路部署：}
    在新竹市典型行區部署密集 IoT UTCI 感測器網路，收集連續夏季逐時實測資料，以此進行模型物理驗證（Physical Validation）並修正現有 EnergyPlus 模擬基準與現地微氣候之系統誤差，完善研究方法論之完整性。

    \item \textbf{最佳化引擎之開發強度下限約束與三目標整合：}
    針對第~\ref{subsec:pareto_morphology}節揭露之退化解問題，於 \texttt{07\_optimization/constraints.py} 之 5 維違規向量中新增最低容積率（FAR $\ge$ FAR$_{\min}$）或建築棟數下限之可行性約束，並將人行動線熱暴露目標 $f_2$（式~\eqref{eq:walkway_exposure}）實際整合進 \texttt{FitnessEvaluator}，完成真實三目標最佳化運行以驗證多目標相互制衡能否從根本上修正單一退化方向之風險。
\end{enumerate}

\noindent\textbf{研究展望：}

本研究所建立之物理資訊時空圖神經網路框架，結合真實都市 GIS 建物擷取流程，為未來「智慧都市設計（Smart Urban Design）」工具鏈提供了一個方法論基礎。隨著感測器網路持續密集化、計算硬體持續進化，以及建築數位孿生（Building Digital Twin）技術的成熟，可以預見在五至十年內，具備物理一致性保證、能實現跨氣候跨尺度預測且錨定於真實都市紋理的多模態都市代理模型，將成為都市設計實踐的標準工具之一。本研究提出的異質圖編碼策略、物理損失約束機制、真實建物擷取方法與代理驅動多目標優化架構，可望為此演進路徑提供具體的技術參照。


%-------------------------------------------------------------------------------
% 參考文獻
%-------------------------------------------------------------------------------

% yclee 使用 APA 參考文獻格式
\addBibToContents
\bibliographystyle{apalike}
% \bibliographystyle{apacite}
% \bibliographystyle{apacite-no-initials}
\bibliography{bib/bibliography} % (將參考文獻加到這個檔案)
% \printbibliography
%-------------------------------------------------------------------------------
% 附錄
%-------------------------------------------------------------------------------

% Start appendix
\appendix

% 附錄之目錄編號不套用「第X章」格式（附錄本身以「附錄」起首，見 thesis.cls），
% 避免目錄出現「第A章 附錄」的重複／不一致字樣
\renewcommand{\cftchappresnum}{}
\renewcommand{\cftchapaftersnum}{}

% Add appendicies to "Table of Contents"
\addAppxToContents

% 請從此開始依序擺放附錄，不須附錄可以刪除或註解掉下一行
\chapter{附錄}

本附錄承載第~\ref{chapter:methodology}章正文中僅保留結論性定義、而將完整推導與參數全表移出主文的技術細節，供有意複現或驗證本研究實作的讀者查閱。

\section{圖索引對齊機制完整計算流程}
\label{appx:index_offset_detail}

第~\ref{subsec:index_offset}節式~\eqref{eq:node_concat}--\eqref{eq:offset}所述之節點串接與索引偏移，於資料載入階段（\texttt{UTCIGraphDataset.get}）依下列順序執行：

\begin{enumerate}
    \item 建物節點特徵 $\mathbf{H}_{\text{obj}} \in \mathbb{R}^{N_{\text{obj}} \times 7}$ 依場景讀取順序取得原始索引 $0, 1, \ldots, N_{\text{obj}}-1$。
    \item 空氣節點特徵 $\mathbf{H}_{\text{air}} \in \mathbb{R}^{N_{\text{air}} \times T \times 9}$ 之 KNN 毗鄰邊（\texttt{contiguity}）先以 \texttt{scipy.spatial.cKDTree} 在空氣節點局部座標系（$0, \ldots, N_{\text{air}}-1$）中求解 $k=8$ 近鄰，此時邊索引尚未偏移。
    \item 串接張量 $\mathbf{H} = [\mathbf{H}_{\text{obj}}; \mathbf{H}_{\text{air}}]$ 完成後，空氣節點於全域索引空間中的實際位置為原索引 $+N_{\text{obj}}$，因此所有 \texttt{contiguity} 邊之 \texttt{edge\_index} 兩端點皆須加上 $N_{\text{obj}}$ 偏移。
    \item \texttt{semantic} 邊（建物 $\leftrightarrow$ 建物）之索引本即位於 $0, \ldots, N_{\text{obj}}-1$ 範圍內，無需偏移。
    \item 動態邊集合 \texttt{dynamic\_edges[t]}（\texttt{shadow}、\texttt{veg\_et}、\texttt{convective}）之來源端或目標端凡屬空氣節點者，於各時步重新計算時同步套用相同偏移規則。
\end{enumerate}

此偏移規則若未正確套用，將導致模型在前向傳播時把空氣節點的特徵誤植於建物節點索引位置，是實作階段最常見但難以察覺的錯誤來源之一，故於此完整記錄以利複現。


\section{RGCN 各關係矩陣維度設定}
\label{appx:rgcn_detail}

第~\ref{subsec:rgcn}節式~\eqref{eq:rgcn}中，$|\mathcal{R}|=5$ 種關係各自維護獨立投影矩陣 $\mathbf{W}_r \in \mathbb{R}^{256 \times 256}$（隱藏維度 $d=256$，即物件與空氣編碼器 128 維輸出經 RGCNBlock 內部升維後之工作維度），加上自迴路矩陣 $\mathbf{W}_{\text{self}} \in \mathbb{R}^{256\times256}$，單層 RGCNBlock 之可訓練參數量為：

\begin{equation}
    |\theta_{\text{RGCNBlock}}| = (|\mathcal{R}| + 1) \times 256^2 = 6 \times 65{,}536 = 393{,}216
    \label{eq:rgcn_params_appx}
\end{equation}

三層堆疊（$L=3$）之 RGCN 編碼器總參數量約為 $1.18 \times 10^6$（不含 LayerNorm 與 Dropout 之可訓練縮放參數）。各層之間以殘差連接直接相加，未使用額外投影對齊維度，故要求所有層之隱藏維度一致固定為 256。


\section{模型完整超參數表}
\label{appx:model_hparams}

第~\ref{subsec:lstm}節所述模型之完整超參數設定，詳見表~\ref{tab:model_params}；訓練最佳化器為 Adam，初始學習率 $1\times10^{-3}$，搭配 \texttt{ReduceLROnPlateau} 排程（驗證集 loss 連續 10 個 epoch 未改善時學習率乘以 0.5），Early Stopping 之耐心值（patience）設為 20 個 epoch。

\begin{table}[htbp]
    \centering
    \caption{PI-ST-GNN 模型超參數設定}
    \label{tab:model_params}
    \small
    \begin{tabular}{llcc}
        \toprule
        \textbf{子模組} & \textbf{超參數} & \textbf{設定值} & \textbf{備註} \\
        \midrule
        物件節點編碼器     & 輸入維度         & 7          & 靜態幾何特徵 \\
                          & 隱層→輸出維度   & 64 → 128   & 兩層 MLP \\
        空氣節點編碼器     & 輸入維度         & 9          & 每時步動態特徵 \\
                          & 隱層→輸出維度   & 64 → 128   & 兩層 MLP \\
        全域情境編碼器     & 輸入維度         & 8          & 全場域統計+差分 \\
                          & 輸出維度         & 192        & 三層 MLP \\
        RGCN 塊            & 輸入→輸出維度   & 256 → 256  & 含殘差+LN \\
                          & 堆疊層數         & 3          & --- \\
                          & 關係數           & 5          & 5 種邊類型 \\
                          & Dropout          & 0.1        & 訓練期間 \\
        Fusion MLP         & 輸入維度         & 256 + 192  & RGCN + 情境 \\
                          & 輸出維度         & 256        & 兩層 MLP \\
        LSTM               & 隱藏維度         & 256        & 單層 \\
                          & 初始隱藏狀態     & RGCN 投影  & 非零初始 \\
        輸出 MLP           & 輸入維度         & 256        & --- \\
                          & 輸出維度         & 1（逐時步）& 純量 UTCI \\
        \midrule
        \textbf{學習率}   & 初始值           & $1 \times 10^{-3}$ & ReduceLROnPlateau \\
        \textbf{批次大小} & ---              & 1（逐場景） & 異質圖結構差異 \\
        \textbf{最大 Epoch} & ---            & 250         & 含 Early Stopping \\
        \bottomrule
    \end{tabular}
\end{table}


\section{物理引導損失函數完整推導}
\label{appx:physics_loss_derivation}

三項物理引導損失均遵循同一套推導邏輯：先將物理原則表述為一條不等式約束，再以\textbf{外點懲罰法}（Exterior Penalty Method）將該不等式約束轉化為可微分的軟性懲罰項，使其能以梯度下降與資料損失聯合最佳化，而不需要對神經網路輸出施加硬性投影。

\subsection{外點懲罰法的一般形式}

對任一應恆成立之不等式約束 $g(\hat{u}) \le 0$，外點懲罰法將其轉化為懲罰項：

\begin{equation}
    \mathcal{L}_{\text{penalty}} = \left[\text{ReLU}\big(g(\hat{u})\big)\right]^p,\qquad p \in \{1, 2\}
    \label{eq:exterior_penalty_general}
\end{equation}

當 $g(\hat{u}) \le 0$（約束滿足）時 $\text{ReLU}(g)=0$，懲罰為零、不產生梯度，訓練不受干擾；當 $g(\hat{u}) > 0$（約束違反）時懲罰隨違反量增加。指數 $p=2$（平方鉸鏈）使懲罰在違反邊界處梯度連續且趨於零，讓最佳化在約束邊界附近平滑；$p=1$（線性鉸鏈）則在違反邊界處梯度不連續但恆定，使懲罰強度不隨違反程度二次放大。本研究三項約束依各自對訓練穩定性之要求，分別採用兩種形式，詳見以下三小節。

\subsection{輻射空間一致性約束（$\mathcal{L}_{\text{rad}}$）完整推導}

\noindent\textbf{物理原則}：於太陽高度角 $\theta_{\text{sol}} > 10^\circ$ 之有效日照時段，完全日照節點之熱壓力應顯著高於完全遮蔭節點，兩者間應存在至少 $m_{\text{rad}}$ 之最小間隔，而非僅要求「不反轉」。

\noindent\textbf{約束表述}：令 $\bar{u}_{\text{sun}}$、$\bar{u}_{\text{shade}}$ 為當前時步日照節點集合 $\mathcal{S}_{\text{sun}} = \{i : \text{SVF}_i > 0.5,\ \text{shadow}_i = 0\}$ 與陰影節點集合 $\mathcal{S}_{\text{shade}} = \{i : \text{SVF}_i < 0.5,\ \text{shadow}_i = 1\}$ 之預測均值，則所需約束為：

\begin{equation}
    g_{\text{rad}} = \bar{u}_{\text{shade}} - \bar{u}_{\text{sun}} + m_{\text{rad}} \le 0
    \quad\Longleftrightarrow\quad
    \bar{u}_{\text{sun}} - \bar{u}_{\text{shade}} \ge m_{\text{rad}}
    \label{eq:rad_constraint}
\end{equation}

\noindent\textbf{套用外點懲罰法}（$p=2$，平方鉸鏈，使邊界處梯度平滑）：

\begin{equation}
    \mathcal{L}_{\text{rad}} = \left[\text{ReLU}\big(g_{\text{rad}}\big)\right]^2
    = \left[\text{ReLU}\big(\bar{u}_{\text{shade}} - \bar{u}_{\text{sun}} + m_{\text{rad}}\big)\right]^2
\end{equation}

此即正文式~\eqref{eq:lrad}。$m_{\text{rad}} = 0.5$ 為正規化 UTCI 空間中之設定值，對應約半個標準差之最小溫差要求。當太陽高度角 $\theta_{\text{sol}} \le 10^\circ$（清晨、黃昏或夜間）時，直射輻射貢獻可忽略，此時 $\mathcal{S}_{\text{shade}}$ 與 $\mathcal{S}_{\text{sun}}$ 均設為空集合，該時步不產生懲罰梯度，避免在低輻射時段對模型施加不具物理意義的約束；訓練時，$\mathcal{L}_{\text{rad}}$ 對所有滿足 $\theta_{\text{sol}}>10^\circ$ 之有效時步取平均。

\subsection{都市熱慣性時序平滑約束（$\mathcal{L}_{\text{temp}}$）完整推導}

\noindent\textbf{物理原則}：都市鋪面之熱儲存效應限制 UTCI 逐時變化率不得超過 $\dot{u}_{\text{max}} = 5\,^\circ\text{C/hr}$。

\noindent\textbf{約束表述}：模擬以整點為間隔（$\Delta t = 1$ hr），約束為 $|u_{i,t+1} - u_{i,t}| \le \dot{u}_{\text{max}} \cdot \Delta t = 5\,^\circ\text{C}$（原始溫度單位）。

\noindent\textbf{單位換算至正規化空間}：模型於 z-score 正規化空間中預測，換算需除以 UTCI 之標準差 $\sigma_{\text{UTCI}}$：

\begin{equation}
    \delta_{\text{temp}} = \frac{\dot{u}_{\text{max}} \cdot \Delta t}{\sigma_{\text{UTCI}}}
    \label{eq:delta_temp_derivation}
\end{equation}

實作中 $\delta_{\text{temp}}$ 取值 $0.625 = 5/8$，對應假設 $\sigma_{\text{UTCI}} \approx 8\,^\circ\text{C}$。然而，本研究訓練資料集（\texttt{ground\_truth\_v2.h5}）實測之 UTCI 標準化統計量為 $\sigma_{\text{UTCI}} = 2.9841\,^\circ\text{C}$（均值 $\mu_{\text{UTCI}} = 35.6146\,^\circ\text{C}$），與 $\delta_{\text{temp}}$ 換算所隱含之 $8\,^\circ\text{C}$ 假設值有明顯落差。若代入實測標準差，依式~\eqref{eq:delta_temp_derivation}反推，$\delta_{\text{temp}}$ 之精確值應為：

\begin{equation}
    \delta_{\text{temp}}^{\text{(實測校正)}} = \frac{5}{2.9841} \approx 1.676
\end{equation}

亦即，若以實際訓練所用之 $\delta_{\text{temp}} = 0.625$ 反推其對應之\textbf{實際容許原始變化率}，可得：

\begin{equation}
    \dot{u}_{\text{max}}^{\text{(實際生效)}} = \delta_{\text{temp}} \times \sigma_{\text{UTCI}} = 0.625 \times 2.9841 \approx 1.865\,^\circ\text{C/hr}
\end{equation}

換言之，本研究實際訓練並報告之模型，其時序平滑約束在原始溫度尺度上生效的容許變化率約為 $1.87\,^\circ\text{C/hr}$，較正文所述之設計目標 $5\,^\circ\text{C/hr}$ 嚴格約 2.7 倍。此一落差源於 $\delta_{\text{temp}}$ 設定時所採用之標準差假設值與訓練資料實測值不一致，特此誠實揭露；由於本研究之全部量化結果（第~\ref{chapter:experiments}章）均基於 $\delta_{\text{temp}}=0.625$ 之模型訓練與評估，此處不回溯更動已訓練模型之設定，而將此列為後續研究應修正之已知實作落差（見第~\ref{sec:limitations}節）。

\noindent\textbf{套用外點懲罰法}（$p=1$，線性鉸鏈，避免訓練初期時序尚不穩定時梯度二次放大而主導總損失）：

\begin{equation}
    \mathcal{L}_{\text{temp}} = \text{mean}_{i,t}\!\left[\text{ReLU}\big(|u_{i,t+1} - u_{i,t}| - \delta_{\text{temp}}\big)\right]
\end{equation}

此即正文式~\eqref{eq:ltemp}，求和範圍 $t=0,\ldots,T-2$（即相鄰時步差分），對應 $T=11$ 個模擬時步共可產生 10 組差分；首個時步（$t=0$，對應 08:00）因無前一時步可比較，不計入此項損失。

\subsection{風場阻滯尾流約束（$\mathcal{L}_{\text{wind}}$）完整推導}

\noindent\textbf{物理原則}：高層建築背風側易形成熱累積之「熱死角」，此類節點之 UTCI 不應顯著偏離場域統計中心。

\noindent\textbf{統計離群定義}：在常態分布假設下，$P(Z > 1.5) \approx 6.7\%$（$Z$ 為標準常態分布），故 $\mu_{\hat{u}} + 1.5\sigma_{\hat{u}}$ 界定了當前時步全場域預測值中最熱的約 6.7\% 尾端，作為「顯著偏離」之操作型定義：

\begin{equation}
    g_{\text{wind}} = \hat{u}_i - \mu_{\hat{u}} - 1.5\,\sigma_{\hat{u}} \le 0,
    \qquad i \in \mathcal{B}_{\text{tall}} = \{i : B_h/50 > 0.4\}
\end{equation}

\noindent\textbf{套用外點懲罰法}（$p=1$，線性鉸鏈，理由同 $\mathcal{L}_{\text{temp}}$）：

\begin{equation}
    \mathcal{L}_{\text{wind}} = \text{mean}_{i \in \mathcal{B}_{\text{tall}}}\!\left[\text{ReLU}\big(\hat{u}_i - \mu_{\hat{u}} - 1.5\,\sigma_{\hat{u}}\big)\right]
\end{equation}

此即正文式~\eqref{eq:lwind}。其中 $\mu_{\hat{u}}$、$\sigma_{\hat{u}}$ 取當前時步\textbf{全場域}（所有節點，非僅 $\mathcal{B}_{\text{tall}}$）預測值之均值與標準差，逐時步重新計算，$\mathcal{B}_{\text{tall}}$ 亦逐時步依當時風向與建築幾何關係重新判定背風側範圍，而非訓練前一次性標記。


\section{熱舒適人行動線熱暴露指標完整推導}
\label{appx:walkway_derivation}

\subsection{無量綱熱壓力懲罰函數 $\Phi(\cdot)$ 之推導}

第~\ref{subsec:from_astar_to_walkway}節式~\eqref{eq:walkway_exposure}所用之熱壓力懲罰函數，其設計目標為：低於舒適閾值時不產生懲罰，高於閾值時懲罰隨超出量遞增，且懲罰量正規化至 $[0,1]$ 便於跨場景比較。此三項需求對應下列分段函數：

\begin{equation}
    \Phi(\text{UTCI}_j) =
    \begin{cases}
        0, & \text{UTCI}_j \le \text{UTCI}_{\text{comfort}} \\[4pt]
        \dfrac{\text{UTCI}_j - \text{UTCI}_{\text{comfort}}}{\text{UTCI}_{\text{max}} - \text{UTCI}_{\text{comfort}}}, & \text{UTCI}_j > \text{UTCI}_{\text{comfort}}
    \end{cases}
    \label{eq:penalty_appx}
\end{equation}

以 ReLU 合併兩種情形，可寫成正文所引用之單行形式：

\begin{equation}
    \Phi(\text{UTCI}_j) = \frac{\text{ReLU}\big(\text{UTCI}_j - \text{UTCI}_{\text{comfort}}\big)}{\text{UTCI}_{\text{max}} - \text{UTCI}_{\text{comfort}}}
\end{equation}

其中 $\text{UTCI}_{\text{comfort}} = 26\,^\circ\text{C}$ 為熱舒適閾值上限；$\text{UTCI}_{\text{max}}$ 為研究時段內場域可能出現之最大 UTCI 值，用以將懲罰量正規化至 $[0,1]$ 區間。以本研究訓練資料集（20 個場景抽樣）實測為例，正午高輻射時段之場域最大 UTCI 可達約 $55$--$56\,^\circ\text{C}$，故 $\text{UTCI}_{\text{max}}$ 建議依實際部署場域之逐時模擬結果動態設定，而非採用固定常數，以避免不同氣候條件或場地尺度下正規化基準失真。

\subsection{時空平均人行動線熱暴露指標之推導}

單一時步、單一路段之熱暴露定義為 $D_{ij}\cdot\Phi(\text{UTCI}_{j,t})$，即以邊之實際步行距離加權之熱壓力懲罰，確保較長路段之熱暴露對總指標的貢獻與其實際行走時間成比例。對整個人行動線網絡 $E_{\text{walk}}$ 取空間平均，再對 $T$ 個模擬時步取時間平均，即得正文式~\eqref{eq:walkway_exposure}：

\begin{equation}
    \bar{\Phi}_{\text{walkway}}
    = \underbrace{\frac{1}{T}\sum_{t=1}^{T}}_{\text{時間平均}}\ \
      \underbrace{\frac{1}{|E_{\text{walk}}|}\sum_{(i,j)\in E_{\text{walk}}} D_{ij}\cdot\Phi\!\left(\text{UTCI}_{j,t}\right)}_{\text{空間平均（距離加權）}}
\end{equation}

此一時空平均指標之量綱與 $f_1$（全域時空平均 UTCI）一致，使 NSGA-II 之 $f_1$、$f_2$ 兩項熱壓力相關目標可在同一最佳化框架下直接比較收斂行為，而不需額外的尺度校正。


\section{NSGA-II 適應度求值器輸入維度修正紀錄}
\label{appx:fitness_dim_fix}

第~\ref{subsec:pareto_morphology}節所呈現之真實 NSGA-II 收斂結果，係於本論文撰寫過程中修正一項適應度求值器之輸入維度缺陷後方能取得。原始 \texttt{07\_optimization/fitness.py} 之 \texttt{FitnessEvaluator} 建構子在初始化 \texttt{GNNInputBuilder}（\texttt{06\_deployment/geometry\_converter.py}）時未傳入 \texttt{dim\_air} 參數，導致其採用套件預設值 \texttt{dim\_air=8}；然而第~\ref{sec:gnn_model}章訓練完成之 V2 checkpoint 之空氣節點編碼器輸入維度為 9（含地表溫度 $T_s$ 特徵，見第~\ref{subsec:nodes}節），兩者不一致造成矩陣乘法維度錯誤（\texttt{RuntimeError: mat1 and mat2 shapes cannot be multiplied}），使適應度求值器對任何染色體個體皆拋出例外，落入例外處理分支回傳懲罰值（UTCI $=60^\circ\text{C}$、綠化率 $=0$），造成整個族群在演化過程中呈現虛假的「零多樣性」收斂假象。本研究於 \texttt{FitnessEvaluator} 建構子與 \texttt{06\_deployment/app.py} 之對應呼叫處新增顯式 \texttt{dim\_air} 參數傳遞（對齊自 checkpoint 權重形狀自動偵測之維度），修正後之適應度求值器方能對每個染色體個體正確執行 GNN 推理，取得第~\ref{subsec:pareto_morphology}節所呈現之真實、非退化之收斂軌跡與帕累托前緣。此一修正記錄於此，以說明圖~\ref{fig:nsga2_real}之收斂數據並非模擬或示意性繪製，而是修正後管線之直接輸出。


\section{Grasshopper 元件使用者操作手冊}
\label{appx:user_manual}

本附錄以第~\ref{sec:nsga2}節所定義之系統為對象，提供一份面向參數化設計／數位建築程式開發者的操作手冊，說明從「周遭環境模型建好之後」到「拿到熱舒適評估結果」之間，使用者實際需要做的每一步。內容對應本研究實際完成之兩個 Grasshopper 元件——\texttt{UTCIPredictor}（單次即時預測）與 \texttt{UTCIOptimizer}（多目標最佳化搜尋）。

\subsection{使用前的準備工作}
\label{appx:user_manual_prep}

在 Grasshopper 中開始操作之前，有兩件事需要先完成一次，之後就不用重複：

\begin{enumerate}
    \item \textbf{啟動後端運算服務}：本系統的 GNN 模型並不是跑在 Grasshopper 元件裡面，而是跑在一支獨立啟動的背景程式（伺服器）上，Grasshopper 元件只是透過網路跟它「問問題、拿答案」。因此每次要使用前，須先在終端機執行一次啟動指令（等同於「打開這套運算服務的電源」），啟動後即可保持開啟，供多次設計迭代重複使用，不需每次重開。
    \item \textbf{在 Rhino 的 Script Editor 環境安裝一個通訊套件}：僅需執行一次 \texttt{pip install websocket-client}，讓 Grasshopper 元件具備與後端服務對話的能力。這與一般安裝 Grasshopper 外掛的方式不同，是在 Rhino 8 內建的 Python 環境中安裝，而非系統上其他的 Python。
\end{enumerate}

完成以上兩步後，Grasshopper 元件即可持續使用，除非重開機或伺服器程式被關閉，否則不需重複準備。

\subsection{準備場地幾何——這是使用者最主要的工作}
\label{appx:user_manual_geometry}

周遭環境模型（基地紋理、既有建物、喬木）建好之後，使用者需要把這些幾何物件，接到元件對應的輸入端口上。可以把每個輸入端口想成一個「貼標籤」的插槽，只要把正確類型的 Rhino 幾何拖進去即可，如表~\ref{tab:gh_inputs} 所整理：

\begin{table}[htbp]
    \centering
    \caption{Grasshopper 元件輸入端口與使用者準備事項對照}
    \label{tab:gh_inputs}
    \small
    \begin{tabular}{lll}
        \toprule
        \textbf{輸入端口} & \textbf{幾何類型} & \textbf{使用者要準備什麼} \\
        \midrule
        site\_boundary     & 封閉曲線     & 一條圈出基地範圍的封閉多邊形 \\
        context\_buildings & 建物量體清單 & 周遭「固定不動」的既有建物 \\
        design\_buildings  & 建物量體清單 & 設計中、要被評估的新建物（僅預測元件使用） \\
        tree\_points       & 點清單       & 喬木種植位置 \\
        tree\_radii        & 數值清單     & 對應每棵樹的樹冠半徑（公尺） \\
        tree\_heights      & 數值清單     & 對應每棵樹的樹高（公尺） \\
        server\_url        & 文字         & 通常不需更動，預設已指向本機服務 \\
        run                & 開關（True/False）& 切成 True 才會真正送出計算請求 \\
        \bottomrule
    \end{tabular}
\end{table}

一個實務上的小提醒：喬木半徑／高度清單的順序必須與喬木位置清單一一對應（第 1 個半徑對應第 1 個位置，以此類推），若順序對不上，系統仍會運算，但每棵樹的遮蔭效果會被套用到錯誤的位置上。

\subsection{操作流程 A：單次即時評估（UTCIPredictor）}
\label{appx:user_manual_predict}

這個元件適合「我現在想看這個設計方案的熱舒適表現如何」的情境，操作步驟如下：

\begin{enumerate}
    \item 依第~\ref{appx:user_manual_geometry}節，把場地邊界、既有建物、設計建物、喬木資訊接上對應輸入端口。
    \item 把 \texttt{run} 接上一個開關元件（Boolean Toggle），切成 True。
    \item 系統會在約數十毫秒內回傳結果——這個速度快到可以邊拖動滑桿調整建築高度、邊即時看到熱舒適熱力圖跟著變化，不需要每次都手動重新觸發。
    \item 若之後修改了建物或喬木的幾何，只要 \texttt{run} 仍保持 True，結果會自動更新；不需要的時候把 \texttt{run} 切回 False 即可暫停運算。
\end{enumerate}

\subsection{操作流程 B：多目標最佳化搜尋（UTCIOptimizer）}
\label{appx:user_manual_optimize}

這個元件適合「我不確定哪種建築配置最好，想讓系統自動幫我搜尋一批候選方案」的情境，除了場地邊界與既有建物外，還需要多準備幾項設定：

\begin{itemize}
    \item \textbf{法規限制}（\texttt{constraints}）：例如容積率上限、建蔽率上限、退縮距離，系統會確保搜尋出來的方案都符合這些限制。
    \item \textbf{搜尋規模}：要優化幾棟建築（\texttt{n\_buildings}）、幾棵樹（\texttt{n\_trees}），以及要嘗試多少組候選方案（\texttt{pop\_size}）、反覆演化幾輪（\texttt{n\_gen}）。這兩個數字越大，找到的方案通常越多樣，但等待時間也越長。
\end{itemize}

操作步驟：

\begin{enumerate}
    \item 設定好上述參數後，把 \texttt{run} 切成 True 啟動搜尋。搜尋過程在背景執行緒中進行，\textbf{不會卡住 Rhino 畫面}，使用者可以在等待的同時繼續操作模型其他部分。
    \item 因為整個搜尋通常需要數十秒到數分鐘才能完成，且過程中會持續產生新一代的候選方案，建議另外接一個 Grasshopper 內建的計時器元件（Timer，觸發間隔建議設為 1 秒），讓這個最佳化元件被重複觸發、持續讀取最新進度——如果沒有接計時器，畫面就會停留在使用者上一次手動觸發當下的結果，不會自動更新。
    \item 如果想中途停止（例如發現參數設錯了），把 \texttt{cancel} 輸入端口切成 True 即可安全中斷，已經算出的候選方案不會遺失。
\end{enumerate}

\subsection{看得懂拿到的結果——輸出如何解讀}
\label{appx:user_manual_output}

本系統沒有另外產生一份 PDF 或文件報告，\textbf{評估結果就是元件輸出端口本身}，即時顯示在 Rhino 場景裡，讓使用者可以直接在自己的模型上比對，而不需切換到其他軟體查看。

\textbf{UTCIPredictor（單次評估）拿到的東西：}

\begin{itemize}
    \item \textbf{熱力圖網格（utci\_mesh）}：一片鋪在場地上、依顏色深淺呈現熱壓力高低的網格，藍色代表涼爽、經黃橘色至深紅色代表越來越炎熱，設計者可以直接看出陰影覆蓋不足或悶熱的角落在哪裡。
    \item \textbf{數值指標（mean\_utci／max\_utci／min\_utci）}：全場地的平均、最高、最低熱壓力數值，方便用具體數字比較不同設計版本的優劣。
    \item \textbf{狀態文字（status）}：一行簡短摘要，例如「OK | Mean: 30.8°C  Max: 42.1°C」，方便快速確認這次運算是否成功、結果大致落在什麼範圍。
\end{itemize}

\textbf{UTCIOptimizer（最佳化搜尋）拿到的東西：}

\begin{itemize}
    \item \textbf{一組候選方案（pareto\_designs／pareto\_geometries）}：系統不會只給「一個最好的答案」，而是給一整批「各有取捨、但都值得參考」的候選方案——例如某方案熱舒適度最好但綠化較少、另一方案綠化最多但熱舒適度普通——連同對應的 3D 建物量體，讓設計者依自己的判斷選擇，而非被單一答案綁住。
    \item \textbf{兩個代表方案（best\_utci\_design／best\_green\_design）}：分別是「全場最涼」與「綠化率最高」的極端方案量體，可作為理解候選方案分布範圍的快速參考點。
    \item \textbf{搜尋過程紀錄（convergence\_data）}：記錄每一輪搜尋的最佳表現，讓使用者可以自行畫出一條「越搜尋、方案越好」的進步曲線，確認搜尋確實有在收斂，而不是原地打轉。
    \item \textbf{即時進度文字（progress\_text／status）}：搜尋進行中會持續回報目前跑到第幾輪、目前最好的方案表現如何。
\end{itemize}

\subsection{常見狀況排除}
\label{appx:user_manual_troubleshoot}

\begin{itemize}
    \item \textbf{狀態顯示「無法連線」}：通常代表附錄~\ref{appx:user_manual_prep}第 1 步的後端服務尚未啟動，或已經被關閉，需要重新啟動一次。
    \item \textbf{狀態顯示「websocket-client 未安裝」}：代表附錄~\ref{appx:user_manual_prep}第 2 步的通訊套件尚未安裝，依指示安裝後重新開啟 Rhino 檔案即可。
    \item \textbf{場地邊界輸入被判定錯誤}：通常是因為封閉曲線的節點數過少（少於 3 個點無法構成一個範圍），確認邊界曲線確實封閉、且形狀合理即可。
    \item \textbf{最佳化畫面卡住不動}：最常見原因是忘記接計時器元件（見附錄~\ref{appx:user_manual_optimize}第 2 步），導致元件沒有被重複觸發去讀取最新進度，而非搜尋真的卡住。
\end{itemize}



%-------------------------------------------------------------------------------
\end{document}